% concatenated from parabolicgap.tex
\documentclass[12pt]{amsart}
\usepackage{amssymb}
\usepackage{amsmath,amscd}
\usepackage{amsfonts,latexsym,verbatim,amscd,mathrsfs,color,array}
\usepackage[all,cmtip]{xy}
\usepackage{mathrsfs}
\usepackage{multirow}
\usepackage{graphicx}
\usepackage{amsfonts, amsthm}
\usepackage{bbm}
\usepackage[hmargin=1in,vmargin=1in]{geometry}
\usepackage{mathdots}
\usepackage{stmaryrd}
\usepackage{MnSymbol}
\usepackage{bbm}
\usepackage[hidelinks]{hyperref}
\usepackage{enumitem}
\usepackage[all]{xy}
\usepackage{tikz-cd}
\usepackage{empheq}

\allowdisplaybreaks

\def\Xint#1{\mathchoice
	{\XXint\displaystyle\textstyle{#1}}%
	{\XXint\textstyle\scriptstyle{#1}}%
	{\XXint\scriptstyle\scriptscriptstyle{#1}}%
	{\XXint\scriptscriptstyle\scriptscriptstyle{#1}}%
	\!\int}
\def\XXint#1#2#3{{\setbox0=\hbox{$#1{#2#3}{\int}$ }
		\vcenter{\hbox{$#2#3$ }}\kern-.6\wd0}}
\def\ddashint{\Xint=}
\def\dashint{\Xint-}

%\setlength{\oddsidemargin}{0.125in}
%\setlength{\evensidemargin}{0.125in}
%\setlength{\textwidth}{6.375in}
%\setlength{\textheight}{8.5in}
%\setlength{\extrarowheight}{0.2cm}
%\topskip 0in
%\topmargin 0.375in
%\footskip 0.25in
\renewcommand{\baselinestretch}{1.08}
%\numberwithin{equation}{section}
\renewcommand{\Re}{\operatorname{Re}}
\renewcommand{\Im}{\operatorname{Im}}
\newcommand{\Sym}{\operatorname{Sym}}
\renewcommand{\O}{\mathrm{O}}
\newcommand{\rank}{\operatorname{rank}}

\newcommand{\na}{\nabla}
\newcommand{\tensor}{\otimes}
\newcommand{\Tr}{\operatorname{Tr}}
\newcommand{\transint}{\cap\kern-0.63em|\kern0.7em}

%\newcommand{\intprod}{\mathbin{\raisebox{\depth}{\scalebox{1}[-1]{$\lnot$}}}}
\DeclareMathSymbol{\intprod}{\mathbin}{MnSyC}{'270}
\DeclareMathOperator{\Pic}{Pic}
\newcommand{\calE}{\mathcal{E}}
\newcommand{\calF}{\mathcal{F}}
\newcommand{\YM}{\mathcal{YM}}
\newcommand{\vol}{\mathrm{Vol}}
\newcommand{\LB}{\left[}
\newcommand{\RB}{\right]}
\newcommand{\gen}[1]{\LA #1 \RA}
\newcommand{\norm}[1]{\left|\!\!\>\left| #1 \right|\!\!\>\right|}
\newcommand{\LA}{\left\langle}
\newcommand{\RA}{\right\rangle}
\newcommand{\m}{{\rm modulo}\:}
\newcommand{\im}{{\rm im}\:}
\newcommand{\p}{{ \partial}}
\newcommand{\GI}{{\mathbb Z}[i]}
\newcommand{\mo}{{\rm mod}\:}
\newcommand{\F}{{\mathbb F}}
\newcommand{\Z}{{\mathbb Z}}
\newcommand{\N}{{\mathbb N}}
\newcommand{\C}{{\mathbb C}}
\newcommand{\CP}{{\mathbb{CP}}}
\newcommand{\RP}{{\mathbb{RP}}}
\renewcommand{\P}{{\mathbb P}}
\renewcommand{\1}{{\mathbf 1}}
\newcommand{\Q}{{\mathbb Q}}
\newcommand{\R}{{\mathbb R}}
\newcommand{\calB}{\mathcal{B}}
%\newcommand{\G}{{\mathbb G}}
\newcommand{\rH}{{\mathrm H}}
\newcommand{\A}{{\mathbb A}}
\renewcommand{\H}{{\mathbb H}}
\newcommand{\HP}{\mathbb{HP}}
\newcommand{\bi}{{\mathbf i}}
\newcommand{\bj}{{\mathbf j}}
\newcommand{\bk}{{\mathbf k}}
\newcommand{\rB}{{\mathrm{B} }}
\newcommand{\Hom}{{\operatorname{Hom} }}
\newcommand{\Fr}{{\mathrm{Fr} }}
\newcommand{\Lie}{{\operatorname{Lie} }}
\newcommand{\Aut}{{\operatorname{Aut} }}
\newcommand{\Vect}{\operatorname{Vect}}
\newcommand{\SU}{{\mathrm{SU} }}
\newcommand{\su}{\mathfrak{su}}
\newcommand{\Sp}{{\mathrm{Sp} }}
\newcommand{\fraksp}{{\mathfrak{sp} }}
\DeclareMathOperator{\Ad}{Ad}
\newcommand{\U}{{\mathrm U}}
\newcommand{\gothu}{\mathfrak{u}}
\newcommand{\SO}{{\mathrm{SO} }}
\newcommand{\so}{\mathfrak{so}}
\renewcommand{\rB}{{\mathrm B}}
\newcommand{\G}{{\mathrm G}}
\newcommand{\rG}{{\mathrm G}}
\newcommand{\GL}{{\mathrm{GL} }}
\newcommand{\gl}{\mathfrak{gl}}
\newcommand{\SL}{{\mathrm{SL} }}
\newcommand{\rK}{{\mathrm K}}
\newcommand{\rHol}{{\mathrm Hol }}
\newcommand{\Spin}{{\mathrm{Spin} }}
\newcommand{\supp}{{\mathrm{supp} \, }}
\newcommand{\rk}{{\mathrm{rk}}}
\newcommand{\Ker}{{\operatorname{Ker}}}
%\newcommand{\cO}{\mathcal{O}}
\renewcommand{\div}{{\operatorname{div} }}
\renewcommand{\o}{{\mathfrak{o} }}
\newcommand{\bbP}{{\mathbb P}}
\newcommand{\del}{{\partial}}
\newcommand{\delbar}{{\overline{\partial}}}
\newcommand{\gm}{{\mathfrak m}}
\newcommand{\gn}{{\mathfrak n}}
\newcommand{\gothg}{{\mathfrak g}}
\newcommand{\gothso}{\mathfrak {so }}
\newcommand{\gothk}{{\mathfrak k}}
\newcommand{\ph}{\varphi}
\newcommand{\gothh}{{\mathfrak h}}
\newcommand{\gothb}{{\mathfrak b}}
\newcommand{\gq}{{\mathfrak q}}
\newcommand{\gothC}{{\mathfrak C}}
\newcommand{\gothG}{{\mathfrak G}}
\newcommand{\gothB}{{\mathfrak B}}
\newcommand{\scrA}{{\mathscr A}}
\newcommand{\calA}{\mathcal{A}}
\newcommand{\scrB}{{\mathscr B}}
\newcommand{\scrD}{{\mathscr D}}
\newcommand{\scrF}{{\mathscr F}}
\newcommand{\scrG}{{\mathscr G}}
\newcommand{\scrI}{{\mathscr I}}
\newcommand{\scrJ}{{\mathscr J}}
\newcommand{\scrL}{{\mathscr L}}
\newcommand{\scrO}{{\mathscr O}}
\newcommand{\actson}{\curvearrowright}
\newcommand{\scrC}{{\mathscr C}}
\newcommand{\scrM}{{\mathscr M}}
\newcommand{\calM}{\mathcal{M}}
\newcommand{\scrN}{{\mathscr N}}
\newcommand{\calN}{\mathcal{N}}
\newcommand{\T}{{\mathcal T}}
\newcommand{\Grass}{{\mbox{Grass}}}
\newcommand{\Spec}{{\mbox{Spec}}}
\newcommand{\Frob}{{\mbox{Frob}}}
\newcommand{\dist}{{\operatorname{dist}}}
\newcommand{\Div}{{\operatorname{Div}}}
\newcommand{\End}{{\operatorname{End} \,}}
\newcommand{\dbar}{{\bar{\partial}}}
\renewcommand{\p}{{\partial}}
\newcommand{\eps}{{\varepsilon}}
\newcommand{\vareps}{{\varepsilon}}
\newcommand{\doublesup}[1][0pt]{%
	\mathrel{\raisebox{#1}{$\sup$}}%
}
%\newcommand{\inj}{{\operatorname{inj} }}
\renewcommand{\d}[1]{\frac{\p}{\p #1}}
\newcommand{\codim}{\operatorname{codim} }
\newtheorem{thm}{Theorem}[section]
\newtheorem{lemma}[thm]{Lemma}
\newtheorem*{lemma*}{Lemma}
\newtheorem{prop}[thm]{Proposition}
\newtheorem{ass}[thm]{Assumption}
\newtheorem{cor}[thm]{Corollary}
\newtheorem{conj}[thm]{Conjecture}
\newtheorem*{conj*}{Conjecture}
\newtheorem{defnlemma}[thm]{Definition/Lemma}
%\newtheorem{defn}{Definition}
%\newtheorem{example}{Example}
%\newtheorem{rmk}{Remark}
%\newenvironment{defn}[1][Definition.]{\begin{trivlist}
%\item[\hskip \labelsep {\bfseries #1}]}{\end{trivlist}}
%\newenvironment{example}[1][Example.]{\begin{trivlist}
%\item[\hskip \labelsep {\bfseries #1}]}{\end{trivlist}}
%\newenvironment{rmk}[1][Remark.]{\begin{trivlist}
%\item[\hskip \labelsep {\bfseries #1}]}{\end{trivlist}}
\newenvironment{claim}{\par\medskip\noindent\textit{Claim.}\space}{\par\medskip}
\newenvironment{claimproof}{\par\noindent\textit{Proof of claim.}\space}{\hfill$\diamond$\medskip\par}

   \newtheoremstyle{others}% name
     {3pt}%      Space above
     {2pt}%      Space below
     {}%         Body font
     {}%         Indent amount (empty = no indent, \parindent = para indent)
     {\bf}% Thm head font
     {.}%        Punctuation after thm head
     {.5em}%     Space after thm head: " " = normal interword space;
           %       \newline = linebreak
     {}%         Thm head spec (can be left empty, meaning `normal')

\theoremstyle{others}
\newtheorem{rmk}[thm]{Remark}
\newtheorem*{rmk*}{Remark}
\newtheorem{defn}[thm]{Definition}
\newtheorem{example}[thm]{Example}
\newtheorem{examples}[thm]{Examples}
\newtheorem{exercise}[thm]{Exercise}
\newtheorem*{question*}{Question}
\newtheorem{question}{Question}

\renewcommand{\thepart}{\Roman{part}} 

\numberwithin{equation}{section}

\setcounter{tocdepth}{1}



%%%%% Anuk's commands

\usepackage{mathtools} 
\usepackage{faktor}
\usepackage{dynkin-diagrams}
\usepackage{mathrsfs}
\usepackage{hyperref}
%\usepackage[all]{xy}
\usepackage{tikz-cd}
\usepackage{bbm}
%\usepackage{latexsym}
%\usepackage{upgreek}
%\usepackage[version=4]{mhchem}

\newcommand{\oo}{\infty}
\newcommand{\into}{\lrcorner}
\newcommand{\inj}{\hookrightarrow}
\newcommand{\surj}{\twoheadrightarrow}
\newcommand{\knp}{\mathbin{\bigcirc\mspace{-15mu}\wedge\mspace{3mu}}}
\newcommand{\op}[1]{\operatorname*{#1}}
\newcommand{\tbf}[1]{\textbf{#1}}
\newcommand{\mbf}[1]{\mathbf{#1}}
\newcommand{\mbb}[1]{\mathbb{#1}}
\usepackage{romanbar}
\newcommand{\rb}{\Romanbar}
\newtheorem{definition}[thm]{Definition}
\newtheorem{proposition}[thm]{Proposition}
\newtheorem{theorem}[thm]{Theorem}
\newtheorem*{thm*}{Theorem}









\begin{document}

\title[Parabolic gap theorems]{Parabolic gap theorems for the Yang-Mills energy}
\author{Anuk Dayaprema}
\address{University of Wisconsin, Madison}
\email{dayaprema@wisc.edu}
\author{Alex Waldron}
\address{University of Wisconsin, Madison}
\email{waldron@math.wisc.edu}
\date{\today}



\begin{abstract}
    We prove parabolic versions of several known gap theorems in classical Yang-Mills theory. %that do not require
    %without %requiring
    %assuming the existence of
    %which do not assume the existence of a Yang-Mills connection {\it a priori}. These follow from smooth convergence properties of Yang-Mills flow with small initial curvature, in several senses.
    On an $\SU(r)$-bundle of charge $\kappa$ over the 4-sphere, we show that the space of all connections with Yang-Mills energy less than $4 \pi^2 \left( |\kappa| + 2 \right)$ deformation-retracts under Yang-Mills flow onto the space of instantons, allowing us to simplify the proof of Taubes's path-connectedness theorem.
    On a compact quaternion-K\"ahler manifold with positive scalar curvature, we prove that the space of pseudo-holomorphic connections whose $\mathfrak{sp}(1)$ curvature component has small Morrey norm deformation-retracts under Yang-Mills flow onto the space of instantons. On a nontrivial bundle over a compact manifold of general dimension, we prove that the infimum of the scale-invariant Morrey norm of curvature is positive.
\end{abstract}

\maketitle


\begin{comment}
\begin{abstract}
    This note provides versions of several known gap theorems in classical Yang-Mills theory %that do not require
    without requiring
     a solution of the Yang-Mills equation {\it a priori}. For a topologically nontrivial bundle over a compact simply-connected Riemannian $n$-manifold, we prove that the infimum of the $L^{n/2}$-norm of curvature (or indeed of the scale-invariant Morrey norm) is positive. On a quaternion-K\"ahler manifold with positive scalar curvature, we prove that a pseudo-holomorphic connection with nearly minimal energy deforms smoothly to an instanton. These follow from existence and convergence properties of Yang-Mills flow with small initial curvature, in several senses.
\end{abstract}
\end{comment}


\tableofcontents


\thispagestyle{empty}


\section{Introduction}

%\subsection{Background} 

%Mathematical gauge theory %, particularly as a branch of geometric analysis,
%\subsection{Background}
Classical Yang-Mills theory is concerned with two closely-related %both first- and second-order
semi-elliptic partial differential equations involving a \emph{connection} (known to physicists as a \emph{gauge field}) on a vector or principal bundle over an orientable Riemannian manifold. Solutions of the first are known as \emph{instantons}, while solutions of the second are known as \emph{Yang-Mills connections}. These equations are of first and second order, respectively, in the connection. %\footnote{This article has no associated data.} %, in the connection. % $A.$ %, in the connection. %Solutions of these first-order systems are referred to as \emph{instantons}, % ; these are the central geometric objects of the field. 
%while solutions of the second-order system are referred to as \emph{Yang-Mills connections}.
%In all cases of interest, the instanton equation directly implies the Yang-Mills equation. %; indeed, instantons typically minimize the above functional.
%Meanwhile, the %question of whether the Yang-Mills equation implies the instanton equation %is much more subtle. %, and
%The converse is more subtle.
%The converse is more subtle, and has generated a significant body of work.

In the 4-dimensional setting, an instanton $A$ satisfies the %first-order
\emph{anti-self-duality equation}
$$F_A^+ = 0,$$
up to reversing orientation.
%which is a first-order PDE in the connection.
%In general dimension, 
A Yang-Mills connection satisfies the \emph{Yang-Mills equation} %the second-order equation
$$D_A^*F_A = 0.$$
%The latter %states that $A$ is a critical point of
%which is of second order in the connection; in fact,
The latter is the Euler-Lagrange equation of the \emph{Yang-Mills energy functional}
$$\YM(A) = \frac12 \int_M |F_A|^2\, dV.$$
Since an instanton minimizes $\YM(\cdot)$ among all connections on the given bundle, it must be Yang-Mills.
%Instantons are Yang-Mills; indeed, they minimize the Yang-Mills functional on the given bundle. %solve the Yang-Mills equation.
However, %the question is then whether any higher critical points exist.
the converse fails in general; for instance,
%we know from 
by the celebrated work of Sibner, Sibner, and Uhlenbeck
\cite{sibnersuhlenbeck}, we know that even the trivial $\SU(2)$-bundle over $S^4$ carries a vast supply of non-minimal Yang-Mills connections. 

%It is nevertheless desirable to determine circumstances guaranteeing that a solution of the Yang-Mills equations will be an instanton. %a solution 
%the (second-order) Yang-Mills equation are instantons. %must be an instanton, as needed for geometric or physical applications.
Positive results in the converse direction are nonetheless common in the literature. Since the work of Bourguignon, Lawson, and Simons \cite{bourguignonlawson, bourguignonlawsonsimons}, these have fallen into two categories: \emph{stability theorems}, stating that a Yang-Mills connection with nonnegative second variation must be an instanton, and \emph{gap theorems}, stating that a Yang-Mills connection with appropriately small self-dual curvature
must be an instanton. %We are concerned here with the latter. 
%This paper is concerned with generalized gap results of a certain type; we first review several existing gap theorems. %in a sense we will explain below.
We begin by reviewing a number of existing gap theorems.


\begin{comment}
This paper is concerned with strengthened gap theorems of a certain type; %This note is concerned with the latter.
%The typical positive result takes the form of a ``gap theorem,'' stating that a Yang-Mills connection with nearly minimal energy must be an instanton.
to begin, we review several standard gap theorems from the literature.
\end{comment}



The $L^2$ gap theorem on $S^4,$ due to Min-Oo \cite{minoogap}, runs as follows. Suppose that $A$ is
a Yang-Mills connection on a vector bundle $E \to S^4$ with
$$\YM(A) < 4 \pi^2 \left( |\kappa(E)| + \delta \right),$$
where $\delta$ is a positive universal constant and $\kappa(E)$ is the characteristic number defined in \S \ref{ss:characteristicnumber} below; then %$F_A^+ \equiv 0,$ {\it i.e.},
$A$ is an instanton. See Donaldson-Kronheimer \cite[\S 2.3.10]{donkron} for a proof, and Dodziuk and Min-Oo \cite{dodziukminoo} for an improved constant.
%The constant $\delta$ % > 0$ 
%in the %Min-Oo's 
%original proof can be made explicit and is less than one.
More recently, in the case that $E$ has structure group $\SU(r)$, Gursky, Kelleher, and Streets \cite{gurskykelleherstreets} have managed to improve the value to $\delta = 2,$ or twice the Yang-Mills energy of the unit-charge instanton. %This is also the first energy where an anti-self-dual, so is the optimal estimate obtainable by general techniques. 
This is consistent with known examples of non-minimal Yang-Mills connections on $S^4$ \cite{bor, sadunsegert, sibnersuhlenbeck}.


%\begin{comment}
More generally, the existence of an $L^2$ gap is known for Riemannian 4-manifolds where the scalar curvature dominates the self-dual Weyl tensor \cite{gerhardtgap, parker4Dgaugetheories},  %parkerinventiones}, 
for compact K\"ahler surfaces with positive scalar curvature \cite{huangkahlersurfaces}, for bundles where the infimum of the Yang-Mills energy is zero \cite{feehanflatgap}, and for certain bundles over compact 4-manifolds with generic metrics %and structure group $\SU(2)$
\cite{feehanfourgap} (based on the Freed-Uhlenbeck Theorem \cite[Theorem 3.17]{freeduhlenbeck}).
%or \cite[Theorem 4.3.17]{donkron}).
Under mild geometric constraints, Vieira extended the Gursky-Kelleher-Streets result to noncompact manifolds with positive Yamabe constant \cite{vieirapositiveyamabegap}. Interestingly, the existence of an $L^2$ gap on a general compact Riemannian 4-manifold remains unknown. %\footnote{An 89-p. ArXiv posting from May 2015 claims to prove the more general statement that the set of critical values of the Yang-Mills functional on a compact 4-manifold is discrete; however, the proof strategy followed there is incorrect.}


\begin{comment}
The Min-Oo gap theorem has been generalized by Gerhardt to 4-manifolds with appropriately positive curvature. The gap theorem is also easy to prove in the case of a flat bundle, and for a generic metric, by appealing to the Freed-Uhlenbeck transversality theorem.
Meanwhile, for bundles with nontrivial first Pontryagin class over general compact Riemannian 4-manifolds, the gap question remains open. %, despite an unresolved claim from May 2015.
%This again demonstrates the subtlety of the problem.
\end{comment}



The notion of instanton can be extended to a higher-dimensional manifold endowed with an $(n-4)$-calibration %, 
$\Psi$ \cite{corrigandevchand, tiancalibrated}: a connection \(A\) is called a (\(\Psi\)-)\emph{instanton} if
\begin{align*}
    \ast (F_A \wedge \Psi) = -F_A.
\end{align*}
In particular, manifolds with special holonomy %(Berger's list)
come with such calibrations, %instantons
and the gap question is basic in higher-dimensional gauge theory. 
No firm gap results exist for $\G_2$- or $\Spin(7)$-manifolds, and the question appears to be a difficult one. However, on a quaternion-K{\"a}hler manifold with positive scalar curvature, Taniguchi established an \(L^\infty\) gap \cite{qksupgap} for pseudo-holomorphic (see \cite[\textsection 3.3]{oliveirawaldron}) Yang-Mills connections, and later an \(L^{\frac{n}{2}}\) gap \cite{qkL2gap}. The \(L^\infty\) result can be seen as a generalization of Bourguignon-Lawson-Simons's %'s 
theorem on \(S^4\), while the \(L^{\frac{n}{2}}\) result can be seen as a generalization of Min-Oo's theorem. \par 




Although higher-dimensional manifolds with generic holonomy lack a natural class of minimizers for \(\mathcal{YM}\) other than the flat connections, the gap question can still be posed at the zero energy.
Bourguignon-Lawson \cite{bourguignonlawson} proved an \(L^\infty\) gap theorem for the full curvature of a connection on a bundle over \(S^n, n \geq 4\). Gerhardt \cite{gerhardtgap} proved an \(L^{\frac{n}{2}}\) gap theorem for compact manifolds where the Riemann curvature satisfies a certain pointwise positivity condition. Zhou \cite{zhougap} has obtained \(L^p\), \(p \geq \frac{n}{2}\), gaps which depend on the constants in the global Sobolev inequality for functions on the manifold. The existence of a (non-explicit) $L^2$ gap near the zero energy is easily proved by combining standard estimates with the well-known \L ojasiewicz-Simon inequality; see Corollary \ref{cor:ellipticgap} %and Remark \ref{rmk:ellipticgap}
below. %For general compact manifolds, the Lojasiewicz-Simon inequality, due to B. Yang \cite{byanguniqueness} in the Yang-Mills setting, and standard analytic facts about Yang-Mills connections in fact imply an \(L^2\) gap, albeit with very small and non-explicit constant; c.f. the \(L^{\frac{n}{2}}\) gap of Feehan \cite[version 8]{feehangap}.



%Despite their prominence in gauge theory, 
%Notice that these gap theorems all share an obvious drawback: %that
%A fundamental drawback of these gap theorems is that
Notice that these gap results all share an obvious drawback, namely, that the connection in question must solve the Yang-Mills equation {\it a priori}. %share the assumption that %these gap theorems all share a limiting assumption: namely, that
 % in the first place.
Although Uhlenbeck's direct method succeeds at producing solutions
in dimension less than four, the Yang-Mills equation is itself extremely difficult to solve in dimension at least four. This is due both to the gauge freedom of the Yang-Mills functional and to the famous ``bubbling phenomenon'' which begins in dimension four and is less-well understood in higher dimensions.
\begin{comment} This is due both to the gauge freedom of the above functional and to the so-called ``bubbling phenomenon;'' meanwhile, from the point-of-view of geometric analysis, these are what make the subject so interesting.
\end{comment}
%Hence, elliptic gap theorems have somewhat limited scope.

The goal of this paper is to broaden the scope of these results by removing the %\emph{a priori}
assumption that the relevant connection be Yang-Mills. As an alternative,
we consider solutions of the Yang-Mills \emph{flow}:
\begin{equation}\label{ymf}
\frac{\p A}{\p t} = - D_A^*F_A.
\end{equation}
%starting from a general connection with energy below the same threshhold. 
%The goal is to establish that (\ref{ymf}) 
%It is natural to approach the gap question from the perspective of Yang-Mills \emph{flow}; this is the downward gradient flow of the Yang-Mills functional, given by
%\begin{equation}\label{ymf}
%\frac{\p A}{\p t} = - D_A^*F_A.
%\end{equation}
In the same circumstances 
as several of the elliptic gap theorems discussed above, we will 
establish that (\ref{ymf})
gives a deformation retraction from an appropriate 
neighborhood in the space of all connections 
%(modulo gauge)
onto the %moduli
subspace of instantons. A result of this kind will be called a \emph{parabolic gap theorem}.

We will also demonstrate a few applications of our main results.


\vspace{5mm}


\begin{comment}

The motivation for developing parabolic gap theorems comes from three directions. %The first is analytic: 
First, a parabolic gap is in fact a nontrivial analytic improvement over an elliptic gap; the former can fail even while the latter holds. %, both in dimension four and (more dramatically) in higher dimensions.
Next, a parabolic gap allows for a direct comparison between the topology of the space of all connections (modulo gauge) %, which can be computed by homotopy-theoretic methods,
and that of the instanton moduli space; this %is \emph{a priori} quite mysterious. This 
is a classical topic going back to Atiyah-Jones \cite{}.
%The last is geometric:
Last, when existence and smooth convergence of (\ref{ymf}) can be established on a special-holonomy manifold, the flow becomes a tool to construct new instantons. This is a goal for future work.

We state and discuss our results over the next three subsections.


\end{comment}

%Our results %and discuss several applications 
%are stated over the next three subsections.


\begin{comment}

We first note that for elementary reasons, the strong gap may fail even while a classical gap holds:



%we will show that a general connection with appropriately small curvature deforms smoothly to an instanton under the flow. %In this way, we need not begin with a solution of the Yang-Mills equation in order to solve the instanton equation.
In this sense, we remove the assumption %from several existing gap theorems
that a Yang-Mills connection exist {\it a priori} on the given bundle. %, giving us a stronger method to construct instantons. %; hence, in theory, we need not begin with a solution of the Yang-Mills equation in order to solve the instanton equation. %In the opposite direction, we can also prove certain non-existence results for general connections with small curvature on a given bundle.
%This approach cuts in two directions: first,

We will also give several applications.
%Meanwhile
Meanwhile, in cases where instantons cannot exist for topological reasons, we %learn that a general connection with small curvature also cannot exist; this gives us 
obtain a gap statement for the infimum of the appropriate curvature norm on the space of all connections. % on the given bundle. \par

\textbf{Theorem 7.3.3 in Lin-Wag book: Gap theorem for harmonic map flow?}


%We now state our results.
%Our results are as follows.
We summarize our results here. %Since they mostly follow from the author's previous work and that of Dayaprema \cite{}, the proofs will scarcely be longer than the statements.

%On the other hand, we learn that if a connection with energy close to the required energy does exist, then it deforms under Yang-Mills flow to an instanton. We expect the latter type of result to be useful for constructing instantons in certain higher-dimensional circumstances.







%We now state our results.



\end{comment}



\subsection{Deformation to instantons on the 4-sphere}
 %We note that the optimal gap value, {\it i.e.} the lowest nontrivial critical value of the Yang-Mills functional on $S^4,$ remains an intriguing open question.

Our first result is a parabolic gap theorem on $S^4$ based on Gursky-Kelleher-Streets's improvement of the Min-Oo gap theorem and the second-named author's thesis. %and
%the improved gap theorem, we can prove the following.


\begin{thm}\label{thm:fourspheregap}
Let $A_0$ be a smooth connection on an $\SU(r)$-bundle $E \to S^4$ with energy
\begin{equation}\label{fourspheregap:energylessthan8pi2}
\YM(A_0) < 4 \pi^2 \left( |\kappa(E)|+ 2 \right).
\end{equation}
If $A(t)$ solves (\ref{ymf}) with $A(0) = A_0,$ then $A(t)$ converges smoothly to an instanton on $E$ as $t \to \infty.$ Moreover, the map $A_0 \mapsto \lim_{t \to \infty} A(t)$ gives a deformation retraction from the space of connections satisfying (\ref{fourspheregap:energylessthan8pi2}) onto the space of instantons. %, with or without modding out by gauge.
\end{thm}

Note that if we ask only for Uhlenbeck convergence along a subsequence, % in the statement,
then the result follows directly from \cite{gurskykelleherstreets} and \cite{lte}; % \cite{lte} and \cite{gurskykelleherstreets}; 
the main point is that the convergence takes place smoothly. %; the proof of convergence is based on \cite{instantons}.
We also note that the statement is true both before or after modding out by the gauge group.

The proof depends on an optimal ``split'' $\eps$-regularity theorem for the self-dual curvature along (\ref{ymf}), which is proved using a type-II blowup argument appealing to Uhlenbeck's compactness and removable singularity theorems. %This gives a supremum bound on the self-dual curvature along the flow.
By the key estimate of \cite{instantons}, this guarantees not only long-time existence but smooth convergence when working %with the round metric 
on $S^4.$
Theorem \ref{thm:fourspheregap} can also be compared with \cite[Theorem 1.3]{gurskykelleherstreets}, which gives the same result on the trivial bundle. %, where the splitting trick is not necessary.

As an application of Theorem \ref{thm:fourspheregap}, we can give a streamlined proof of the following classic theorem of Taubes.
\begin{thm}[Taubes \cite{taubespathconnected}, 1984]\label{thm:pathconnected}
For any $\SU(r)$-bundle $E \to S^4,$ the moduli space of instantons on $E$ is path-connected.
\end{thm}
\noindent In our variant, described in \S \ref{ss:pathconnected}, all of the hard analysis in Taubes's proof is replaced by Theorem \ref{thm:fourspheregap}. %; in this sense, our proof is simpler.
We can also remove the requirement that $r = 2$ or $3,$ since the GKS gap theorem \cite{gurskykelleherstreets} renders the use of Taubes's index bounds \cite{taubesstability} unnecessary. For general $r,$ connectedness of the $\SU(r)$ moduli spaces on $S^4$ was previously observed by Donaldson \cite[p. 454]{donaldsoninstantonsandgit} using his algebraic description of framed instantons.
%{\color{red} Note that Taubes also proves the result for $\SU(3)$-bundles, but since the GKS gap is not known in that case, we do not obtain a substantial simplification.}

Last, it is interesting to point out that Theorem \ref{thm:pathconnected} provides a purely differential-geometric proof of the completeness of the ADHM construction \cite{adhm}%, originally established by algebro-geometric methods
---see Corollary \ref{cor:adhm} below. We include this observation (without claiming any originality)
in order to highlight the significance of Taubes's theorem. % \cite{taubespathconnected}. %Theorem \ref{thm:pathconnected}.



 %Methods based on Yang-Mills flow should are hoped to yield further results about instanton moduli spaces in the future. %\footnote{Taubes was  too modest to point this out in \cite{taubespathconnected}.}

%Note that this result is a generalization of Corollary \ref{} of Gusrsky-Kelleher-Streets.
%Notice that the bundle in this statement can be nontrivial, by contrast with Corollary bla of \cite{}. Also note that since our normalizations differ by a factor of 4, the value on the RHS is twice the instanton energy.

%  This follows by combining the convergence theorem in my thesis with the gap theorem of Gursky-Kelleher-Streets.


%We can also use this result to give a streamlined proof of the ADHM Theorem \cite{ADHM} following an argument due to Taubes \cite{taubespathconnected}. This can be summarized as follows...







\vspace{5mm}
 
\subsection{Deformation to instantons on quaternion-K\"ahler manifolds}

Let $M$ be a quater-nion-K{\"a}hler manifold of dimension $4m$. For \(m \geq 2\), this means a Riemannian manifold whose holonomy lies in
$$\Sp(m)\Sp(1) \subset \SO(4m).$$ When \(m = 1\), we further impose that \(M\) be Einstein %with nonzero scalar curvature 
and half-conformally flat. %with vanishing self-dual Weyl curvature. %, which is automatic for $k \geq 2$ (see \cite{salamonqkinventiones}). \par

On a quaternion-K{\"a}hler manifold, we have the orthogonal decomposition \cite[\textsection 3.3]{oliveirawaldron}
\begin{align*}
    \Lambda^2 T^\ast M &= \Lambda_{\mathfrak{sp}(m)}^2 \oplus \Lambda_{\mathfrak{sp}(1)}^2  \oplus \Lambda_{(\mathfrak{sp}(m) \oplus \mathfrak{sp}(1))^\perp}^2 \\
    &=: \Lambda_-^2 \oplus \Lambda_+^2 \oplus \Lambda_\perp^2.
\end{align*} Hence, given a connection \(A\), the curvature decomposes as
\begin{align*}
    F_A = F_A^- + F_A^+ + F_A^\perp.
\end{align*} In this paper, an \emph{instanton} on a quaternion-K{\"a}hler manifold is a connection satisfying \(F_A = F_A^-\), i.e., whose curvature takes values in \(\Lambda_-^2(\mathfrak{g}_E)\). These are minimizers of the Yang-Mills functional and hence are Yang-Mills themselves. Connections whose curvatures take values either in \(\Lambda_+^2(\mathfrak{g}_E)\) or in \(\Lambda_{\perp}^2(\mathfrak{g}_E)\) are sometimes also referred to as instantons in the literature. \par 

More generally, a \emph{pseudo-holomorphic connection}, i.e., an \(\text{Sp}(m)\text{Sp}(1)\)-compatible connection per Oliveira-Waldron \cite[Def. 2.1]{oliveirawaldron}, is a connection whose curvature takes values in 
$$\Lambda_{-}^2(\mathfrak{g}_E) \oplus \Lambda_{+}^2(\mathfrak{g}_E).$$
We have the following result, which can be seen as a generalization of Theorem \ref{thm:fourspheregap}. Here the gap is measured in the $M^{2,4}$ Morrey norm (defined in (\ref{preliminaries:morreynormdefinition}) below), which is a natural generalization of the conformally invariant \(L^2\) norm in four dimensions.

%The \(L^2\)-norm on 2-forms has the special feature of being conformally invariant in four dimensions. A natural generalization in higher dimensions is the $M^{2,4}$ Morrey norm defined in (\ref{preliminaries:morreynormdefinition}) below. 

\begin{thm}\label{thm:quaternionkahlergap}
Let $M$ be a compact quaternion-K\"ahler manifold of dimension \(4m\). There exists a constant $\delta_0 > 0,$ depending only on the geometry of $M,$ as follows.

Suppose that $A_0$ is a pseudo-holomorphic connection on $E \to M,$  with %\footnote{See the preliminaries for a definition of the \(M^{2, 4}\) Morrey norm.} 
\begin{align}\label{quaternionkahler:energylessthandelta}
    \left\|F_{A_0}^+ \right\|_{M^{2, 4}} < \delta_0.
\end{align}
Then the solution \(A(t)\) of (\ref{ymf}) with $A(0) = A_0$ exists for all time. If $M$ has positive scalar curvature, then $A(t)$ converges smoothly as \(t \to \infty\) to an instanton $A_\infty,$ satisfying $F_{A_\infty} = F_{A_\infty}^{-}.$ Moreover, the space of %\(\mathfrak{sp}(m)\)-
instantons is a neighborhood deformation retract (see Definition \ref{defn:nbddefret} below) of the space of pseudo-holomorphic connections.
%Moreover, the map $A_0 \mapsto \lim_{t \to \infty} A(t)$ gives a neighborhood deformation retraction from the space of connections satisfying (\ref{quaternionkahler:energylessthandelta}) onto the space of \(\mathfrak{sp}(k)\)-instantons.
\end{thm}












\vspace{5mm}

\subsection{Deformation to flat connections on general Riemannian manifolds}
We have the following parabolic gap statement for the Morrey norm of the full curvature. 
\begin{thm}\label{thm:flatgap}
Let %$G$ be a compact Lie group and
\(M\) be a compact %, oriented 
Riemannian manifold of dimension $n \geq 4.$ There exists a constant $\delta_1 > 0,$ depending only on $\rk(E)$ and the geometry of $M,$ as follows. Suppose that \(A_0\) is a connection on \(E \to M\) with
%Given any %$H^k$ 
%connection $A_0$ with
\begin{equation}\label{flatgap:Morreysmall}
\| F_{A_0} \|_{M^{2, 4}} < \delta_1.
\end{equation}
Then the solution $A(t)$ of (\ref{ymf}) with $A(0) = A_0$ exists for all time and converges smoothly as \(t \to \infty\) to a flat connection. 
% as \(t \to \infty\).
%\begin{equation}
%    \| F_{A(t)} \|_{L^\infty} < \left( C_M t^{-1} + C_? t^{-\frac{1}{1-2\tilde{\theta}}}\right) \delta.
%\end{equation}
\begin{comment}
There exists an $H^{k + 1}$ gauge transformation $u$ and a smooth, flat connection $A_\infty$ on $E$ such that
for any $s \in \N$ and $t \geq 1,$ we have
\begin{equation}\label{detailedMorreygap:polynomialconvergence}
    \| u(A(t)) - A_\infty \|_{H^s} < C_{M,s} \delta t^{-sthg}.
\end{equation}
\end{comment}
Moreover, the space of flat connections is a neighborhood deformation retract of the space of all connections.
%Moreover, the flow gives a neighborhood deformation retraction in the $H^k$ topology from the space of all connections satisfying (\ref{detailedMorreygap:laterMorreysmall}) onto the space of flat connections.
\end{thm}

%These results will follow from known estimates for Yang-Mills flow in dimension four, and 
We include two proofs of this result, the first based on the parabolic Morrey estimates of the first-named author \cite{estimatespaper} and the second based on Hamilton's monotonicity formula for Yang-Mills flow \cite{hamiltonmonotonicity}. %, or from the epsilon-regularity theorem of Chen-Shen%. 
%We also have 
Theorem \ref{thm:flatgap} has the following consequence:

\begin{cor}
Let $E \to M$ be a vector bundle over a compact %, oriented 
Riemannian manifold of dimension $n \geq 4$ and let $\mathcal{A}_E$ denote the space of all connections on $E.$ %There
There exists a flat connection on \(E\) if and only if
\begin{align*}
    \inf\left\{\|F_A\|_{M^{2, 4}} \mid A \in \mathcal{A}_E \right\} = 0.
\end{align*}
\begin{comment}
Suppose that $M$ is simply connected and $P \to M$ is not topologically trivial (or, for general $M,$ that $P$ does not admit a flat connection). We then have
$$\inf_{A \in \calA_P} \|F_A\|_{M^{2, 4}} > \delta_0.$$
Here $\eps > 0$ is the constant of the previous Theorem.
\end{comment}
\end{cor}

%We note that in case $P$ has a nonvanishing Chern or Pontryagin class, positivity of the infimum follows directly from Chern-Weil theory. %Theory. 
%However, in higher dimensions, characteristic classes do not detect triviality, so the statement is less obvious.
The corresponding result for maps between compact Riemannian manifolds is \cite[Corollary 1.4]{estimatespaper}, which %generalizes 
is related to results of White \cite[Cor. 2 and 3 on pg. 128]{whiteinfima} on the infimum of the $L^{n}$-norm of the derivative on the space of maps. % between closed Riemannian manifolds.

By contrast, work of Naito \cite{naito} on the existence of finite-time blowup for (\ref{ymf}) on spheres of dimension $n \geq 5$ implies that there does not in general exist a parabolic \(L^2\) gap. The method of \cite{naito} can also be found in prior work of Chen and Ding in the setting of harmonic map flow \cite[Theorem 1.1]{chending}. We recover the following version of Naito's result for arbitrary compact base manifolds, proved by Wang and Zhang \cite{wangzhangfinitetimeblowupymf}.

\begin{cor}[{\cite[Theorem 1.1]{wangzhangfinitetimeblowupymf}}]\label{cor:Naito}
Let $G$ be a compact Lie group. Supposing that $\pi_1(M)$ has no nontrivial representations in $G,$ there exists a constant \(\delta_2 > 0\), depending only on $G$ %the bundle \(P\)
and the geometry of \(M\), as follows. Given any connection $A_0$ on a topologically nontrivial vector bundle \(E \to M\) with structure group $G,$ if 
\begin{equation}\label{smallym}
\YM(A_0) < \delta \leq \delta_2,
\end{equation}
then the maximal time of smooth existence of the solution of (\ref{ymf}) with $A(0) = A_0$ is less than $\delta_2^{-1}\delta^{\frac{2}{n-4}}$.
%In particular, supposing that the structure group is bla or bla, and possibly under some topological assumption, there always exists a finite-time blowup. (Just need to check whether you can always obtain a nontrivial bundle by grafting in one from the sphere, because then you can obviously make a connection with small YM energy on a nontrivial bundle.) %
\end{cor}

%As observed by Naito \cite{naito}, any bundle over $S^n,$ $n > 4,$ carries a connection satisfying (\ref{smallym}).


\vspace{5mm}




\subsection{Acknowledgements} A.W. was partially supported by NSF DMS-2004661 during the preparation of this article. A.D. was supported by NSF DMS-2037851 during the preparation of this article.

\vspace{10mm}


\section{Preliminaries}\label{section:Preliminaries}

\subsection{Principal and vector bundles} Let $G$ be a compact Lie group and $\rho : G \to \GL_K(V)$ a faithful representation of dimension $r$ over $K = \R$ or $\C.$ We assume that $V$ is endowed with a (Hermitian) inner product preserved by $G,$ so that $\rho(G) < \O(r)$ if $K = \R$ and $\rho(G) < \U(r)$ if $K = \C.$ Let $\gothg$ denote the Lie algebra of $G$, which we identify with the space of left-invariant vector fields on \(G.\)

Let $M$ be a smooth manifold. Recall that a \emph{principal $G$-bundle over $M$} is a smooth fiber bundle $G \inj P \to M$ with fiber \(G\) and a right $G$-action \(P \times G \to P\) satisfying the following conditions. The action should be smooth, commute with the fiber bundle projection, and act freely and transitively on the fibers. Given a principal $G$-bundle $P \to M$ and a space $F$ on which $G$ acts on the left, we may form the \emph{associated fiber bundle}
$$P \times _G F = \{ P \times F \} / \left( x,y \right) \sim \left( x g, g^{-1} y \right).$$
%which is a fiber bundle with fiber $F.$
We form the following associated fiber bundles:
\begin{enumerate}
\item The bundle of \emph{gauge transformations}:
$$\mathcal{G}_P = P \times_{\Ad} G.$$
By abuse of notation, we also refer to its sheaf of smooth sections by $\mathcal{G}_P;$ these correspond to local automorphisms of $P.$

\item The bundle of \emph{infinitesimal gauge transformations} (a.k.a. the \emph{adjoint bundle}):
$$\gothg_P = P \times_{\mathrm{Ad}} \gothg.$$
This is a vector bundle with fiber $\gothg.$

\item The vector bundle
$$E = P \times_\rho V,$$
which has rank $r$ over $K$ and inherits a bundle metric from the inner product on $V.$
\end{enumerate}

The maps $\rho : G \to \GL_K(V)$ and \(d \rho : \mathfrak{g} \to \mathfrak{gl}_K(V)\) yield embeddings of $\mathcal{G}_P$ and $\gothg_P,$ respectively, into $\End E = E \otimes_K E^*$. Hence, an (infinitesimal) gauge transformation of $P$ is equivalent to a section of $\End E$ that lies fiberwise in the image of $\rho(G)$ (resp. $\rho(\gothg)$). %and an infinitesimal gauge transformation is equivalent to a skew-symmetric (resp. skew-Hermitian) section of $E \otimes E^*.$
We endow both $\mathcal{G}_P$ and $\gothg_P$ with the induced metric from $\End E$; in particular, sections of the former are orthogonal (resp. unitary), while sections of the latter are skew-symmetric (resp. skew-Hermitian), for $K = \R$ or $\C,$ respectively.

The principal bundle $P$ can be reconstructed from the associated bundle $E$ by applying the inverse of the map $\rho$ to the transition functions. In fact, any real or complex vector bundle whose transition functions belong to the image of a representation of $G$ arises from the associated bundle construction for some principal \(G\)-bundle. We may therefore define a %rank $r$ real or complex 
\emph{vector bundle with structure group $(G, \rho)$}, or \emph{$(G,\rho)$-bundle}, in the following two equivalent ways:

\begin{enumerate}

\item The associated bundle $E = P \times_\rho V$ for a principal $G$-bundle $P$ and $r$-dimensional real or complex representation, $V,$ or

\item A vector bundle of rank $r$ over $\R$ or $\C$ and a cover of \(M\) by trivializations for which the transition functions belong to $\rho(G).$

\end{enumerate}
The choice of representation $\rho$ is usually obvious; for example, a rank $r$ real vector bundle ``with structure group $\SO(r)$'' has all transition functions in the image of the standard rank $r$ real representation, while a rank $r$ complex vector bundle ``with structure group $\SU(r)$'' has all transition functions in the image of the standard rank $r$ complex representation.

In this paper, we prefer to work with vector rather than principal bundles, which entails no loss of generality since $G$ is compact. Henceforth, for a $(G,\rho)$-bundle $E,$ we will denote the images of $\mathcal{G}_P$ and $\gothg_P$ inside $\End E$ by $\mathcal{G}_E$ and $\gothg_E,$ respectively.

%A {\it gauge transformation} is a (local) section of \(\mathcal{G}_P\).  The associated vector bundle \(\mathfrak{g}_P := P \times_\rho \mathfrak{g}\) is called the {\it bundle of infinitesimal gauge transformations}, or the {\it adjoint bundle}.

\vspace{5mm}

\subsection{Connections, curvature, and gauge transformations}

Recall that a connection on a principal \(G\)-bundle \(\pi : P \to M\) is a smooth distribution \(\mathcal{H} \subset TP\) which satisfies the following:
\begin{enumerate}
    \item \(TP \cong \mathcal{H} \oplus \ker{d\pi}\).
    \item For all \(g \in G\) and \(p \in P\), \((R_g)_\ast(\mathcal{H} \cap T_pP) = \mathcal{H} \cap T_{pg} P\).
\end{enumerate}
%Consequently, 
Given a tangent vector \(X\) at \(x \in M\) and a point \(p \in \pi^{-1}(x)\), there is a unique \(X^H \in \mathcal{H} \cap T_pP\), called the horizontal lift of \(X\), such that \(d\pi(X^H) = X\). \par 
A connection $A$ on \(P\) induces a covariant derivative \(\nabla_A\) on the vector bundle $E = P \times_\rho V \to M$ as follows. Let \(s\) be a smooth section of \(E\). Since \(G\) acts freely and transitively on the fibers of \(P\), \(s\) determines a unique, \(G\)-invariant section \(\tilde{s}\) of \(P \times V \to P\), which may equivalently be viewed as a \(G\)-equivariant function \(P \to V\). Given \(X \in T_xM\), \((d\tilde{s})(X_p^H)\) determines an element of \(\pi^{-1}(x) \times V\), and by \(G\)-equivariance, \((d\tilde{s})(X_p^H)\) lie in the same equivalence class of \(\pi^{-1}(x) \times_\rho V\) as \(p\) varies in \(\pi^{-1}(x)\). Hence we define \(\nabla_A\) by
\begin{align*}
    (\nabla_A s)(X) := [(d\tilde{s})(X^H)].
\end{align*} \par
%Conversely, given a \((G, \rho)\)-bundle \(E\), suppose we have a covariant derivative $\nabla_A$ whose connection forms take values in $d\rho(\gothg)$ in local frames for which the transition functions belong fiberwise to $\rho(G).$ 
Conversely, given a \((G, \rho)\)-bundle \(E\), suppose we have a covariant derivative $\nabla_A$ such that, in local frames for which the transition functions belong fiberwise to $\rho(G)$, the connection 1-forms take values in $d\rho(\gothg)$. We obtain a connection on \(P\) as follows. Let \(U\) be a trivialized neighborhood of \(x \in M\), and let \(X \in T_xM\). By abuse of notation, we denote the connection form with respect to the trivialization by \(A.\) Since \(\rho\) is faithful, we may define a horizontal vector at \((x, e) \in U \times G\) by \((X, (d\rho)^{-1}(A(X)) )\). We obtain a horizontal distribution over \(U \times G\) by right-translation, and since the transition functions of \(E\) belong to \(\rho(G)\) and \(\rho\) is faithful, this prescription glues together to give a connection on \(P\). Thus we refer to such covariant derivatives as \((G, \rho)\)-connections, and in the sequel we restrict our attention to \((G, \rho)\)-connections on \(E\) instead of principal connections on \(P\). \par 
The curvature of a connection \(A\) is a section \(F_A\) of \(\Lambda^2T^\ast M \otimes \mathfrak{g}_E\) defined locally by 
\begin{align*}
    F_A = dA + A \wedge A.
\end{align*} Given another \((G, \rho)\)-connection \(\tilde{A}\), whence \(a = \tilde{A} - A\) is a globally well-defined \(\mathfrak{g}_E\)-valued 1-form, we have
\begin{align*}
    F_{\tilde{A}} = F_A + D_Aa + a \wedge a,
\end{align*}
where $D_A$ denotes the covariant exterior derivative.
A gauge transformation \(u \in \Gamma(\mathcal{G}_E)\) acts on connections by the prescription \begin{align*}
    \nabla_{u(A)} s = u(\nabla_A(u^{-1}(s))).
\end{align*} Locally, the connection form transforms by
\begin{align}\label{preliminaries:connectiontransformationlaw}
    A \mapsto u A u^{-1} - du u^{-1},
\end{align} and the curvature transforms by
\begin{align}\label{preliminaries:curvatureisgaugeequivariant}
    F_{u(A)} = uF_Au^{-1}.
\end{align}

\vspace{5mm}


\subsection{The Yang-Mills functional}

Suppose now that $M$ is endowed with a Riemannian metric $g.$ We define the induced inner product on $k$-forms by requiring that
$$|e^1 \wedge \cdots \wedge e^k| = 1$$
for orthonormal covectors $e^1, \ldots, e^k.$ The Hodge star operator $* : \Omega^k \to \Omega^{n-k}$ is an isometry defined by
\[
\LA \alpha, \beta \RA \, dV = \alpha \wedge * \beta
\]
which extends linearly to forms valued in any vector bundle.
Combining this with the norm on $\End E,$ for $\omega, \eta$ (skew-symmetric/skew-Hermitian) sections of $\Omega^k(\gothg_E),$ we have
\begin{align}\label{preliminaries:forminnerproduct}
\LA \omega, \eta \RA \, dV = - \Tr \omega \wedge * \eta,
\end{align}
where $\Tr = \Tr_K.$ The Yang-Mills functional is defined by 
\begin{align}\label{preliminaries:ym}
   \mathcal{YM}(A) & = - \frac12 \int_M \Tr F_A \wedge * F_A \\
& = \frac12 \int_M |F_A|^2 \, dV . \nonumber
\end{align}



\vspace{5mm}

\subsection{The characteristic number}\label{ss:characteristicnumber}

We follow the convention of Donaldson-Kronheimer \cite[p. 42]{donkron} for the characteristic number $\kappa(E) \in \frac14 \Z$ of a real or complex vector bundle \(E\) over a closed, oriented 4-manifold:
\begin{align*}
    \kappa(E) = \begin{cases} \langle c_2(E) - \frac12 c_1(E)^2, [M]\rangle & K = \C \\
    -\frac14 \langle p_1(E), [M]\rangle & K = \R.
    \end{cases}
\end{align*}
The coefficient $-\frac14$ guarantees %the following 
that for an $\text{SU}(2)$-bundle $E,$ %considered as a complex bundle,
we have $\kappa(E) = \kappa(\gothg_E)$ %where $\gothg_E$ is considered as a real bundle 
\cite[(2.1.37)]{donkron}. %In other words, for an $\text{SU}(2)$-bundle, the characteristic number of the rank-2 complex bundle associated to the standard representation of $\SU(2)$ agrees with that of the rank-3 real bundle associated with the adjoint representation $\SU(2) \to \SO(3).$
\begin{comment}
given an $\SU(r)$-bundle $E,$ considered as a complex bundle, % of rank $r,$ 
we have $\kappa(E) = \kappa(\gothg_E),$ where $\gothg_E$ is considered as a real bundle \cite[(2.1.37)]{donkron}.  % of rank ${r + 1 \choose 2}.
For example, given a principal $\text{SU}(2)$-bundle $P,$ the characteristic number of the rank-2 complex bundle associated to the standard representation of $\SU(2)$ agrees with that of the rank-3 real bundle associated with the adjoint representation $\SU(2) \to \SO(3).$ %This number 
\end{comment}
The number \(\kappa(E)\) has the property that it is an integer over the 4-sphere, according to Atiyah-Hitchin-Singer \cite[pg. 450]{ahs}.

According to Chern-Weil theory \cite[Appendix C]{milnorstasheff}, for any %smooth 
connection $A$ on $E,$ we have
\begin{align*}
    \kappa(E) = \frac{1}{8\pi^2 n_{K}} \int_M \Tr_K F_A \wedge F_A,
\end{align*}
where
\begin{align}\label{preliminaries:nK}
    n_{K} = \begin{cases} 1 & K = \C \\
    4 & K = \R.
    \end{cases}
\end{align} \par 
In four dimensions, the Hodge star operator \(\ast\) acts as an involution on 2-forms. Hence we have the pointwise orthogonal decomposition \(\Lambda^2T^\ast M = \Lambda_-^2 \oplus \Lambda_+^2\) into negative and positive eigenspaces of \(\ast\), and this decomposition extends to \(\mathfrak{g}_E\)-valued 2-forms. Since \(\ast F_A = F_A^+ - F_A^-\), we deduce from (\ref{preliminaries:forminnerproduct}) and (\ref{preliminaries:ym}) that
\begin{align*} 
    \YM(A) = 4 \pi^2 n_{K} \kappa(E) + \|F^+_A\|_{L^2}^2.
\end{align*} \par

Suppose now that \(M\) is a closed quaternion-K{\"a}hler manifold of dimension \(4m\). Let \(\Omega\) denote the Kraines form on \(M\) \cite{kraines}. Since \(d\Omega = 0\), it follows from Chern-Weil theory that the integral
\begin{align*}
    \int_M \text{Tr}_K(F_A \wedge F_A) \wedge \Omega^{m-1}
\end{align*} is a topological invariant of the bundle. A consequence is that for a pseudo-holomorphic connection \(A\), we have 
\begin{align}\label{preliminaries:qkchernweil}
    \mathcal{YM}(A) = \lambda(E) + c\|F^+\|_{L^2}^2,
\end{align} where \(\lambda(E)\) is a constant depending only on the bundle, and \(c\) is a constant depending only on \(m\) \cite[Lemma 2.3 (ii)]{oliveirawaldron}.
\begin{rmk}\label{rmk:gkscomparison}
%{\color{red} Note that in the case $K = \R,$ this convention differs by a factor of $4$ from that of Gursky, Kelleher, and Streets \cite{gurskykelleherstreets}. Our convention for the norm of a 2-form also differs by a factor of $2.$ The overall effect is that for \(G = SU(2)\) and \(\rho\) the standard representation, whereas if \(G = SO(3)\). Our convention is chosen in order that the gap constant $\delta_{G, \rho}$ of Definition \ref{defn:delta} be given in multiples of the unit instanton energy. }
Due to the different choices of norms, our convention for $\kappa(E)$ agrees with that of Gursky-Kelleher-Streets \cite{gurskykelleherstreets} in the complex case but differs by half in the real case. On the other hand, our convention for $\|F\|^2$ is half that of \cite{gurskykelleherstreets} in the complex case and agrees in the real case.
\end{rmk}




\vspace{5mm}

\begin{comment}
    \subsection{Sobolev and Morrey norms}
\end{comment}
\subsection{Sobolev spaces of connections}
 Fix a smooth reference \((G, \rho)\)-connection \(\nabla_{\text{ref}}\) on $E.$ Let \(\nabla_{\text{ref}}\) also denote the induced connection on \(\mathfrak{g}_E\) as well as any tensor product connection obtained by coupling the Levi-Civita connection on a tensor bundle over \(M\) with \(\nabla_{\text{ref}}\) on \(E\) or \(\mathfrak{g}_E\). For \(p \in [1, \infty)\), we define the \(W^{k, p}\) Sobolev norm of a section of \(E, \mathfrak{g}_E,\) or any natural tensor product bundle, by
\begin{align}
    \|s\|_{W^{k, p}} := \left(\sum_{j = 0}^k \int_M |\nabla_{\text{ref}}^{(k)} s|^p \, dV \right)^{\frac{1}{p}}.
\end{align} When \(p = 2\), we abbreviate
\begin{align*}
    \|s\|_{H^k} &:= \|s\|_{W^{k, 2}} \\
    \|s\| &:= \|s\|_{L^2}.
\end{align*}



 % on sections, gauge transformations, and connections



%\begin{align*}
%    \|A\|_{W^{k, p}} := \|\nabla - \nabla_{ref}\|_{W^{k, p}}.
%\end{align*}
\begin{comment}
We next define Morrey spaces of sections of a vector bundle (equipped with a fiber metric) on a Riemannian manifold $M$ of dimension \(n\).
        \begin{defn}\label{preliminaries:morreynormdefinition}
		 Given \(q \in [1, \infty]\) and \(\lambda \in [0, n]\), the \emph{Morrey space} $M^{q, \lambda}$ is the Banach space of measurable sections $s$ for which the Morrey norm $\|s\|_{q, \lambda}$ is finite, where
    \begin{align*}
        \|s\|_{q, \lambda} := 
        \begin{cases}
            \sup_{x \in M, 0 < R \leq 1} \left(R^{\lambda-n} \int_{B_R(x)} |s|^q \, dV\right)^{\frac{1}{q}} & q < \infty \\
            \|s\|_\infty & q = \infty.
        \end{cases}
    \end{align*}
	\end{defn}  \par
\end{comment}
%\begin{comment}
We next define the Morrey norm on a Riemannian manifold $M$ of dimension \(n\).
        \begin{defn}\label{preliminaries:morreynormdefinition}
		 Given \(q \in [1, \infty]\) and \(\lambda \in [0, n]\), the $M^{q, \lambda}$ \emph{Morrey norm} of a function \(f\) is given by
    \begin{align*}
        \|f\|_{q, \lambda} := 
        \begin{cases}
            \sup_{x \in M, 0 < R \leq 1} \left(R^{\lambda-n} \int_{B_R(x)} |f|^q \, dV\right)^{\frac{1}{q}} & q < \infty \\
            \|f\|_\infty & q = \infty.
        \end{cases}
    \end{align*} The Morrey norm of a section \(s\) of a bundle with fiber metric is just the Morrey norm of the function \(|s|\).
	\end{defn}  \par
%\end{comment}
 
\begin{comment}
\vspace{5mm}

\subsection{The Sobolev topology on the space of connections (modulo gauge)}
\end{comment}
To define the Sobolev topology on the space of connections $\calA_E,$ %first fix a smooth reference connections $\nabla_{ref}$ on $E.$ 
recall that the covariant derivative associated to $A$ can be written uniquely as
\[
\nabla_A = \nabla_{\text{ref}} + A
\]
for a \(\mathfrak{g}_E\)-valued 1-form $A$ on \(M\); abusing notation, we denote this 1-form by the same letter as the connection itself. The \(W^{k, p}\) norm of $A$ is simply defined as that of the corresponding 1-form. We may also define the space of Sobolev connections \(W^{k, p}(\calA_E)\) as the completion of the space of smooth connections $\calA_E$ with respect to the \(W^{k, p}\) norm. By compactness of \(M\), the ``Sobolev topology'' on this space is independent of the choice of a smooth reference connection. 
\begin{comment}
Moreover, by Donaldson-Kronheimer \cite[...]{donkron}, for $k$ sufficiently large, the inclusion \(W^{k + 1, p}(\calA_E) \hookrightarrow W^{k, p}(\calA_E)\) is a homotopy equivalence.
\end{comment}
 \begin{rmk}\label{rmk:Hqcontrolsmorrey}
      For \(k \geq \frac{n}{2}-1\), Sobolev embedding implies \(H^{k-1} \inj L^p\) for some \(p \geq \frac{n}{2}\); we may also simultaneously replace both \(\geq\) by \(>\). On a compact manifold, H{\"o}lder's inequality implies that for such \(p\), the $M^{2, 4}$ Morrey norm is bounded above by a constant times the $L^p$ norm.
      %\(L^p \inj M^{2, 4}\). 
      Since the assignment \(A \mapsto F_A\) is a continuous map from \(H^k \to H^{k-1}\) for \(k \geq \frac{n}{2}-1\), it follows that the \(M^{2, 4}\) norm of the curvature is a continuous, gauge-invariant function on \(H^k(\mathcal{A}_E)\).
  \end{rmk}

The transformation law (\ref{preliminaries:connectiontransformationlaw}) implies that %or $p \in (1, \infty)$ and
for $k > \frac{n}{2} - 1,$ the action of the gauge group $\mathcal{G}_E$ on $\calA_E$ extends continuously to the spaces of Sobolev connections \(H^k(\mathcal{A}_E)\) and gauge transformations \(H^{k+1}(\mathcal{G}_E)\). In fact, this is a smooth action by a Hilbert Lie group on a Hilbert manifold; see \cite[\textsection 6 and Appendix A]{freeduhlenbeck} or %. Moreover, we can mod out \(H^k(\mathcal{A}_E)\) by the action of \(H^{k+1}(\mathcal{G}_E)\) to obtain the Hausdorff topological space \(\mathcal{B}_E\)
\cite[\textsection 4.2.1]{donkron}. \par 
Fix a point \(x_0 \in M\). Let $(\mathcal{G}_E)_{x_0}$ denote the fiber of the bundle of gauge transformations over $x_0.$ We define a framed connection to be an element of
\begin{align*}
    \tilde{\mathcal{A}}_E := \mathcal{A}_E \times (\mathcal{G}_E)_{x_0}.
\end{align*}The gauge group acts on \(\tilde{\mathcal{A}}_E\) by \(\sigma\cdot(A, u) = (\sigma(A), \sigma(x_0)u)\). Since \(H^{k+1}\) gauge transformations are continuous for \(k+1 > \frac{n}{2}\), we have a well-defined, smooth %continuous 
action of \(H^{k+1}(\mathcal{G}_E)\) on the \(H^k\) framed connections \(H^k(\tilde{\mathcal{A}}_E)\). Fixing $k > \frac{n}{2} - 1,$ we define the moduli space of framed connections by
$$\tilde{\mathcal{B}}_E:=H^k(\tilde{\mathcal{A}}_E)/H^{k+1}(\mathcal{G}_E).$$
This is a Banach manifold independent of \(k\), in the sense that for $\ell \geq k,$ the induced map \(H^\ell(\tilde{\mathcal{A}}_E)/H^{\ell+1}(\mathcal{G}_E) \to H^k(\tilde{\mathcal{A}}_E)/H^{k+1}(\mathcal{G}_E)\) %are weakly homotopy-equivalent to a common topological space 
is a weak homotopy equivalence \cite[\S 5.1.1]{donkron}. We also let $\calB_E = \tilde{\calB}_E / G$ denote the unframed moduli space of connections.






\vspace{5mm}

\subsection{Weitzenb{\"o}ck formulae} Recall that the adjoint \(D_A^\ast\) of the exterior covariant derivative, acting on \(\mathfrak{g}_E\)-valued \(k\)-forms, is given by
\begin{align*}
    D_A^\ast = (-1)^{n(k+1)+1} \ast D_A \ast.
\end{align*} We define the Hodge Laplacian \(\Delta_A\) by 
\begin{align*}
    \Delta_A = D_A^\ast D_A + D_A D_A^\ast.
\end{align*} The Hodge Laplacian and the connection Laplacian \(\nabla_A^\ast \nabla_A\) differ by algebraic terms involving the curvatures of \(A\) and \(g\). Let $R_{ijkl}, R_{ij},$ and $R$ respectively denote the Riemann, Ricci, and scalar curvature of $g.$ On \(\mathfrak{g}_E\)-valued 2-forms, we have from \cite[(4.6)]{oliveirawaldron} that
\begin{align*}
    (\Delta_A \omega)_{ij} = (\nabla_A^\ast \nabla_A \omega)_{ij} - [F_i^{\ k}, \omega_{kj}] + [F_j^{\ k}, \omega_{ki}] + R_{i \ j}^{\ k \ l}\omega_{lk} - R_{j \ i}^{\ k \ l}\omega_{lk} + R_i^{\ l}\omega_{lj} + R_j^{\ l}\omega_{il},
\end{align*}
\begin{comment}
    \begin{align*}
    (\Delta_A \omega)_{ij} = -\nabla^k\nabla_k \omega_{ij} - [F_i^{\ k}, \omega_{kj}] + [F_j^{\ k}, \omega_{ki}] + R_{i \ j}^{\ k \ l}\omega_{lk} - R_{j \ i}^{\ k \ l}\omega_{lk} + R_i^{\ l}\omega_{lj} + R_j^{\ l}\omega_{il},
\end{align*}
\end{comment}
 where we use the Einstein summation convention of summing over repeated indices. By the algebraic Bianchi identity,
 \begin{align*}
    R_{ikjl} = - R_{lijk} - R_{klji} = R_{jkil} + R_{ijkl},
\end{align*}
 \begin{comment}
\begin{align*}
    R_{i \ j}^{\ k \ l} = - R_{\ ij}^{l \ \ k} - R_{\ \ ji}^{kl} = R_{j \ i}^{\ k \ l} + R_{ij}^{\ \ kl},
\end{align*}
 \end{comment}
 so the Weitzenb{\"o}ck identity takes the form
\begin{align*}
    (\Delta_A \omega)_{ij} = (\nabla_A^\ast \nabla_A \omega)_{ij}
    %-\nabla^k\nabla_k \omega_{ij} 
    - [F_i^{\ k}, \omega_{kj}] + [F_j^{\ k}, \omega_{ki}] - R_{ij}^{\ \ kl}\omega_{kl} + R_i^{\ l}\omega_{lj} + R_j^{\ l}\omega_{il}.
\end{align*} \par
We now specialize the general Weitzenb{\"o}ck formula to the quaternion-K{\"a}hler setting. On a quaternion-K{\"a}hler manifold of dimension \(4m\), Salamon \cite[Theorem 3.1]{salamonqkinventiones} showed that the Riemannian curvature tensor can be decomposed into a scalar multiple of the curvature tensor of \(\mathbb{HP}^m\) plus the curvature tensor of a hyperk{\"a}hler manifold. To be precise, let \({}^1 I, {}^2 I, {}^3 I\) be a local frame of 
%\(\mathfrak{sp}(1) \subset \Lambda^2T^\ast M\)
\(\Lambda_+^2 \subset \Lambda^2T^\ast M\) satisfying the quaternion relations. Define the tensor
\begin{align*}
    R_{ijkl}^{\mathbb{HP}^m} := g_{il}g_{jk} -  g_{jl}g_{ik} + \sum_{\alpha = 1}^3 {}^\alpha I_{il}{}^\alpha I_{jk} - {}^\alpha I_{jl}{}^\alpha I_{ik} - 2 \ {}^\alpha I_{ij}{}^\alpha I_{kl}.
\end{align*} This represents the curvature tensor of \(\mathbb{HP}^m\) with scalar curvature \(16m(m+2)\). Thus by Salamon, we have
\begin{align*}
    R_{ijkl} = \frac{R}{16m(m+2)} R_{ijkl}^{\mathbb{HP}^m} + R_{ijkl}^{\text{hyp}},
\end{align*} where \(R_{ijkl}^{\text{hyp}}\) is a section of $\text{Sym}^2\left(\Lambda_-^2\right).$ %\(S^2(\mathfrak{sp}(m))\). 
Next we consider the action of \(R_{ijkl}^{\mathbb{HP}^m}\) on a form \(\omega \in \Omega_+^2.\) %\(\omega \in \Gamma(\mathfrak{sp}(1))\). 
%Letting \(\omega_{ij} = f \ {}^1 I_{ij} + g \ {}^2 I_{ij} + k \ {}^3 I_{ij}\), we compute
We compute
\begin{align*}
    R_{ij}^{\ \ kl} \omega_{kl} = \frac{R}{16m(m+2)}(-2\omega_{ij} + (-2 + 8m)\omega_{ij} + 4\omega_{ij}) = \frac{R}{2(m+2)} \omega_{ij}.
\end{align*} Therefore the Riemannian zeroth-order term in the Weitzenbock formula for \(\Omega_+^2(\mathfrak{g}_E)\) is
\begin{align*}
    \left(\frac{R}{4m} + \frac{R}{4m} -\frac{R}{2(m+2)}\right) \operatorname{id} = \frac{R}{m(m+2)}\operatorname{id}. %\frac{1}{m+2}\frac{R}{m} \operatorname{id} > 0.
\end{align*} In particular, suppose \(A\) is a pseudo-holomorphic connection. Since \([\mathfrak{sp}(m), \mathfrak{sp}(1)] = 0\) in \(\Lambda^2T^\ast M\), we have for \(\omega \in \Omega_+^2(\mathfrak{g}_E)\) that
\begin{align}\label{preliminaries:qkweitzenbock}
    (\Delta_A \omega)_{ij} = (\nabla_A^\ast \nabla_A \omega)_{ij} - [(F_A^+)_i^{\ k}, \omega_{kj}] + [(F_A^+)_j^{\ k}, \omega_{ki}] + \frac{R}{m(m+2)}\omega.
\end{align}
\begin{comment}
\begin{align}\label{preliminaries:qkweitzenbock}
    \Delta_A \omega = \nabla_A^\ast \nabla_A \omega - \llbracket F_A^+, \omega\rrbracket + \frac{R}{m(m+2)}\omega.
\end{align} 
\end{comment}


\vspace{5mm}



\subsection{Yang-Mills flow}
Given an initial connection $A_0,$ we wish to produce a solution $A(t)$ of (\ref{ymf}) on a maximal time interval. The gauge-invariance of the Yang-Mills functional \(\mathcal{YM}\) implies that the flow equation (\ref{ymf}), which is the downward gradient flow of \(\mathcal{YM}\), is not strictly parabolic. Nonetheless, we may follow the strategy of R{\aa}de \cite[Theorem 1]{rade}, based on DeTurck's first simplification\footnote{It is also possible to obtain well-posedness of the flow via the better-known Donaldson-DeTurck ``trick,'' see e.g. \cite{kozonomaedanaito4dglobal, ohtatarucaloric, struweym}, but obtaining the optimal result is no more straightforward than in R\aa de's setup.} \cite{deturckfirsttrick} of Hamilton's existence result for Ricci flow, to obtain a unique, short-time solution of (\ref{ymf}), as well as continuity of the solution with respect to the initial data in \(H^k(\mathcal{A}_E)\) for \(k > \lceil\frac{n}{2}-1\rceil\). The main well-posedness results are detailed in Theorem \ref{thm:wellposedness}. \par 
The key observation underpinning R{\aa}de's strategy is that if \(A(t)\) evolves by (\ref{ymf}), then the curvature evolves by the parabolic equation
\begin{align}\label{preliminaries:evolutionofcurvature}
    \frac{\partial F_A}{\partial t} + \Delta_A F_A = 0.
\end{align} Thus we obtain a solution of (\ref{ymf}) by solving for \((a(t), \Omega(t)) \in \Omega^1(\mathfrak{g}_E) \oplus \Omega^2(\mathfrak{g}_E)\) satisfying
\begin{align*}
    \frac{\partial a}{\partial t} &= -D_{\nabla_{\text{ref}} + a}^\ast \Omega \\
    \frac{\partial \Omega}{\partial t} &= -\Delta_{\nabla_{\text{ref}} + a} \Omega,
\end{align*} and then showing that if \(\Omega(0) = F_{\nabla_{\text{ref}} + a(0)}\), then \(\Omega(t) = F_{\nabla_{\text{ref}} + a(t)}\). This is explained in the short-time existence/uniqueness result, Theorem \ref{thm:shorttimeexistence}. \par
To prove the (neighborhood) deformation retraction statements of Theorems \ref{thm:fourspheregap}, \ref{thm:quaternionkahlergap}, and \ref{thm:flatgap}, we must \textit{a priori} work with Sobolev solutions of (\ref{ymf}), i.e., where \(A(t)\) is in \(H^k\). However, as we show in the appendix, any \(H^k\) solution is gauge-equivalent by a constant-in-time \(H^{k+1}\) gauge transformation to a smooth solution. Thus, in most instances we need only consider smooth solutions of (\ref{ymf}). \par 
Last, we note that the (neighborhood) deformation retractions of the previous three theorems are continuous on spaces of connections and not just on gauge-equivalence classes. To obtain convergence of solutions of (\ref{ymf}) and continuity of the infinite-time limit in \(H^k\) without modding out by gauge, we prove in the appendix a Sobolev distance estimate (Theorem \ref{thm:Hkestimates}) for solutions of (\ref{ymf}) with bounded curvature, as well as an \(H^{-1}\) {\L}ojasiewicz-Simon inequality in Lemma \ref{lemma:lojasiewicz}, cf. \cite[(9.1)]{rade}.


\vspace{10mm}







\section{Deformation to instantons on the 4-sphere}





\subsection{Optimal split $\eps$-regularity} We first give an optimal version of the split $\eps$-regularity estimate relied upon in \cite{instantons}.



\begin{defn}\label{defn:delta} Let $\calB_{S^4,G}$ denote the space of all gauge-equivalence classes of smooth connections carried by principal $G$-bundles $P \to S^4,$ including all topologies. %where we take the disjoint union over all topological bundles $P \to S^4.$
Given a representation $\rho: G \to \GL_K(r),$ where $K = \R \text{ or } \C,$ define
\begin{equation*}
\delta_{G, \rho} = \inf \left\{ \frac{1}{4 \pi^2 n_{K}} \left\| F^+_{B_\rho} \right\|^2 \mid \LB B \RB \in \calB_{S^4,G}, D^*F_B = 0, F^+_B \not\equiv 0 \right\}.
\end{equation*}
Here, $B_\rho$ is the connection induced on the vector bundle $E = P \times_\rho V$ by the representation $\rho.$
\end{defn}
\begin{rmk}\label{rmk:deltaGrho} We make several comments on the foregoing definition.
\begin{enumerate}[itemsep=2mm]

\item Evidently, $\delta_{G, \rho}$ is the largest number such that the implication
\begin{equation*}%\label{delta:implication}
\begin{split}
D_A^* F_A = 0, \quad \|F^+_A\|^2 < 4 \pi^2 n_K \delta_{G, \rho} %\YM(A) < 4 \pi^2 n_{K} \left( \kappa(E) + \delta_{G, \rho} \right) 
\quad \Rightarrow \quad F^+_A = 0
\end{split}
\end{equation*}
holds generally for connections $A$ on $(G, \rho)$-bundles $E \to S^4.$ In particular, if $\kappa(E) \geq 0,$ any Yang-Mills connection with
$ \YM(A) < 4 \pi^2 n_{K} \left( \kappa(E) + \delta_{G, \rho} \right) $ is anti-self-dual.


\item We have $\delta_{G, \rho} \geq \delta_{\SO(r), \rho_{std}}$ for $K = \R$ and $\delta_{G, \rho} \geq \delta_{\SU(r), \rho_{std}}$ for $K = \C.$ %where $r$ is the dimension of $\rho.$
By the Min-Oo gap theorem, %$\delta_{G, \rho} \geq \delta_0 > 0$
these are bounded below by a positive universal constant.

\item $\delta_{G, \rho}$ determines $\delta_{G, \rho'}$ for all other representations $\rho'$ by an algebraic constant; the difference comes only from the norm on $\rho(\gothg) \subset V \otimes V^*.$

\item
For $G = \SU(r)$ and the standard representation, according to \cite[Corollary 1.2]{gurskykelleherstreets} (see Remark \ref{rmk:gkscomparison}), we have
$$\delta_{\SU(r), \rho_{std}} = 2.$$
This corresponds to the general case $\gamma_1 = \frac{4}{\sqrt{3}}$ in the notation of \cite{gurskykelleherstreets}.

\item For $G = \SO(r)$ and the standard representation, we also have
$\delta_{\SO(r), \rho_{std}} = 1$ for $r \geq 4$ and $\delta_{\SO(3), \rho_{std}} = 2.$

\end{enumerate}
\end{rmk} 




\begin{thm}\label{thm:strongepsreg}
%   Let \(M\) be a Riemannian 4-manifold without boundary,
    %with bounded geometry,
%    \(P \to M\) a principal \(\mathrm{SU}(2)\)-bundle, $x_0 \in M,$ and $0 < R < R_0,$ depending only on the geometry of $M$ near $x_0.$
Given \(\gamma > 0\), there exist \(k \in \N\) and a constant $C_{(\ref{thm:strongepsreg})} > 0$, which depends on \(\gamma \) and the structure group \((G, \rho)\), as follows. 

Let $B = B_1(0)$ be a geodesic ball for a smooth metric $g$ which is sufficiently close in $C^k$ to the Euclidean metric in normal coordinates. %Let $g$ be a smooth metric on $B = B_1(0) \subset \R^4$ sufficiently close in $C^k$ to the Euclidean metric. %Let \(G\) be a compact Lie group. Denote
%\begin{align*}
%    \eta = \eta(G) := \inf\{\|F_A^+\|^2 \mid A \text{ is a non-ASD connection on a principal } G-\text{bundle over } S_\text{std}^4\}.
%\end{align*}
Suppose that \(A(t)\) is a smooth solution of (\ref{ymf}) on a $(G, \rho)$-bundle $E \to B$ for $t \in \left( -1, 0 \right)$, with 
    \begin{align}\label{strongepsreg:energybound}
        \sup_{-1 < t < 0} \YM(A(t)) \leq \gamma^{-1}
    \end{align}
    and
    \begin{align}\label{strongepsreg:GKSbound}
        \sup_{-1 < t < 0} \|F_{A(t)}^+\|_{L^2(B)}^2 \leq \delta \leq 4 \pi^2 n_{K} \delta_{G, \rho} - \gamma,
    \end{align}
    where self-duality and norms are defined by the metric $g.$
    Then 
    \begin{equation}\label{epsilonregularitybound}
        \sup_{B_{1/2} \times \LB -1/4, 0 \right)} | F_{A(t)}^+(x,t) | \leq C_{(\ref{thm:strongepsreg}) } \sqrt{\delta}.
    \end{equation}
 %   Here $C_{\ref{lemma:strongepsreg}}$ depends on $E $ and $\gamma.$
\end{thm}

The proof relies on the following two lemmas. In the sequel, we denote $F^+(t) := F_{A(t)}^+.$

\begin{lemma}\label{lemma:energynotsmall}
    There exists a constant $0 < c < \frac{1}{10}$, depending only on the structure group \((G, \rho)\), as follows. Given a solution $A(t)$ satisfying the assumptions of Theorem \ref{thm:strongepsreg}, if $0 < r \leq 1$ is such that
    \begin{equation}\label{energynotsmall:assn}
       \int_{B_{r}} \left|F^+\left( -(cr)^2 \right) \right|^2 \leq c^2,
    \end{equation}
    then
    \begin{equation}\label{energynotsmall:est}
        \sup_{ B_{r/2} \times \LB - \frac12 (cr)^2, 0 \right) } |F^+| \leq \frac{1}{cr^2}.
    \end{equation}
\end{lemma}
\begin{proof} 
By parabolic rescaling, we may assume without loss of generality that $r = 1.$ %parabolically rescaling by the factor $r,$ we 
The assumption (\ref{energynotsmall:assn}) reads
\begin{equation}
    \int_{B} |F^+(-c^2)|^2 \leq c^2.
\end{equation} In the sequel, \(c\) may decrease in each appearance but will depend only on the structure group \((G, \rho)\). We also let $C > 0$ denote a constant depending only on $(G, \rho)$ which may increase in each appearance.
Recall that the curvature evolves by
\begin{align*}
    \partial_t F + D D^\ast F = 0.
\end{align*} Letting \(\varphi \in C_c^\infty(B)\), we take the inner product with \(\varphi^2 F^+\), integrate in space, and then integrate by parts. This yields
\begin{align*}
    \frac{1}{2}\frac{d}{dt} \int_B \varphi^2 |F^+|^2 &= \int_B \langle \varphi^2 F^+, \partial_t (F^+)\rangle = \int_B \langle \varphi^2 F^+, (\partial_t F)^+\rangle = \int_B \langle \varphi^2 F^+, \partial_t F\rangle \\
    &= -\int_B \langle \varphi^2 F^+, D D^\ast F\rangle \\
    &= -\int_B \varphi^2\langle D^\ast F^+, D^\ast F\rangle + 2\int_B \langle \nabla \varphi \intprod F^+,  \varphi D^\ast F\rangle.
\end{align*} Recall that \(D^\ast F = 2D^\ast F^+\). Hence we obtain
\begin{align*}
    \frac{d}{dt}\int_B \varphi^2 |F^+|^2 + \int_B \varphi^2 |D^\ast F|^2 &= 4\int_B  \langle  \varphi D^\ast F, \nabla \varphi \intprod F^+\rangle.
\end{align*}
The Cauchy-Schwarz and Young's inequalities imply
\begin{align*}
    \left|4\int_B \varphi \langle D^\ast F^+, (\nabla \varphi) \intprod F^+\rangle\right| \leq \frac{1}{2}\int_B \varphi^2 |D^\ast F|^2 + 8\int_B |\nabla \varphi|^2 |F^+|^2.
\end{align*} Integrating in time and rearranging yields 
\begin{align}\label{localenergyinequality}
    \int_B \varphi^2 |F^+(b)|^2 + \frac{1}{2}\int_a^b\int_B \varphi^2 |D^\ast F|^2 \leq \int_B \varphi^2 |F^+(a)|^2 + 8\int_a^b\int_B |\nabla \varphi|^2|F^+|^2.
    \end{align}
    Let \(\varphi\) be a cutoff for \(B_{3/4}.\) Since the metric is assumed to be close to Euclidean, we have $\|\nabla \varphi\|_\infty \leq C.$ Letting \(a = -c^2\) and \(b = 0\) in (\ref{localenergyinequality}), we obtain
    \begin{comment}
    In particular, letting \(\varphi\) be a cutoff for \(B_{3/4}\), \(a = -c^2\), and \(b = 0\), we obtain
    \end{comment}
    \begin{align}\label{localenergybound}
    \sup_{-c^2 \leq t \leq 0} \int_{B_{3/4}} |F^+(t)|^2 \leq c^2 + C \delta_{G, \rho} c^2 \leq C c^2.
    \end{align}
    For $c$ sufficiently small, it follows from standard $\eps$-regularity (Prop. 2.3 of \cite{instantons}) that %for some universal constant \(C\)
    $$\sup_{B_{1/2} \times \LB -\frac12 c^2, 0 \right) }|F^+| \leq \frac{C c^2}{c^2} = C \leq \frac{1}{c}.$$
    Undoing the parabolic rescaling yields (\ref{energynotsmall:est}).
\end{proof}


We will use the following well-known point selection lemma. Denote %In the sequel, we will use the following notation for parabolic cylinders
\begin{align*}
    P_r(x, t) = B_r(x) \times [t - r^2, t],
\end{align*} where \(B_r(x)\) is a geodesic ball of radius \(r\) centered at \(x\). For the following lemma, we do not make any assumptions about the metric on \(B_r(x)\).




\begin{lemma}\label{lemma:comparablesup}
    Let \(\lambda > 0\) and let \(f\) be a non-negative continuous function on \(P_r(x_0, t_0)\) such that \(f(x_0, t_0) = \lambda^{-2} > 16r^{-2}\). Then there exists a spacetime point \((\tilde{x}, \tilde{t})\) such that, denoting \(\mu^{-2} := f(\tilde{x}, \tilde{t})\), we have \(\mu \leq \lambda\), \((\tilde{x}, \tilde{t}) \in P_{2\lambda}(x_0, t_0)\), and     \begin{align} \label{comparablesupbound}
\sup_{P_\mu(\tilde{x}, \tilde{t})} f \leq 4\mu^{-2}.
    \end{align}
\end{lemma}
\begin{proof}
    Suppose for %a 
    contradiction that such a spacetime point does not exist. Then in particular the desired conclusion does not hold for \((x_0, t_0)\), so we have
    \begin{align*}
        \sup_{P_\lambda(x_0, t_0)} f > 4\lambda^{-2}.
    \end{align*} Assume by induction that we have found \(n+1\) spacetime points \((x_i, t_i)\), \(i = 0, \ldots, n\) satisfying
    the following conditions:
    \begin{align*}
       (\ast) \begin{cases}
            (x_{i+1}, t_{i+1}) \in \overline{P_{\mu_i}(x_i, t_i)} \ \ \ \text{ where } \mu_i^{-2} := f(x_i, t_i), \\
            f(x_{i + 1}, t_{i+1}) = \sup_{P_{\mu_i}(x_i, t_i)} f > 4\mu_i^{-2}.
        \end{cases}
    \end{align*} Note that \(\mu_{i+1} < 2^{-1}\mu_i\), and hence \(\mu_n < 2^{-n}\lambda\). We let \(d\) denote the parabolic distance \(d((p, t), (q, s)) = \max\{d_g(p, q), \sqrt{|t-s|}\}\), for which parabolic cylinders are metric balls. Observe that 
    \begin{align*}
        d((x_0, t_0), (x_n, t_n)) \leq \sum_{i = 0}^n d((x_{i+1}, t_{i+1}), (x_i, t_i)) \leq \sum_{i = 0}^n \mu_i < \lambda \sum_{i = 0}^n 2^{-i} < 2\lambda.
    \end{align*} Thus for \(0 \leq i \leq n-1\), \((x_i, t_i)\) is in \(P_{2\lambda}(x_0, t_0)\). Additionally, \(P_{\mu_i}(x_i, t_i)\) is compactly contained in \(P_r(x_0, t_0)\) since \(\lambda < r/4\). Our assumption that there is no spacetime point satisfying the conclusion of the lemma implies the existence of a spacetime point \((x_{n+1}, t_{n+1})\) satisfying \((\ast)\). Therefore, by induction we obtain an infinite sequence of spacetime points \(\{(x_i, t_i)\}\) satisfying \((\ast)\). Since this sequence is contained in the compact set \(\overline{P_{2\lambda}(x_0, t_0)}\) and since \(f\) is continuous, and yet \(\mu_i^{-2}\) tends to infinity, we reach a contradiction, as desired.
\end{proof}


\begin{proof}[Proof of Theorem \ref{thm:strongepsreg}]
We first derive the desired estimate without the $\delta$ on the right-hand side.

Suppose for contradiction that there exists a sequence of solutions $A_i$ on the unit parabolic cylinder (with metrics tending to Euclidean in $C^\infty$) such that
$$\sup_{B_{1/2} \times \LB -\frac14, 0 \RB} |F^+_{A_i}(x,t) | \geq c^{-4i-1},$$%.$$
where $c$ is the same constant from Lemma \ref{lemma:energynotsmall}. We abbreviate $F_i := F_{A_i}.$ After recentering and rescaling by a factor no less than $\frac12,$ we may assume
$$|F^+_i(0,0)| \geq c^{-4i-1}.$$
Let
$$t_j = -c^{4j}$$
%$$t_i = -c^{4i}$$
and
$$I_j = \LB t_{j-1}, t_{j} \RB.$$
%$$I_i = \LB t_{i-1}, t_{i} \RB.$$
In view of (\ref{strongepsreg:GKSbound}) and the convergence of the metrics, (\ref{localenergyinequality}) implies that we can assume
$$\int_{P_1(0, 0)} |D_i^*F_i|^2 \, dV dt \leq C \delta.$$
Since the time-intervals $I_j$ are essentially disjoint, for each $i$ there exists $j_i \in \{\lceil \frac{i}{2} \rceil, \ldots, i\},$ such that
\begin{align}\label{tensiongoingtozero}
    \int_{B_1(0) \times I_{j_i}} |D_i^*F_{i}|^2 \, dV dt \leq \frac{2C \delta}{i}.
\end{align} Denote
\begin{align*}
    \tau_i := t_{j_i}.
\end{align*} By the contrapositive of Lemma \ref{lemma:energynotsmall}, we must %also
have $$\int_{B_{c^{-1}\sqrt{-\tau_i}}(0)} |F^+_i(\tau_i)|^2 > c^2.$$
In particular, there exist %exists
points \(\tilde{x}_{i}\) in \(B_{c^{-1}\sqrt{-\tau_i}}(0)\) such that
\begin{equation}
\lambda_i^{-2} := |F^+_i(\tilde{x}_{i}, \tau_i)| > \frac{c^{3}}{C(-\tau_i)}.
%\frac{c^3}{-\tau_i}.
\end{equation}
\begin{comment}
$$\int_{B_{c^{-1}\sqrt{-t_{j_i}}}(0)} |F^+_i(t_{j_i})|^2 > c^2.$$
In particular, there exist %exists
points \(\tilde{x}_{j_i}\) in \(B_{c^{-1}\sqrt{-t_{j_i}}}(0)\) such that
\begin{equation}
\lambda_j^{-2} := |F^+_i(\tilde{x}_{i,j}, t_j)| > \frac{c^3}{-t_j}.
\end{equation}
\end{comment}
\begin{comment}
Also, in view of (\ref{strongepsreg:GKSbound}) and the convergence of the metrics, (\ref{localenergyinequality}) implies that we can assume
%Also, by (\ref{localenergyinequality}), we can assume that
$$\int_{P_1(0, 0)} |D_i^*F_i|^2 \, dV dt \leq C \delta.$$
Since the time-intervals $I_j$ are essentially disjoint, for each $i$ there exists $j_i \in \{\lceil \frac{i}{2} \rceil, \ldots, i\},$ such that
\begin{align}\label{tensiongoingtozero}
    \int_{B_1(0) \times I_{j_i}} |D_i^*F_{i}|^2 \, dV dt \leq \frac{2C \delta}{i}.
\end{align}
\end{comment}
Denote $r_i^2 := (c^{-4}-1)(-\tau_i).$
%\(r_i^2 := t_{j_i}-t_{j_i-1}\).
Since for \(c\) sufficiently small (but still depending only on the structure group \((G, \rho)\))
\begin{align*}
    r_i^2\lambda_{i}^{-2} > r_i^2\frac{c^{3}}{C(-\tau_i)} = -\tau_i\left(\frac{1}{c^4}-1\right)\frac{c^{3}}{C(-\tau_i)} > 16,
\end{align*}
\begin{comment}
\begin{align*}
    r_i^2\lambda_{j_i}^{-2} > r_i^2\frac{c^3}{-t_{j_i}} = -t_{j_i}\left(\frac{1}{c^4}-1\right)\frac{c^3}{-t_{j_i}} > 16,
\end{align*}
\end{comment}
it follows that \(|F_i^+|\) satisfies the hypotheses of Lemma \ref{lemma:comparablesup} on $P_{r_i}(\tilde{x}_i, \tau_i).$ %\(P_{r_i}(\tilde{x}_{j_i}, t_{j_i})\).
Thus there exists a spacetime point $(x_i, s_i) \in P_{2\lambda_i}(\tilde{x}_i, \tau_i)$
%\((x_i, s_i) \in P_{2\lambda_{j_i}}(\tilde{x}_{j_i}, t_{j_i})\)
such that, denoting \(\mu_i^{-2} := |F_i^+(x_i, s_i)|\),
\begin{align}\label{F+comparablesup}
    &\sup_{P_{\mu_i}(x_i, s_i)} |F_i^+| \leq 4\mu_i^{-2}.
\end{align} Moreover, we have \(\mu_i \leq \lambda_{i} \leq \frac{r_i}{4}\), so
\begin{align*}
    P_{\mu_i}(x_i, s_i) \subset P_{\mu_i + 2\lambda_{i}}(\tilde{x}_i, \tau_i) \subset P_{r_i}(\tilde{x}_i, \tau_i). 
    %P_{\mu_i + 2\lambda_{i}}(\tilde{x}_{j_i}, t_{j_i}) \subset P_{r_i}(\tilde{x}_{j_i}, t_{j_i}).
\end{align*} In particular, we must have
\begin{align*}
    [s_i - \mu_i^2, s_i] \subset [\tau_i - r_i^2, \tau_i] = I_{j_i}.
    %[t_{j_i} - r_i^2, t_{j_i}].
\end{align*}
%Additionally, by construction \(j_i \geq \lceil \frac{i}{2}\rceil\), so
By construction, $j_i \geq \lceil \frac{i}{2}\rceil.$ Thus
\begin{align*}
    \sqrt{-\tau_i} \leq \sqrt{c^{4\frac{i}{2}}} = c^i
\end{align*} and consequently
\begin{align*}
    \lambda_i \leq \frac{r_i}{4} = \frac{1}{4}\sqrt{(c^{-4}-1)(-\tau_i) } \leq c^{i-2}.
\end{align*} Therefore
\begin{align*}
    d_g(x_i, 0) \leq d_g(x_i, \tilde{x}_i) + d_g(\tilde{x}_i, 0) \leq 2\lambda_i + c^{-1}\sqrt{-\tau_i} 
    %\leq c^{-3}\sqrt{-\tau_i} 
    \leq c^{i-3} \to 0
\end{align*}
as $i \to \infty.$
\begin{comment}
\begin{align*}
    d_g(x_i, 0) \leq d_g(x_i, \tilde{x}_{j_i}) + d_g(\tilde{x}_{j_i}, 0) \leq 2\lambda_{j_i} + c^{-1}\sqrt{-t_{j_i}} \leq c^{-2}\sqrt{-t_{j_i}} \leq c^{i-3} \to 0.
\end{align*}
\end{comment}
Hence, for \(i\) large enough, we deduce that
\begin{align}\label{containedwhereD*Fsmall}
    P_{\mu_i}(x_i, s_i) \subset B_{1/2}(x_i) \times [s_i-\mu_i^2, s_i] \subset B_1(0) \times I_{j_i}.
\end{align}\par
 We now parabolically rescale \(B_{1/2}(x_i) \times [s_i-\mu_i^2, s_i]\) about \((x_i, s_i)\) by the factor \(\mu_i\). This yields a sequence of solutions of (\ref{ymf}), denoted by \(\tilde{A}_i\), on \(B_{(2\mu_i)^{-1}}(0) \times [-1, 0]\), where the metric on \(B_{(2\mu_i)^{-1}}(0)\) is approaching the Euclidean metric in \(C_\text{loc}^\infty\). In view of the bounds (\ref{strongepsreg:energybound}) and (\ref{tensiongoingtozero}) and the containment (\ref{containedwhereD*Fsmall}), we may apply standard Uhlenbeck compactness theory (see e.g. \cite[Theorem 1.3 and \textsection 5]{waldronuhlenbeck}) as follows. After passing to a subsequence of the \(\tilde{A}_i\), there exists a finite collection of points \(y_k\) in \(\R^4\), a finite-energy smooth connection \(A_\infty\) on a \((G, \rho)\)-bundle over \(\R^4 \setminus \{y_k\}\),  and bundle maps \(u_i\) defined on an exhaustion of \(\R^4 \setminus \{y_k\}\), such that \(u_i(\tilde{A}_i)\) converges to a Yang-Mills connection \(A_\infty\) in \(C_{\text{loc}}^\infty(\R^4\setminus \{y_k\} \times (-1, 0))\). %Moreover, the bound (\ref{tensiongoingtozero}) and the containment (\ref{containedwhereD*Fsmall}) imply that \(A_\infty\) is Yang-Mills.
 
Writing \(\tilde{F}_i^+ = F_{u_i(\tilde{A}_i)}^+\),
    note that \begin{align}\label{capturedenergy}
    |\tilde{F}_i^+(0, 0)| = 1,
\end{align}
so in particular by Lemma \ref{lemma:energynotsmall}, we have
\begin{equation*}
\int_{B_{1/2}} \left| \tilde{F}_i^+ \left( -\frac12 \right) \right|^2 \geq c^2.
\end{equation*}
Meanwhile, (\ref{F+comparablesup}) implies
\begin{align}\label{blowupbounded}
    \sup_{P_1(0, 0)} |\tilde{F}_i^+| \leq 4.
\end{align}
Therefore our convergence is strong in $L^2$ on $B_1(0),$ and we still have
\begin{equation}\label{strongepsreg:F+Ainftynonzero}
\int_{B_{1/2}} |F_{A_\infty}^+ |^2 \geq c^2
\end{equation}
in the limit. \par
 On the other hand, by Uhlenbeck's removable singularity theorem, we may extend the bundle in question over the point at infinity to obtain a smooth Yang-Mills connection \((\tilde{E}, A_\infty)\) on the 4-sphere. It follows from (\ref{strongepsreg:GKSbound}), conformal invariance of the energy under the rescaling, Fatou's lemma, and (\ref{strongepsreg:F+Ainftynonzero}) that
 $$c^2 \leq \|F_{A_\infty}^+\|^2 \leq 4\pi^2 n_{K} \delta_{G, \rho} - \gamma.$$
 But then the implication of Remark \ref{rmk:deltaGrho}(1) gives $F^+_{A_\infty} \equiv 0,$ which is a contradiction. %this violates the Gursky-Kelleher-Streets gap theorem \cite[Corollary 1.2 (1)]{gurskykelleherstreets}.
 
 We have established the desired bound (\ref{epsilonregularitybound}) without \(\delta\) on the right-hand side.
    To obtain the bound with \(\delta\), we proceed as follows. Here \(C_{(\ref{thm:strongepsreg})}\) may increase from line to line but will keep the same dependencies. The bound without \(\delta\) implies
    \begin{align*}
        \sup_{P_{3/4}(0, 0)} |F^+| \leq C_{(\ref{thm:strongepsreg})}.
    \end{align*} Hence the self-dual curvature satisfies the differential inequality
    \begin{align*}
        (\partial_t - \Delta) |F^+| \leq C_{(\ref{thm:strongepsreg})}|F^+|.
    \end{align*} Parabolic Moser iteration (see \cite[Lemma 19.1]{ligeomanalysis}) then yields the bound
    \begin{align*}
        \sup_{P_{1/2}(0, 0)} |F^+| \leq C_{(\ref{thm:strongepsreg})}\sqrt{\delta},
    \end{align*} as desired.
\end{proof}




\vspace{5mm}





\subsection{Proof of Theorem \ref{thm:fourspheregap}} %We now prove the convergence theorem stated in the introduction.
We will use the following result, which is a version of the central estimate of \cite[Theorem 2.5]{instantons}.


\begin{thm}[Cor. 7.3 of \cite{oliveirawaldron}]\label{thm:F+boundimpliesFbound}
    Let \(M\) be a closed Riemannian \(4\)-manifold, and suppose \(A(t)\) is a solution of (\ref{ymf}) on \([0, T)\). Let \(0 \leq t_1 \leq t_2 < T\), and set
    \begin{align*}
       E_0 &:= \|F(0)\|^2 \\
       K &:= \int_{t_1}^{t_2} \|F^+(t)\|_\infty \, dt.
    \end{align*} There exist positive constants \(C\), which is universal, and  \(R_0\), depending on \(E_0\) and the geometry of \(M\), such that
    \begin{align}\label{F+boundimpliesFbound:bound}
        \|F(t_2)\|_\infty \leq \max\left\{R_0^{-2}, e^{C\max\left\{1, K \right\}}\|F(t_1)\|_\infty \right\}.
    \end{align}
\end{thm}

\begin{thm}[Cf. Theorem \ref{thm:fourspheregap}]\label{thm:fourspheregeneralgap}
Let \(k \geq 2\), and let $E \to S^4$ be a $(G, \rho)$-bundle. Suppose \(A_0\) is an \(H^k\) connection on \(E\) satisfying
\begin{equation}\label{fourspheregeneralgap:energylessthandelta}
\YM(A_0) < 4 \pi^2 n_{K}\left( |\kappa(E)|+ \delta_{G, \rho} \right).
\end{equation}
Here $\kappa(E)$ is defined in \S \ref{ss:characteristicnumber} and $\delta_{G, \rho}$ is given in Definition \ref{defn:delta}.
Then the solution \(A(t)\) of (\ref{ymf}) with \(A(0) = A_0\) exists for all time and converges in \(H^k\) exponentially as \(t \to \infty\) to an instanton on $E$. Moreover, the map $A_0 \mapsto \lim_{t \to \infty} A(t)$ for \(A_0\) satisfying (\ref{fourspheregeneralgap:energylessthandelta}) is homotopic to the identity relative to the %moduli 
space of instantons.
\end{thm}
\begin{proof}
Define the set
\begin{align*}
    \calN := \{A_0 \in H^k(\mathcal{A}_E) \mid A_0 \text{ satisfies } (\ref{fourspheregeneralgap:energylessthandelta})\},
\end{align*}which is open in the \(H^k\) topology. Given \(A_0\ \in \calN\), item (\ref{wellposedness:gaugeequivalenttosmooth}) of Theorem \ref{thm:wellposedness} implies that there exists an \(H^{k+1}\) gauge transformation \(u\) such that \(u(A(t))\) is a smooth solution of (\ref{ymf}) for \(t > 0\). Thus if \(u(A(t))\) exists for all time and converges smoothly to an instanton, \(A(t)\) will exist for all time and converge in \(H^k\) to an instanton. Thus to prove the first part of the theorem, we assume that \(A(t)\) is smooth. \par 
Because $\mathcal{YM}(A(t))$ is nonincreasing, the assumption (\ref{fourspheregeneralgap:energylessthandelta}) is preserved. %Since \(S^4\) is compact, there exists some radius \(R_0 > 0\) 
There exists a universal \(R_0 > 0\) such that for any point \(x \in S^4\), rescaling \(B_{R_0}(x)\) by the factor \(R_0\) yields a ball satisfying the metric hypothesis of Theorem \ref{thm:strongepsreg}, whence for \(t > 0\)
\begin{align}\label{fourspheregeneralgap:selfdualsupisbounded}
    |F^+(x, t)| \leq \frac{4\pi C_{(\ref{thm:strongepsreg})}\sqrt{\delta_{G, \rho}}}{\min\{1, t\}}.
\end{align} It then follows from (\ref{F+boundimpliesFbound:bound}) and item (\ref{wellposedness:blowupcharacterization}) of Theorem \ref{thm:wellposedness} that \(A(t)\) exists for all time. Moreover, the full curvature blows up at most exponentially at infinite time. We will now argue that the solution in fact converges smoothly. \par 

Applying Uhlenbeck's compactness and removable-singularity theorems as in the proof of Theorem \ref{thm:strongepsreg}, we may extract an Uhlenbeck limit \(A_\infty'\) from the sequence of connections $A(n),$ $n \in \N,$ which is a Yang-Mills connection on a possibly different \((G, \rho)\)-bundle over $S^4.$ By Fatou's lemma and the global energy inequality, we have  
$$\left\|F^+_{A_\infty'}\right\|^2 \leq \|F^+(0)\|^2 < 4\pi^2n_{K}\delta_{G, \, \rho}.$$
It now follows from the implication of item (1) of Remark \ref{rmk:deltaGrho} that $A_\infty'$ must be anti-self-dual:
\begin{align}\label{fourspheregeneralgap:UhlenbecklimitisASD}
    F^+_{A_\infty'} \equiv 0.
\end{align}

However, the Weitzenb{\"o}ck formula implies that anti-self-dual connections on $S^4$ have vanishing self-dual cohomology; it therefore follows directly from \cite[Corollary 3.8]{instantons} that $A(t)$ converges exponentially in the smooth topology. For purposes of exposition, we provide the following simpler argument in the present case.


     In the sequel, \(C\) denotes a positive universal constant which may increase from line to line.
     %, and \(\varepsilon_0\) denotes the epsilon-regularity constant, e.g. the \(k = 0\) case in \cite[Prop. 3.1 (a)]{lte}. %
     The bound (\ref{fourspheregeneralgap:selfdualsupisbounded}) implies that the Uhlenbeck convergence \(\left| F^+(n_j) \right| \to \left| F_{A_\infty'}^+ \right| \) is strong in \(L^2\). By (\ref{fourspheregeneralgap:UhlenbecklimitisASD}), \(\|F^+(n_j)\| \to 0\), so \(\|F^+(t)\| \to 0\) since the self-dual energy is decreasing. Now, the round metric on \(S^4\) has vanishing Weyl tensor and positive scalar curvature. Hence the Weitzenb{\"o}ck formula for self-dual forms implies
        \begin{align*}
         \frac{1}{2}\frac{d}{dt} \|F^+\|^2 &= - (\nabla^\ast \nabla  F^+, F^+) + (\llbracket F^+, F^+\rrbracket, F^+) - 4 \|F^+\|^2 \\
        &=  -\| \nabla  F^+\|^2 + \left(\llbracket F^+, F^+\rrbracket, F^+ \right) - 4\|F^+\|^2.
    \end{align*} By Cauchy-Schwarz and \cite[Lemma 2.30]{bourguignonlawson},
    \begin{align*}
        \left( \llbracket F^+, F^+\rrbracket, F^+ \right) \leq C \|F^+\|_4^2\|F^+\|.
    \end{align*}By Kato's inequality and the Sobolev inequality on functions,
    \begin{align*}
        \|F^+\|_4^2 \leq C(\|\nabla F^+\|^2 + \|F^+\|^2),
    \end{align*} where \(C\) is independent of the reference connection (see e.g., \cite[(1.16)]{instantons}). Since \(\|F^+(t)\| \to 0\) as \(t \to \infty\), there is some  time \(\tau_0 > 1\) such that for \(t \geq \tau_0\) we have
    \begin{align}\label{fourspheregeneralgap:selfdualenergysmall}
        \|F^+(t)\| < \min\left\{\frac{1}{4C},  1\right\},
    \end{align}whence
    \begin{align*}
        \frac{d}{dt} \|F^+\|^2 \leq -6\|F^+\|^2.
    \end{align*} By Gronwall's inequality, we have for \(t \geq \tau_0\)
    \begin{align}\label{fourspheregeneralgap:selfdualexponentialdecay}
        \|F^+(t)\| \leq \|F^+(\tau_0)\|e^{-3(t - \tau_0)} \leq e^{-3(t-\tau_0)}.
    \end{align} Then  (\ref{fourspheregeneralgap:selfdualexponentialdecay}) and $\varepsilon$-regularity for \(F^+\) yield
    \begin{align}\label{fourspheregeneralgap:F+exponentialdecay}
        \|F^+(t)\|_\infty \leq Ce^{-3(t-\tau_0)}, \ \ \ t \geq \tau_0 + 1.
    \end{align} Thus \(\|F^+(t)\|_\infty\) is in \(L^1([\tau_0 + 1, \infty))\), so it follows from (\ref{F+boundimpliesFbound:bound}) that
    \begin{align}\label{fourspheregeneralgap:fullcurvaturesupbound}
        \sup_{t \geq \tau_0 + 2} \|F(t)\|_\infty < \infty.
    \end{align}Note from this bound that the Uhlenbeck convergence \(A(n_j) \to A_\infty'\) is in fact smooth convergence modulo gauge, and \(A_\infty'\) is a connection on the original bundle.
    \begin{comment}
    Observe that for any integers \(N_2 > N_1 \geq 4\), the global energy identity implies
    \begin{align*}
        \int_{\tau_0 + N_1}^{\tau_0+N_2} \|D^\ast F(s)\| \, ds &\leq \sum_{j=N_1}^{N_2-1} \int_{\tau_0 + j}^{\tau_0 + j + 1} \|D^\ast F(s)\| \, ds \\
        &\leq \sum_{j=N_1}^{N_2-1} \left(\int_{\tau_0 + j}^{\tau_0 + j + 1} \|D^\ast F(s)\|^2 \, ds\right)^{\frac{1}{2}} \\
        &= \sum_{j=N_1}^{N_2-1} \left(\|F(\tau_0 + j + 1)\|^2 - \|F(\tau_0 + j)\|^2\right)^{\frac{1}{2}} \\
        &\leq \sum_{j = N_1}^{N_2-1} e^{-3j} \\
        &\leq Ce^{-3N_1}.
    \end{align*}
    \end{comment} 
    As in \cite[Theorem 3.7]{instantons}, we have for \(t > t_1 \geq \tau_0\)
    \begin{align}\label{fourspheregeneralgap:rhosmall}
        \int_{t_1}^t \|D^\ast F(s)\| \, ds \leq Ce^{-3(t_1 - \tau_0)}.
    \end{align} It then follows from (\ref{fourspheregeneralgap:fullcurvaturesupbound}) and Theorem \ref{thm:Hkestimates} that as \(t \to \infty\), \(A(t)\) converges exponentially in \(H^k\) for all \(k\) to some smooth instanton \(A_\infty\). It also follows that \(A_\infty\) is smoothly gauge equivalent to \(A_\infty'\) \cite[\textsection 2.3.7]{donkron}. \par 
     Finally, we explain the deformation-retraction property of the flow. Given a connection $A$ and $t \geq 0,$ let $\varphi_A(t)$ denote the solution of (\ref{ymf}) starting at $A$ evaluated at time $t.$ By the first part of the theorem, we may define the following map
    \begin{align*}
        \Psi &: \calN \times [0, 1] \to \calN \\
        \Psi(A, s) &:= 
        \begin{cases}
            \varphi_A\left(\frac{s}{1-s}\right) & 0 \leq s < 1 \\
            \lim_{t \to \infty} \varphi_A(t) & s = 1.
        \end{cases}
    \end{align*} Continuity of \(\Psi\) at \((A, s)\) for \(s \in [0, 1)\) follows from item (\ref{wellposedness:wellposedness}) of Theorem \ref{thm:wellposedness}. Continuity at \((A, 1)\) follows from Theorem \ref{thm:continuityatinfinitetime}, since instantons are in particular local minimizers of the Yang-Mills energy. Therefore, \(\Psi\) is continuous, and since the instantons are fixed under (\ref{ymf}), \(\Psi\) is a deformation retraction. Finally, the Sobolev multiplication theorem implies that the space of \(H^{k+1}\) gauge transformations acts smoothly on the space of \(H^k\) connections. Since this action preserves \(\calN\) and commutes with the mapping \(A \mapsto \varphi_A(\cdot)\), \(\Psi\) descends to the quotient of \(\calN\) by the action of \(H^{k+1}\) gauge transformations, as desired.
\end{proof}

\begin{rmk}
   Recall that for \(\CP^2\) with the Fubini-Study metric, there are no ASD instantons on the \(SU(2)\) bundle with \(c_2 = 1\) \cite[4.1.3, Example (iii)]{donkron}. Hence, there is no parabolic gap. This is in contrast to the elliptic \(L^2\) gap of Huang \cite{huangkahlersurfaces} for positive scalar curvature K{\"a}hler surfaces.
\end{rmk}



\vspace{5mm}


\subsection{Path-connectedness}\label{ss:pathconnected} We now describe our shortcut in Taubes's proof of his path-connectedness theorem. %path-connectedness of the $\SU(r)$ instanton moduli spaces on $S^4.$
%The original proof has two separate ingredients, Theorems 1.3 and 1.4 in \cite{taubespathconnected}, and we can use Theorem \ref{thm:fourspheregap} to simplify both of the proofs.



Fix $r \geq 2$ and let $\tilde{\mathcal{B}}_k$ denote the moduli space of framed connections on the charge-$k$ $\SU(r)$-bundle on $S^4,$ and let $\calM_k \subset \tilde{\calB}_k$ denote the moduli space of framed instantons. Taubes's proof has two separate ingredients, Theorems 1.3 and 1.4 in \cite{taubespathconnected}. The first contains Taubes's Lusternik-Schnirelman theory for the Yang-Mills functional,\footnote{It is interesting to note that the Lusternik-Schnirelman part of Taubes's proof also relies on a flow (different from Yang-Mills flow) to give a deformation retraction from a smaller neighborhood onto the space of instantons. See \cite[Proposition 3.1]{taubespathconnected}.} while the second contains his key trick for proving path-connectedness. %In the $\SU(2)$ case, 
We can replace %\cite[Theorem 1.3]{taubespathconnected}
the Lusternik-Schnirelman arguments entirely with Theorem \ref{thm:fourspheregap} and also simplify one step in the proof of the second result. %\cite[Theorem 1.4]{taubespathconnected}. %The key construction in Taubes's proof amounts to the following theorem.




\begin{comment}
\begin{thm}[Taubes \cite{taubespathconnected}, Theorem 1.3]
Suppose that there exists a path $\phi \in C^0\left( \LB 0,1 \RB; \mathcal{B}_k \right)$ with $\phi(\{0,1\}) \subset \calM_k,$ for which
\begin{equation}\label{leq2}
\sup_{y \in \LB 0,1 \RB} \YM(\phi(y)) < 4 \pi^2 \left( k + 2 \right).
\end{equation}
Then $\phi$ is homotopic relative to $\{0,1\}$ to a map into $\calM_k.$
\end{thm}
\begin{proof} By Theorem \ref{thm:fourspheregap}, the subset of $\tilde{\calB}_k$ satisfying (\ref{leq2}) deformation retracts onto $\calM_k,$ so the statement is clear.
\end{proof}
\end{comment}


\begin{thm}[Taubes \cite{taubespathconnected}, Theorem 1.4]\label{thm:taubespathconnected}
Suppose that for every $1 \leq k' < k,$ $\calM_{k'}$ is path-connected. Then the subset
\begin{equation}\label{leq2}
\calN_k = \{ \LB A \RB \in \tilde{\calB}_k \mid \YM(A) < 4 \pi^2 \left( k + 2 \right) \}
\end{equation}
is also path-connected.
\end{thm}
\begin{proof} We adopt Taubes's notation and recall the proof in \cite[\S 6]{taubespathconnected}, while obviating the use of \cite[Prop. 6.4]{taubespathconnected}. 

%Assume for the sake of induction that $\calM_{k-1}$ (and so too $\calN_{k-1},$ by Theorem 1.1) is path-connected.

Let $m_0, m_1 \in \calM_k.$ By \cite[Lemma 6.3]{taubespathconnected}, we may assume without loss of generality that both $m_0$ and $m_1$ have nonvanishing anti-self-dual curvature %self-dual curvature
at the south pole $s \in S^4.$ In this case,
according to \cite[Prop. 6.2]{taubespathconnected}, it is possible to choose $b_0, b_1 \in \mathcal{M}_1$ such that $\YM((m_i - b_i)_\rho) < 4 \pi^2 \left( k + 1 \right)$
%$\YM(m_i - b_i)_\rho < 4 \pi^2 \left( k + 1 \right)_\rho$ 
for $i = 0,1$; in particular, we have
$$(m_i - b_i)_\rho \in \calN_{k-1}.$$
Here, $(m_i - b_i)_\rho,$ $i = 0,1,$ are the connections of charge $k -1$ obtained by Taubes's ``subtraction'' procedure, and $\rho > 0$ is a small scale parameter. By hypothesis, $\calM_{k-1}$ is path-connected, so by Theorem \ref{thm:fourspheregap}, $\calN_{k-1}$  is as well. Hence, there exists a path $\gamma(t) \in \calN_{k-1},$ $t \in \LB 0,1 \RB,$ with $\gamma(i) = (m_i - b_i)_\rho$ for $i = 0,1.$

Next, since $\calM_1$ is path-connected (by hypothesis), we can choose $\phi(t) \in \calM_1$ with $\phi(i) = b_i$ for $i = 0,1.$ According to \cite[Lemma 6.5]{taubespathconnected}, %(which relies only on continuity of the energy in the gluing parameter),
for $r$ sufficiently small, the glued path
$$\psi(t):= \left( \gamma(t) - \alpha \phi(t) \right)_r$$
has energy less than $(k + 1) + 1 = k + 2,$ so lies entirely within $\calN_k.$ Here, $\alpha$ is the inversion map across the equator of $S^4$. Its endpoints are
$$\psi(i) = \left( (m_i - b_i)_\rho - \alpha t_\rho b_i \right)_r,$$
for $i = 0,1,$ where $t_\rho$ denotes pullback under dilation by the factor $\rho.$

Finally, by \cite[Lemma 6.6]{taubespathconnected}, it is possible to connect each endpoint $\psi(i)$ to $m_i$ within $\calN_k$ by ``cancelling'' $b_i$ against $t_\rho \alpha b_i,$ while again staying entirely within $\calN_k.$ The result is a path connecting $m_0$ and $m_1$ within $\calN_k,$ as desired.
\end{proof}

\begin{proof}[Proof of Theorem \ref{thm:pathconnected}] First note that $\calM_1$ %\cong \SO(3) \times B^5$
is indeed path-connected---see the proof of \cite[Proposition 6.2]{taubespathconnected} on p. 381. Since $\calM_k \simeq \calN_k$ by our Theorem \ref{thm:fourspheregap}, path-connectedness of $\calM_k$ now follows from Theorem \ref{thm:taubespathconnected} by induction.
\end{proof}



\begin{cor}[Atiyah-Drinfeld-Hitchin-Manin \cite{adhm}]\label{cor:adhm} For $r = 2,$ the space $\calM_k$ consists only of instantons arising from ADHM data. 
\end{cor}
\begin{proof} Let $\calM'_k \subset \calM_k$ denote the subset of instantons that arise from the ADHM construction (see \cite[Ch. II]{atiyahgeometry} for an elementary description). It is established in Donaldson-Kronheimer \cite[\S 3.4.2-4]{donkron} that $\calM'_k$ is a closed submanifold of $\tilde{\calB}_k$ of dimension $8k.$ Meanwhile, according to the Atiyah-Hitchin-Singer Theorem \cite{ahs}, the space $\calM_k$ is also a closed submanifold of $\tilde{\calB}_k$ of the same dimension. Therefore $\calM'_k$ is both closed and open relative to $\calM_k;$ since $\calM_k$ is connected, the two must be equal. %On the other hand, it is possible to check (see Atiyah?) that the space of ADHM data of charge $k$ is also smooth of the same dimension. The map taking ADHM data to $\calM_k$ is also proper (how to see this?). Hence its image is open and closed in $\calM_k,$ which is path-connected, therefore equal to all of $\calM_k.$
\end{proof}

\vspace{10mm}

\section{Deformation to instantons on quaternion-K\"ahler manifolds}




\begin{comment}
    We will use the following generalization of Theorem \ref{thm:F+boundimpliesFbound}.
  \begin{thm}[Cor. 7.3 of \cite{oliveirawaldron}]\label{thm:qKF+boundimpliesFbound}
    Let \(M\) be a quaternion-K{\"a}hler manifold of real dimension \(n\), and suppose \(A(t)\) be a smooth solution of (\ref{ymf}) on \([0, T)\) with \(A(0)\) pseudoholomorphic. Let \(0 \leq t_1 \leq t_2 < T\), and set
    \begin{align*}
       E &:= \|F(0)\|^2 \\
       K &:= \frac{\frac{n}{2}+4}{3}\int_{t_1}^{t_2} \|F^+(t)\|_\infty \, dt.
    \end{align*} There exists positive constants \(C_n\), depending on \(n\), and  \(R_0\), depending on \(E\) and the geometry of \(M\), such that
    \begin{align}\label{qkF+boundimpliesFbound:bound}
        \|F(t_2)\|_\infty \leq \max\left\{R_0^{-2}, e^{C_n\max\left\{1, K \right\}}\|F(t_1)\|_\infty \right\}.
    \end{align}
\end{thm}
\end{comment}




We first relate the \(M^{2, 4}\) norm to the monotone %monontone 
quantity \(\Phi\) from \cite[Def. 5.5]{oliveirawaldron}, which we now recall.  Let \(\varphi \in C_c^\infty([0, 1))\) be a standard cutoff for \([0, \frac{1}{2}]\), let \(x, x_1 \in M\), and let
$$\rho_1 = \min\{\text{inj}(M), 1\}.$$
Set \(\varphi_{x_1, \rho_1}(y) := \varphi\left(\frac{d(x_1, y)}{\rho_1}\right)\), where \(d\) is the Riemannian distance. For a function \(f\) on \(M\), we set
    \begin{align*}
        \Phi_{x_1, \rho_1}(f, R, x) := \frac{R^{4-n}}{(4\pi)^{\frac{n}{2}}} \int_M |f(y)|^2 e^{-\frac{d(x, y)^2}{4R^2}} \varphi_{x_1, \rho_1}(y) \, dV(y).
    \end{align*}
\begin{lemma}\label{lemma:morreycontrolsmonotonicity}
    There exists \(C_{(\ref{morreycontrolsmonotonicity:bound})} \geq 1\), depending only on the geometry of \(M\), such that for \(R \in (0, 1], q \in [2, \infty],\) and \(\lambda \in [0, n]\), we have
    \begin{align}\label{morreycontrolsmonotonicity:bound}
        \Phi_{x, \rho_1}(f, R, x) \leq C_{(\ref{morreycontrolsmonotonicity:bound})} R^{4-\frac{2\lambda}{q}} \|f\|_{q, \lambda}^2.
    \end{align}
\end{lemma}
\begin{proof}
Since \(M\) is compact, for \(R \in (0, 1]\) the function \((4\pi R^2)^{-\frac{n}{2}}e^{-\frac{d(x, y)^2}{4R^2}}\) approximates the heat kernel \(H(x, y, R^2)\) of \(M\). In particular, the proof of \cite[Proposition 3.2]{estimatespaper} in the case \(s_1 = \frac{q}{2}, s_2 = \infty,\) and \(t = R^2\) implies that 
\begin{align*}
    \Phi_{x, \rho_1}(f, R, x) \leq R^4C_{(\ref{morreycontrolsmonotonicity:bound})}(R^2)^{-\frac{\lambda}{q}}\||f|^2\|_{\frac{q}{2}, \lambda} \leq C_{(\ref{morreycontrolsmonotonicity:bound})}R^{4-\frac{2\lambda}{q}} \|f\|_{q, \lambda}^2.
\end{align*}
\end{proof}
    \begin{comment}
    {\color{red} WARNING, ARGUMENT INCOMPLETE, FOUND BETTER ONE}
     We also define
    \begin{align*}
        \|f\|_{q, \lambda, \rho_1} := \sup_{x \in M, 0 < R \leq \rho_1} \left(R^{\lambda - n} \int_{B_R(x)} |f|^q\right)^{\frac{1}{q}}.
    \end{align*}
    In the sequel, \(C\) may increase from line to line.
     For \(\rho_0 \leq 1\) sufficiently small, there exists \(C\) such that for all \(x \in M\) and \(R \in (0, K_{\rho_0} \rho_0]\), \(B_{2R}(x) \setminus B_R(x)\) can be covered by \(C\) balls of radius \(R\). For \(R \in (0, \rho_0]\), we estimate \(\Phi_{x, \rho_0}(f, R, x)\) by breaking the integral into the following three regions:
    \begin{enumerate}
        \item \(B_R(x)\);
        \item \(B_{K_RR}(x) \setminus B_R(x)\);
        \item \(B_{\rho_0}(x) \setminus B_{K_RR}(x)\).
    \end{enumerate} 
    \emph{Region (1):} We have
    \begin{align*}
        R^{4-n} \int_{B_R(x)} |f|^2 e^{-\frac{d(x, y)^2}{4R^2}} \, \mu_g(y) \leq \|f\|_{2, 4, \rho_0}^2.
    \end{align*} 
    \emph{Region (2):} Let \(J\) be the smallest positive integer such that
    \begin{align*}
        K_R \leq 2^J.
    \end{align*} Let \(i\) be an integer with \(0 \leq i \leq J\). 
    For \(r \geq 0\),
    \begin{align*}
        e^{-r} \leq C r^{-\frac{n}{2}},
    \end{align*} so for \(y\) in the annulus \(B_{2^{i+1}R}(x) \setminus B_{2^iR}(x)\) we have 
    \begin{align*}
        e^{-\frac{d(x, y)^2}{4R^2}} \leq C\left(\frac{(2^iR)^2}{4R^2}\right)^{-\frac{n}{2}} = C2^{-n(i-1)}.
    \end{align*}
    Hence
    \begin{align*}
        R^{4-n}\int_{B_{K_RR}(x) \setminus B_R(x)} |f|^2 e^{-\frac{d(x, y)^2}{4R^2}} &\leq R^{4-n}\sum_{i = 0}^{J-1}  \int_{B_{2^{i+i}R}(x) \setminus B_{2^iR}(x)} |f|^2 e^{-\frac{d(x, y)^2}{4R^2}} \\
        &\leq \sum_{i = 0}^{J-1} C2^{-n(i-1)}R^{4-n} \int_{B_{2^{i+i}R}(x) \setminus B_{2^iR}(x)} |f|^2  \\
        &\leq \sum_{i = 0}^J C 2^{-n(i-1)} C2^{(n-4)i} \|f\|_{2, 4, \rho_0}^2 \\
        &= 2^nC^2 \|f\|_{2, 4, \rho_0}^2 \sum_{i = 0}^J 2^{-4i} \\
        &\leq 2^{n+1}C^2\|f\|_{2, 4, \rho_0}^2.
    \end{align*} 
    \emph{Region (3):} Note that for \(y\) such that \(d(x, y) \geq K_RR\),
    \begin{align*}
        R^{4-n}e^{-\frac{d(x, y)^2}{4R^2}} \leq 1.
    \end{align*} Therefore,
    \begin{align*}
        R^{4-n} \int_{B_{\rho_0}(x) \setminus B_{K_RR}(x)} |f|^2 e^{-\frac{d(x, y)^2}{4R^2}} \leq \int_{B_{\rho_0}(x)} |f|^2 \leq \rho_0^{n-4}\|f\|_{2, 4, \rho_0}^2. 
    \end{align*} \par 
    Summing up the estimates for regions (1)-(3), we deduce that
    \begin{align*}
        \Phi_{x, \rho_0}(f, R, x) \leq C\|f\|_{2, 4, \rho_0}^2,
    \end{align*}  as desired.
    \end{comment}



We recall the definition of a neighborhood deformation retract.
\begin{defn}\label{defn:nbddefret}
    Given a topological space \(X\), a closed subset \(B \subset X\) is called a neighborhood deformation retract of \(X\) if there is a continuous function \(f : X \to [0, 1]\), with \(B = f^{-1}(0)\), and a homotopy \(G : X \times [0, 1] \to X\) satisfying \(G_0 = \text{id}\), \(G_s|_B = \text{id}\) for \(s \in [0, 1]\), and \(G_1(x) \in B\) if \(f(x) < 1\).
\end{defn}
 We will prove the following detailed version of Theorem \ref{thm:quaternionkahlergap}.
\begin{thm}[Cf. Theorem \ref{thm:quaternionkahlergap}]\label{thm:detailedquaternionkahlergap}
  Let $M$ be a compact quaternion-K{\"a}hler manifold of dimension \(n = 4m\) and let $k \geq 2m.$%\(k > \lceil\frac{n}{2}-1\rceil\).  
  There exists \(\delta_0 \in (0, 1)\) and \(C_{(\ref{detailedquaternionkahlergap:estimatepapercriticalbounds})} > 1\),
  %0\), 
  depending on the geometry of \(M\), as follows.

  Let $A_0$ be an \(H^k\) pseudo-holomorphic connection on a \((G, \rho)\)-bundle \(E \to M\), for which
\begin{align}\label{detailedquaternionkahlergap:F+initiallysmall}
      \|F_{A_0}^+\|_{2, 4} < \delta \leq \delta_0.
  \end{align} 
  Then the solution of (\ref{ymf}) starting from \(A_0\) exists for all time, and we have the estimate 
  \begin{align}\label{detailedquaternionkahlergap:estimatepapercriticalbounds}
      \|F^+(t)\|_{2, 4} + \min\{1, t\}\|F^+(t)\|_\infty \leq C_{(\ref{detailedquaternionkahlergap:estimatepapercriticalbounds})}\delta.
  \end{align} 
  Suppose further that \(M\) has positive scalar curvature. Then the solution converges exponentially to an instanton as $t \to \infty,$ and the space of \(H^k\) %\(\mathfrak{sp}(m)\)-
  instantons is a neighborhood deformation retract of the space of \(H^k\) pseudo-holomorphic connections.
  
  Finally, let $E_0 > 0,$ $q \in (2, \infty],$ and \(\lambda \in [0, 4)\). Let $\varepsilon_0, R_0$ be as in \cite[Theorem 6.2]{oliveirawaldron}. Assuming that $\delta > 0$ is sufficiently small, depending also on $E_0,q,$ and $\lambda,$ the following estimate is true. % for \(R_0 > 0\) small depending on \(E\) and the geometry of \(M\), and \(\varepsilon_0 > 0\) small depending on \(E_0\) and \(n\). 
  In addition to (\ref{detailedquaternionkahlergap:F+initiallysmall}), suppose that
  \begin{align}\label{detailedquaternionkahlergap:hypotheses1}
      \|F_{A_0}\|_{2,4} + \| F^+_{A_0} \|_{q,\lambda} < E_0.
  \end{align}
 Given $x \in M,$ let $R \in \left( 0, R_0 \RB$ be such that
  \begin{align}\label{detailedquaternionkahlergap:hypotheses2}
      \Phi_{x, \rho_1} \left( |F_{A_0}|, R, x \right) \leq \frac{\eps_0}{2}.
  \end{align}
  We then have
  \begin{align}\label{detailedquaternionkahlergap:regularityscale}
      |F(x,t)| \leq \frac{C_n}{R^2}
  \end{align}
  for $(x,t) \in B_{R/2}(x) \times \LB R^2, \infty \RB.$
\end{thm}

  \begin{proof}%[Proof of Theorem \ref{thm:detailedquaternionkahlergap}]
  In the sequel, \(C > 0\) is a constant depending on the geometry of \(M\) which may increase in each appearance. We first prove the bound (\ref{detailedquaternionkahlergap:estimatepapercriticalbounds}). Since the hypothesis (\ref{detailedquaternionkahlergap:F+initiallysmall}) and the conclusion (\ref{detailedquaternionkahlergap:estimatepapercriticalbounds}) are gauge-invariant, we WLOG assume that \(A_0\) is smooth, in view of item \ref{wellposedness:gaugeequivalenttosmooth} of Theorem \ref{thm:wellposedness}. By \cite[Prop. 4.1]{oliveirawaldron}, \(A(t)\) is pseudo-holomorphic as long as it exists. Let \([0, T)\) be the interval of existence of \(A(t)\). By the Weitzenb{\"o}ck formula (\ref{preliminaries:qkweitzenbock}) for \(\Omega_+^2(\mathfrak{g}_E)\) and Kato's inequality,
  \begin{align*}
      \left(\frac{\partial}{\partial t} - \Delta \right) |F^+| \leq B_1|F^+|^2 + B_2|F^+|,
  \end{align*} where \(B_1 > 0\) is universal and \(B_2 > 0\) depends on the geometry of \(M\). Thus by \cite[Theorem 1.2]{estimatespaper}, (\ref{detailedquaternionkahlergap:estimatepapercriticalbounds}) holds provided
  \begin{align*}
      \|F_{A_0}^+\|_{2, 4} + \sup_{0 \leq t < T} \|F^+(t)\| < \delta_0
  \end{align*} for \(\delta_0\) small enough. Now, on a compact manifold of dimension \(n \geq 4\), the \(L^2\) norm is dominated by the \(M^{2, 4}\) norm up to a factor of \(C\). Then since \(\|F^+(t)\|\) is nonincreasing in view of (\ref{preliminaries:qkchernweil}), we have that
    \begin{align*}
        \sup_{0 \leq t < T} \|F^+(t)\| \leq C\|F_{A_0}^+\|_{2, 4}.
    \end{align*} Thus (\ref{detailedquaternionkahlergap:estimatepapercriticalbounds}) indeed holds. Since \(\|F^+(t)\|_\infty \in L_{\text{loc}}^1((0, T))\), we deduce from \cite[Cor. 7.3]{oliveirawaldron} and item (\ref{wellposedness:blowupcharacterization}) of Theorem \ref{thm:wellposedness} that \(A(t)\) exists for all time, with \(\|F(t)\|_\infty\) blowing up at most exponentially as \(t \to \infty\). \par 
    We next specialize to the case that \(M\) has positive scalar curvature. 
    In the sequel, \(c > 0\) is a constant depending on \(M\) which may decrease in each appearence but importantly will remain positive. The evolution equation for \(F\) and the Weitzenb{\"o}ck formula (\ref{preliminaries:qkweitzenbock}) imply
    \begin{align*}
        \frac{1}{2}\frac{d}{dt} \int |F^+|^2 &= -\int \langle \nabla^\ast \nabla  F^+, F^+\rangle + \int \langle \llbracket F^+, F^+\rrbracket, F^+\rangle - 2c \int |F^+|^2 \\
        &=  -\int| \nabla  F^+|^2 + \int \langle \llbracket F^+, F^+\rrbracket, F^+\rangle - 2c \int |F^+|^2 \\
        &\leq -\int| \nabla  F^+|^2 + C\int |F^+|^3 - 2c\int |F^+|^2.
    \end{align*} By H{\"o}lder's inequality,
    \begin{align*}
        \int |F^+|^3 \leq \left(\int |F^+|^{\frac{2n}{n-2}}\right)^{\frac{n-2}{n}}\left(\int |F^+|^{\frac{n}{2}}\right)^{\frac{2}{n}}.
    \end{align*}By Kato's inequality and the Sobolev inequality on functions,
    \begin{align*}
       \|F^+\|_{\frac{2n}{n-2}}^2 \leq C(\|\nabla F^+\|^2 + \|F^+\|^2),
    \end{align*} where again \(C\) is independent of the reference connection. Hence if \(\delta_0\) is sufficiently small, (\ref{detailedquaternionkahlergap:estimatepapercriticalbounds}) implies that for \(t \geq 1\),
    \begin{align*}
        \frac{d}{dt} \int |F^+|^2 \leq -2c \int |F^+|^2.
    \end{align*} Thus by Gronwall's inequality, we have for \(t \geq 1\)
    \begin{align}\label{detailedquaternionkahlergap:F+energyexponentialdecay}
        \|F^+(t)\| \leq \|F^+(0)\|e^{-ct} \leq C\delta e^{-ct}.
    \end{align}  Then  (\ref{detailedquaternionkahlergap:estimatepapercriticalbounds}) and parabolic Moser iteration yield 
    \begin{align}\label{detailedquaternionkahlergap:F+supexponentialdecay}
        \|F^+(t)\|_\infty < C\delta e^{-ct}, \ \ \ t \geq 1.
    \end{align} Thus \(\|F^+(t)\|_\infty\) is in \(L^1([1, \infty))\), so it follows from \cite[Cor. 7.3]{oliveirawaldron} that
    \begin{align}\label{detailedquaternionkahlergap:Fsupbound}
       \sup_{t \geq 1} \|F(t)\|_\infty < \infty. 
    \end{align} Furthermore, it follows from (\ref{preliminaries:qkchernweil}) and (\ref{detailedquaternionkahlergap:F+supexponentialdecay}) that we have, analogously to (\ref{fourspheregeneralgap:rhosmall}),
    \begin{align*}
        \int_{t'}^{t} \|D^\ast F(s)\| \, ds \leq C e^{-ct'}, \ \ \ 2 \leq t' < t \leq \infty.
    \end{align*} Therefore, we have as in the proof of Theorem \ref{thm:fourspheregeneralgap} that \(A(t)\) converges in \(C^\infty\) to an instanton exponentially as \(t \to \infty\). \par 
    We next show that the subset of instantons is a neighborhood deformation retract of the space of pseudo-holomorphic connections. 
    %Recall that given a topological space \(X\), a closed subset \(B \subset X\) is a neighborhood deformation retract of \(X\) if there is a continuous function \(f : X \to [0, 1]\), with \(B = f^{-1}(0)\), and a homotopy \(G : X \times [0, 1] \to X\) satisfying \(G_0 = \text{id}\), \(G_s|_B = \text{id}\) for \(s \in [0, 1]\), and \(G_1(x) \in B\) if \(f(x) < 1\). \par 
    Let
    \begin{align*}
        X_1 &:= \{A \in H^k(\mathcal{A}_E) \mid A \text{ is pseudo-holomorphic}\} \\
        X_0 &:= \{A \in X_1 \mid F_A^+ = 0\} \\
        U_1 &:= \{A \in X_1 \mid \|F_A^+\|_{2, 4} < C_{(\ref{detailedquaternionkahlergap:estimatepapercriticalbounds})}^{-1} \delta_0\} \\
        U_2 &:= \{A \in X_1 \mid \|F_A^+\|_{2, 4} < \delta_0\}.
    \end{align*} Recall that $C_{(\ref{detailedquaternionkahlergap:estimatepapercriticalbounds})} > 1.$ Let \(h : [0, \infty) \to [0, 1]\) be a smooth, nondecreasing function such that \(h(0) = 0\) and \(h([C_{(\ref{detailedquaternionkahlergap:estimatepapercriticalbounds})}^{-1}\delta_0, \infty)) = 1\). Let \(\tilde{h} : [0, \infty) \to [0, 1]\) be a smooth, nonincreasing function such that \(\tilde{h}([0, C_{(\ref{detailedquaternionkahlergap:estimatepapercriticalbounds})}^{-1}\delta_0]) = 1\) and \(\tilde{h}([\delta_0, \infty)) = 0\). Let $\varphi_A(t)$ be as in the proof of Theorem \ref{thm:fourspheregeneralgap}. Define
    \begin{align*}
        f &: X_1 \to [0, 1] \\
        f(A) &:= h(\|F_A^+\|_{2, 4}) \\
        H &: X_1 \times [0, 1] \to [0, 1] \\
        H(A, s) &:= \min\{\tilde{h}(\|F_A^+\|_{2, 4}), s\} \\
        \Psi(A, s) &:= 
        \begin{cases}
            \varphi_A\left(\frac{H(A, s)}{1-H(A, s)}\right) & 0 \leq s < 1, \\
            \lim_{t \to \infty} \varphi_A(t) & H(A, s) = 1.
        \end{cases}
    \end{align*} Note that \(\Psi(U_1 \times [0, 1]) \subset U_2\) by (\ref{detailedquaternionkahlergap:estimatepapercriticalbounds}). By item (\ref{wellposedness:wellposedness}) of Theorem \ref{thm:wellposedness} and Prop. \ref{thm:continuityatinfinitetime}, \(\Psi\) is continuous. Thus \(X_0\) is a neighborhood deformation retract of \(X_1\) with homotopy \(\Psi\) and associated function \(f\). \par 
    We finally establish the refined estimate (\ref{detailedquaternionkahlergap:regularityscale}) on the full curvature of \(A(t)\) using \cite[Theorem 6.2]{oliveirawaldron}. Recall the constants \(\varepsilon_0, R_0\) from the statement of \cite[Theorem 6.2]{oliveirawaldron}. By (\ref{detailedquaternionkahlergap:hypotheses1}) and Lemma \ref{lemma:morreycontrolsmonotonicity}, %(\ref{lemma:morreycontrolsmonotonicity}),
    the required energy and entropy bounds on $A_0$ \cite[(6.7-8)]{oliveirawaldron} are satisfied. Moreover, by (\ref{detailedquaternionkahlergap:hypotheses2}) we have
    \begin{align*}
        \Phi_{x, \rho_1}(|F_{A_0}|, R, x) \leq \frac{\varepsilon_0}{2}.
    \end{align*} Thus, to obtain hypothesis \cite[(6.9)]{oliveirawaldron}, we just need to show
    \begin{align*}
        \kappa \int_0^\infty \|F^+(t)\|_\infty \, dt < \frac{\varepsilon_0}{2},
    \end{align*} where \(\kappa\) depends only on \(n\). By assumption
    \begin{align*}
        \|F_{A_0}^+\|_{q, \lambda} \leq E_0,
    \end{align*} and we have
    \begin{align*}
        \alpha := \frac{\lambda}{2q} < 1.
    \end{align*} Thus by \cite[(1.7)]{estimatespaper}, we have for some \(C' > 0\), depending on \(E_0, \alpha,\) and the geometry of \(M\), that
    \begin{align*}
        \|F^+(t)\|_{\infty} \leq \min\{C't^{-\alpha}, C\delta t^{-1}\} + C\delta. 
        %C't^{-\alpha} \wedge C\delta t^{-1} + C\delta.
    \end{align*}  Since \(\alpha < 1\), we may choose \(T' \in (0, 1)\) such that 
    \begin{align*}
        \int_0^{T'} \|F^+(t)\|_\infty \, dt < \frac{\varepsilon_0}{4\kappa}.
    \end{align*} On the other hand, if \(\delta\) is small enough depending on \(E_0, \alpha,\) and the geometry of \(M\),  (\ref{detailedquaternionkahlergap:F+supexponentialdecay}) implies
    \begin{align*}
        \int_{T'}^\infty \|F^+(t)\|_\infty < \frac{\varepsilon_0}{4\kappa}.
    \end{align*} Hence
    \begin{align*}
        \int_0^\infty \|F^+(t)\|_\infty \, dt < \frac{\varepsilon_0}{2\kappa}. %\frac{\varepsilon_0}{2}.
    \end{align*} The desired bound (\ref{detailedquaternionkahlergap:regularityscale}) now follows from \cite[Theorem 6.2]{oliveirawaldron}.
    \begin{comment}
    We first bound \(\Phi\). Since \(k > \frac{n}{2}-1\), it follows from Remark \ref{rmk:Hqcontrolsmorrey} that the \(H^k\) norm controls some Morrey norm \(M^{q, \lambda}\) with \(q > 2\) and \(\lambda < 4\). In particular, we have
    \begin{align*}
        \alpha := \frac{\lambda}{2q} < 1.
    \end{align*} Hence for some \(E_0 > 0\),
    \begin{align*}
        \left\|F_{A_0}\right\|_{q, \lambda} + \|F_{A_0}\| \leq E_0.
    \end{align*} Let \(R_0 \in (0, 1]\) be as in the statement of \cite[Theorem 6.2]{oliveirawaldron}. We have that \(R_0\) depends on \(K\) and the geometry of \(M\). By Lemma \ref{morreycontrolsmonotonicity:bound}, we have for \(R \in (0, R_0)\)
    \begin{align*}
        \Phi_{x, \rho_1}(|F_{A_0}|, R, x) \leq C_{(\ref{morreycontrolsmonotonicity:bound})} R^{4(1-\alpha)}K^2.
    \end{align*} Thus if \(R\) is small enough depending on \(\alpha, K,\) and the geometry of \(M\), 
    \begin{align*}
        \Phi_{x, \rho_1}(|F_{A_0}|, 2R, x) \leq 1.
    \end{align*} Thus \(\varepsilon > 0\) depends only on \(n\). Hence by taking \(R\) smaller if need be, we have
    \begin{align*}
        \Phi_{x, \rho_1}(|F_{A_0}|, R, x) < \frac{\varepsilon_0}{2}.
    \end{align*}\par
    \end{comment} 
\end{proof}

\vspace{10mm}

%\section{Morrey-norm gap theorem in general dimension}
\section{Deformation to flat connections}


\begin{proof}[First proof of Theorem \ref{thm:flatgap}] 
As in the proof of Theorem \ref{thm:quaternionkahlergap}, we WLOG assume \(A_0\) is smooth. By \cite[Theorem 1.2]{estimatespaper}, we have for \(\delta_1 > 0\) small enough, depending only on the geometry of \(M\), that if
\begin{align*}
    \|F_{A_0}\|_{2, 4} < \delta \leq \delta_1,
\end{align*} then
\begin{align}\label{flatgap:estimatepapercriticalbounds}
    \|F(t)\|_{2, 4} + \min\{1, t\}\|F(t)\|_\infty \leq C_{(\ref{flatgap:estimatepapercriticalbounds})} \delta, \quad t > 0.
\end{align} Here, \(C_{(\ref{flatgap:estimatepapercriticalbounds})}\) only depends on the geometry of \(M\). The long-time existence statement of the theorem now follows in view of item (\ref{wellposedness:blowupcharacterization}) of Theorem \ref{thm:wellposedness}. \par
We next prove convergence to a flat connection. %Let $\mathcal{F}$ denote the (possibly empty) space of flat connections on $E \to M.$
It suffices to prove the statement for vector bundles of a given rank $r.$ Suppose for contradiction that there does not exist \(\delta_1 > 0\), now depending on $r$ in addition to \(M\), such that for any smooth connection \(A_0\) on a bundle $E \to M$ of rank $r$ with 
\begin{align*}
    \|F_{A_0}\|_{2, 4} < \delta_1,
\end{align*} the solution \(A(t)\) of (\ref{ymf}) starting at \(A_0\) satisfies:
\begin{center}
    \(A(1)\) is contained in a gauge-invariant neighborhood \(U \supset \mathcal{F}\) as in Lemma \ref{lemma:uniformlojasiewicz}. \ \ \((\ast)\)
\end{center}
Then there exists a sequence of smooth connections \(A_i\) with 
\begin{align}\label{flatgap:morreygoingtozero}
    \|F_{A_i}\|_{2, 4} \searrow 0, 
\end{align} but for which the solution \(\tilde{A}_i(t)\) of (\ref{ymf}) starting at \(A_i\) does not satisfy \((\ast)\). In view of (\ref{flatgap:estimatepapercriticalbounds}), (\ref{flatgap:morreygoingtozero}) implies 
\begin{align*}
    \lim_{i \to \infty} \sup_{\frac{1}{2} \leq t < \infty} \|F_i(t)\|_\infty = 0.
\end{align*} Thus by standard strong Uhlenbeck compactness arguments for Yang-Mills flow, e.g. \cite[Theorem 1.3]{waldronuhlenbeck}, we may pass to a subsequence, also labeled by \(i\), for which there exist smooth gauge transformations \(u_i\) such that \(u_i(\tilde{A}_i(1))\) smoothly converges to a flat connection \(A_\infty\) on $E \to M.$ In particular, a neighborhood \(U\) of $A_\infty$ as stipulated in Lemma \ref{lemma:uniformlojasiewicz} exists, and furthermore, gauge-invariance of \(U\) implies \(\tilde{A}_i(1) \in U\) for \(i\) large enough. Thus \((\ast)\) holds, which yields the desired contradiction. \par

We may therefore let $\delta_1$ be such that $(\ast)$ holds for $A(1),$ so that the desired convergence to a flat connection follows from Lemma \ref{lemma:uniformlojasiewicz}. The argument that $\mathcal{F}$ is a neighborhood deformation retract of \(H^k(\mathcal{A}_E)\) is essentially identical to the argument given in the proof of Theorem \ref{thm:detailedquaternionkahlergap}, so we omit the details.
\begin{comment}
    We finally explain the neighborhood deformation retract property. Let
    \begin{align*}
        X_1 &:= H^k(\mathcal{A}_E) \\
        X_0 &:= \{A \in X_1 \mid F_A = 0\} \\
        U_1 &:= \{A \in X_1 \mid \|F_A\|_{M^{2, 4}} < C_{(\ref{flatgap:estimatepapercriticalbounds})}^{-1} \delta_1\} \\
        U_2 &:= \{A \in X_1 \mid \|F_A\|_{M^{2, 4}} < \delta_1\}.
    \end{align*} Let \(h : [0, \infty) \to [0, 1]\) be a smooth, increasing function such that \(h(0) = 0\) and \(h([C_{(\ref{flatgap:estimatepapercriticalbounds})}^{-1}\delta_1, \infty)) = 1\). Let \(\tilde{h} : [0, \infty) \to [0, 1]\) be a smooth, decreasing function such that \(\tilde{h}([0, C_{(\ref{flatgap:estimatepapercriticalbounds})}^{-1}\delta_1]) = 1\) and \(\tilde{h}([\delta_1, \infty)) = 0\). Define
    \begin{align*}
        f &: X_1 \to [0, 1] \\
        f(A) &:= h(\|F_A\|_{M^{2, 4}}) \\
        H &: X_1 \times [0, 1] \to [0, 1] \\
        H(A, s) &:= \min\{\tilde{h}(\|F_A\|_{M^{2, 4}}), s\} \\
        \Psi(A, s) &:= 
        \begin{cases}
            \varphi_A\left(\frac{H(A, s)}{1-H(A, s)}\right) & 0 \leq s < 1 \\
            \lim_{t \to \infty} \varphi_A(t) & H(A, s) = 1.
        \end{cases}
    \end{align*} Note that \(\Psi(U_1 \times [0, 1]) \subset U_2\) by (\ref{flatgap:estimatepapercriticalbounds}). By item (\ref{wellposedness:wellposedness}) of Theorem \ref{thm:wellposedness} and Prop. \ref{thm:continuityatinfinitetime}, \(\Psi\) is continuous. Thus \(X_0\) is a neighborhood deformation retract of \(X_1\) with homotopy \(\Psi\) and associated function \(f\).
\end{comment}
\end{proof}

As mentioned in the introduction, we can also give a proof of Theorem \ref{thm:flatgap} using monotonicity and $\varepsilon$-regularity in place of \cite[Theorem 1.2]{estimatespaper}.

\begin{proof}[Second proof of Theorem \ref{thm:flatgap}]
     As before, we WLOG assume \(A_0\) is smooth. We first establish long-time existence of \(A(t)\). Let $\rho_1 := \min\{\text{inj}(M), 1\}.$ Suppose \(A(t)\) exists on \([0, T]\) with \(T \in (0, \rho_1^2]\). Let \(\varepsilon_0, R_0\) be as in the statement of $\varepsilon$-regularity \cite[Theorem 6.2]{oliveirawaldron} (with \(\gamma = 1\)) with \(E_0 = E = 1\). By Lemma \ref{lemma:morreycontrolsmonotonicity} and $\varepsilon$-regularity \cite[Theorem 6.2]{oliveirawaldron} (with \(\gamma = 1\)), we have that if \(\delta_1\) is sufficiently small, depending on the geometry of \(M\), then \(\|F(T)\|_\infty < \infty\). Consequently, we may assume \(T \geq R_0^2\), and we have 
    \begin{align*}
        \|F(R_0^2)\|_\infty \leq \frac{C_n}{R_0^2}.
    \end{align*}
    Thus the fact that the energy is non-increasing implies that for \(\delta_1\) sufficiently small,
    \begin{align*}
        \sup_{x \in M, t \geq R_0^2} \Phi_{x, \rho_1}(|F(t)|, R_0, x) < \varepsilon_0.
    \end{align*} Hence, it follows from \cite[Theorem 6.2]{oliveirawaldron} again that \(A(t)\) is uniformly bounded at any time \(t \geq R_0^2\). Therefore, \(A(t)\) exists for all time. \par 
    The convergence to a flat connection follows as in the first proof of Theorem \ref{thm:flatgap}, which requires \(\delta_1\) to be small depending on $\rk(E)$ in addition to \(M\). \par 
    Finally, to show that the subset of flat connections is a neighborhood deformation retract, we just need the existence of some \(C_0\), depending only on the geometry of \(M\), such that
    \begin{align*}
        \|F(t)\|_{M^{2, 4}} \leq C_0 \delta,  \ \ \ t \geq 0.
    \end{align*} This bound follows from monotonicity, e.g. \cite[Theorem 5.7]{oliveirawaldron} (with \(\gamma = 1\)), for short time, and from $\varepsilon$-regularity for later times.
\end{proof}

\begin{cor}\label{cor:lowYMandtimeoneimpliesconvergence}
    Given $E \to M,$ with \(M\) of dimension \(n \geq 4\), there exists $\delta_2 > 0$, depending only on $\rk(E)$ and the geometry of \(M\), as follows. Let $A_0$ be an \(H^k\) connection with $\YM(A_0)  < \delta \leq \delta_2,$ and suppose that the solution $A(t)$ of (\ref{ymf}) with $A(0) = A_0$ exists for at least time $T > \delta_2^{-1}\delta^{\frac{2}{n-4}}$. Then $T = \infty$ and \(A(t)\) %the flow 
    converges to a flat connection.
\end{cor}
\begin{proof}%[Proof of Corollary \ref{cor:lowYMandtimeoneimpliesconvergence}]
    Suppose \(A(t)\) exists on \(\left[0, \delta_2^{-1}\delta^{\frac{2}{n-4}}\right]\). It follows from monotonicity \cite[Theorem 5.7]{oliveirawaldron} (with \(\gamma = 1\)) and the definition of \(\Phi\) that if \(\delta_2\) is sufficiently small, then for some \(\tau\), we have
    \begin{align*}
        \left\|F\left(\tau\right)\right\|_{M^{2,4}} \leq C\sup_{x \in M, 0 < R \leq \rho_1} \Phi_{x, \rho_1}\left(\left|F(\tau)\right|, R, x\right) < \delta_1.
    \end{align*} Thus the corollary now follows from Theorem \ref{thm:flatgap}.
\end{proof}



\begin{proof}[Proof of Corollary \ref{cor:Naito}]
%Suppose for a contradiction that the solution is global. Then if the energy is initially small, the monotonicity formula implies that at say time \(t = 1\), the Morrey norm of the curvature will be small. Then by Theorem \ref{thm:Morreygap}, there exists a flat connection on \(E\), which is a contradiction.
This follows from the contrapositive of Corollary \ref{cor:lowYMandtimeoneimpliesconvergence}.
\end{proof}





\begin{cor}\label{cor:ellipticgap}
    If $A$ is a %smooth 
    Yang-Mills connection on $E \to M$ with $\YM(A) < \delta_2,$ then $A$ is flat.
\end{cor}

\begin{proof}[First proof of Corollary \ref{cor:ellipticgap}]
If $A = A_0$ is initially Yang-Mills, then $A(t) \equiv A_0,$ so this follows directly from Corollary \ref{cor:lowYMandtimeoneimpliesconvergence}.
\end{proof}



\begin{proof}[Second proof of Corollary \ref{cor:ellipticgap}]\label{ellipticgapsecondproof}
It is easy to give a purely elliptic proof of Corollary \ref{cor:ellipticgap} along the lines of the second proof of Theorem \ref{thm:flatgap}. Supposing for %a 
contradiction that there is no gap, there exists a sequence \(A_i\) of %smooth 
Yang-Mills connections with $F_i \not \equiv 0$ and \(\|F_i\| \searrow 0\), which we may assume are smooth after gauge transforming. As in \cite[Fact 2.2 and Lemma 3.1]{nakajimahigherdimensions}, it follows that \(\|F_i\|_\infty \searrow 0\), together with bounds on all covariant derivatives. By strong Uhlenbeck compactness, we may pass to a subsequence, also labeled by \(i\), such that there exist smooth gauge transformations \(u_i\) for which \(u_i(A_i)\) smoothly converge to a flat connection \(A_\infty\). But then the Lojasiewicz inequality, (\ref{lojasiewicz:lojasiewiczinequality}) or \cite[Lemma 12]{byanguniqueness}, implies that \(F_{u_i(A_i)} \equiv 0 \equiv F_{A_i} \) for \(i\) large enough, contradicting our assumption.
\end{proof}







\begin{rmk}\label{rmk:ellipticgap}
Note that Corollary \ref{cor:Naito} %and 4.3
gives another instance in which the parabolic $L^2$ gap fails, in this case at finite time, even while the elliptic $L^2$ gap holds.
We also note that Corollary \ref{cor:ellipticgap} evidently implies the \(L^{\frac{n}{2}}\) gap result published by Feehan \cite{feehanflatgap}.
\end{rmk}






%\begin{rmk}
%Corollary \ref{cor:ellipticgap} gives an elliptic \(L^2\) gap near the flat connections. However, by Naito's construction of solutions of Yang-Mills flow with initial data whose curvature is arbitrarily small in \(L^2\) and yet the solution blows up in finite time \cite{naito}, there can be no general parabolic \(L^2\) gap.
%\end{rmk}


\begin{comment}

\begin{rmk}
    One can prove Theorem \ref{thm:Morreygap} using the monotonicity formula and epsilon-regularity theorem for \(F\) in place of the results from {\color{red} ESTIMATES PAPER}; we explain this below. However, in the quaternion-K{\"a}ler setting, we do not have a monotonocity formula for \(F^+\) like we do for the full curvature. Hence we use the results of {\color{red} ESTIMATES PAPER}
\end{rmk}


\end{comment}


\vspace{10mm}


\appendix





\section{Short-time existence and well-posedness of Yang-Mills flow}

  The goal of this section is to prove an existence, uniqueness, and well-posedness result for Yang-Mills flow, Theorem \ref{thm:wellposedness}. First we give our definition of a solution to (\ref{ymf}) in the Sobolev setting.

\begin{defn}[Cf. {\cite[Def. 2.1]{struweym}}]
Let $G$ be a compact Lie group, let $
\rho : G \to \GL_K(V)$ be a faithful orthogonal/unitary representation over a $K$-vector space $V$, let \(E \to M\) be a \((G, \rho)\)-bundle over a closed Riemannian \(n\)-manifold \(M\), let \(k\) be an integer such that \(k > \frac{n}{2}-1\), and let \(\tau_0 \in (0, \infty)\). We say that \(A(t)\) is a weak solution of (\ref{ymf}) on \([0, \tau_0]\) with initial data \(A_0\) if, writing $\nabla_{A(t)} = \nabla_{\text{ref}} + a(t)$ and $\nabla_{A_0} = \nabla_{\text{ref}} + a_0,$ the following hold:
\begin{enumerate}
    \item We have \(\left(a(t), F_{A(t)}\right) \in U(\tau_0),\) where \(U(\tau_0)\) is a certain Hilbert space defined in (\ref{shorttimeexistence:definitionofU}) below.
\item \(a(0) = a_0\).
\item  For all \(\varphi \in C_c^\infty(M \times (0, \tau_0), \Omega^1(\mathfrak{g}_E))\) we have
\begin{align*}
    \int_0^{\tau_0} \int_M \left\langle a(t), -\frac{\partial \varphi}{\partial t}\right\rangle +\langle F_{A(t)}, D_{A(t)} \varphi\rangle = 0.
\end{align*}
\end{enumerate}
\end{defn}
 
 \begin{thm}\label{thm:wellposedness}
     Given $K > 0,$ there exists $\tau> 0,$ depending on $K, k, \nabla_{\text{ref}},$ and the geometry of $M,$ such that for each $A_0 \in H^k\left( \mathcal{A}_E \right)$ with $\|A_0 \|_{H^k} < K,$ there exists $T \in [\tau,  \infty]$ and a unique weak solution $A(t)$ of (\ref{ymf}) with $A(0) = A_0,$ satisfying $a(t) \in C^0_{loc}(\left[0, T \right), H^k)$ and $\left( a(t), F_{A(t)} \right) \in U^+(T),$ where \(U^+(T)\) is defined in (\ref{shorttimeexistence:localhilbertspace}) below. We have the additional regularity
     \begin{equation}
        F_{A(t)} \in L_{\text{loc}}^2([0, T), H^k), \quad D_{A(t)}^\ast F_{A(t)} \in L_{\text{loc}}^2([0, T), H^{k-1}).
        \end{equation}
        The following properties also hold:
     \begin{enumerate}
         %\item $A(t)$ is the unique weak solution of (\ref{ymf}) in $C^0_{loc}(\left[0, T \right), H^k)\) with \(A(0) = A_0\).
         %\item \(A(t)\) is smooth on \(M \times [0, t_0]\) if \(A_0\) is smooth.
         \item \label{wellposedness:wellposedness} The flow is well-posed: given $\eps > 0$ and $0 < T_0 < T,$ there exists $\delta > 0$ such that for $A_0' \in B_{\delta}(A_0) \subset H^k \left( \mathcal{A}_E \right),$ the corresponding solution $A'(t)$ exists on $\LB 0, T' \right),$ with $T' > T_0,$ and $\| A'(t) - A(t) \|_{H^k} < \eps$ for all $t \in \LB 0 , T_0 \RB.$
         \item \label{wellposedness:gaugeequivalenttosmooth} There exists a (constant-in-time) gauge transformation \(u \in H^{k+1}(\mathcal{G}_E)\) such that \(u(A(t))\) is smooth and solves (\ref{ymf}) classically for \(t > 0\).
         \item \label{wellposedness:blowupcharacterization} Either $T = \infty$ or \(\limsup_{t \nearrow T} \|F_{A(t)}\|_\infty = \infty\).
     \end{enumerate}
 \end{thm}
The first part of Theorem \ref{thm:wellposedness} is the content of Theorem \ref{thm:shorttimeexistence} in \textsection \ref{radegendim}. Item (\ref{wellposedness:wellposedness}) readily follows from Theorem \ref{thm:shorttimeexistence} and is proven in \textsection \ref{item1proof}. We prove item (\ref{wellposedness:gaugeequivalenttosmooth}) in \textsection \ref{gaugeeqtosmth} using the estimates of Theorem \ref{thm:Hkestimates} from \textsection \ref{sobolevdistest}. The proof of item (\ref{wellposedness:blowupcharacterization}) in \textsection \ref{ltecriterion} is then short given the rest of Theorem \ref{thm:wellposedness}.


\vspace{5mm}




\subsection{R{\aa}de in general dimension} \label{radegendim}
     %\section{Well-posedness of Yang-Mills flow in general dimension}
%\subsection{Short-time existence}
In dimensions two and three, R{\aa}de provided a robust short-time existence theory for Yang-Mills flow with \(H^1\) initial data \cite{radethesis, rade}. Assuming stronger regularity on the initial data, R{\aa}de's method extends straightforwardly to higher dimensions. %We remark that short-time existence theory for (\ref{ymf}) has also been developed by Struwe \cite{struweym} on compact 4-manifolds, by Oh-Tataru \cite{ohtatarucaloric} on $\R^4,$ and by Kozono-Maeda-Naito \cite{kozonomaedanaito4dglobal} on compact manifolds of general dimension, all using the Donaldson-DeTurck gauge-fixing strategy. 
\begin{thm}\label{thm:shorttimeexistence}
     Let \(n \geq 2\), and let \(q \in \R\) satisfy \(q \geq 1\) and \(q > \frac{n}{2} - 1\). Let \(E\) be a \((G, \rho)\)-bundle over a closed Riemannian \(n\)-manifold \(M\). Given a connection \(A_0 \in H^q(\mathcal{A}_E)\), there exists \(\tau_0 > 0\), depending on \(q, \|A_0\|_{H^q},\) \(\nabla_{\text{ref}}\), and the geometry of \(M\), such that the following holds. There is a unique weak solution \(A(t) \) of (\ref{ymf}) with \(A(0) = A_0,\) $a(t) \in C^0([0, \tau_0], H^q),$ and $\left(a(t), F_{A(t)} \right)$ is in the space $U(\tau_0)$ defined in (\ref{shorttimeexistence:definitionofU}) below.
     %\(\left(a(t), F_{A(t)} \right) \in U(\tau_0)\) 
     The curvature satisfies \(F_{A(t)} \in L^2([0, \tau_0], H^q)\).
     The solution \(A(t)\) is smooth on \(M \times [0, \tau_0]\) if \(A_0\) is smooth.
 \end{thm}
 
 
 \begin{proof}
     %Fix  a smooth reference \((G, \rho)\)-connection \(\nabla_{\text{ref}}\), and denote its curvature by \(F_{\text{ref}}\) and its exterior covariant derivative by \(D_{\text{ref}}\). 
     Denote the curvature of the fixed smooth reference $(G, \rho)$-connection $\nabla_{\text{ref}}$ by $F_{\text{ref}},$ and the induced covariant exterior derivative by $D_{\text{ref}}.$ %Let \(A(t) \) be a path %family of connections and write \( \nabla_{A(t)} = \nabla_{\text{ref}} + a(t) \).
     If we set \(\Omega(t) = F_{A(t)}\), then (\ref{ymf}) and %the flow equation 
     the corresponding evolution equation for the curvature can be expressed as the following system (where \(\#\) denotes various multilinear (not necessarily associative) operations whose precise forms will not be germane):
     \begin{comment}
         Suppose \(A(t) = A_{\text{ref}} + a(t)\) is a smooth solution of Yang-Mills flow. Then the flow equation and the Bianchi identity imply that \((A, F_A)\) solves the system
 \begin{align*}
     &\frac{d A}{dt} + D_A^\ast F_A = 0 \\
     &\frac{d F_A}{dt} + \Delta_A F_A = 0.
 \end{align*} This can be expanded into the system
 \begin{align*}
     &\frac{d a}{dt} + D_{\text{ref}}^\ast F_A = a \# F_A \\
     &\frac{d F_A}{dt} + \nabla_{\text{ref}}^\ast \nabla_{\text{ref}} F_A = \nabla_{\text{ref}} a \# F_A + a \# \nabla_{\text{ref}} F_A + a \# a \# F_A + F_{\text{ref}} \# F_A + Rm \# F_A,
 \end{align*} where \(\#\) denotes various multilinear (not necessarily associative) operations whose precise form will not be germane. 
     \end{comment}
 \begin{empheq}[left=\empheqlbrace]{equation}
  \begin{split}\label{shorttimeexistence:yangmillssystem}
     \hspace{0.6em} \frac{d}{dt} a &+ D_{\text{ref}}^\ast \Omega = a \# \Omega \\
     \frac{d}{dt} \Omega &+ \nabla_{\text{ref}}^\ast \nabla_{\text{ref}} \Omega = F_{\text{ref}} \# \Omega + Rm \# \Omega + \nabla_{\text{ref}} a \# \Omega + a \# \nabla_{\text{ref}} \Omega  + a \# a \# \Omega    \\
     a(0) &= a_0 \\
     \Omega(0) &= \Omega_0.
  \end{split}
\end{empheq} As in \cite[(18)]{radethesis}, we shall show existence/uniqueness of solutions to the above system for a family \(a = a(t)\) of \(\mathfrak{g}_E\)-valued 1-forms and a family \(\Omega = \Omega(t)\) of \(\mathfrak{g}_E\)-valued 2-forms, where \(\Omega\) need not be the curvature of \(\nabla_{\text{ref}} + a\); at the %very 
end, we will show that %this is the case 
\(\Omega\) is indeed the curvature of \(\nabla_{\text{ref}} + a\) if \(\Omega_0\) is the curvature of \(\nabla_{\text{ref}} + a_0\), %as in 
just like \cite[pg. 22]{radethesis}. \par
 We next describe the function spaces we shall work in. Let \(\tau_0 > 0\). For real numbers \(r\) and \(s\), we will use the shorthand \(H^r\) to denote the Hilbert Sobolev space \(H^r(T^\ast M \otimes \mathfrak{g}_E)\) or \(H^r(\Lambda^2 T^\ast M \otimes \mathfrak{g}_E)\) and \(H^{r, s}\) to denote \(H^r([0, \tau_0], H^s)\) \cite[\textsection D.1, pg. 38]{radethesis}. We will use \(H_P^{r, s}\) to denote the parabolic subspace of \(H^{r, s}\) \cite[pg. 42]{radethesis}. Set  
 \begin{align*}
     \mu &:= q - \frac{n}{2}+1 > 0 \\
     \varepsilon &:= \frac{\min\left\{1, \mu\right\}}{10} > 0.
 \end{align*} Set
\begin{align}\label{shorttimeexistence:definitionofU}
X &= \left\{(b_0, \Psi_0) \mid b_0 \in H^q, \Psi_0 \in H^{q-1}\right\} \nonumber \\
U(\tau_0) &= \left\{(b, \Psi) \mid b \in H^{\frac{1}{2} + \varepsilon, q - 2\varepsilon} \cap H^{\frac{1}{2}, q}, \Psi \in  H^{\frac{1}{2} + \varepsilon, q-1-2\varepsilon} \cap H^{-\frac{1}{2}, q +1}\right\} \\
U_P(\tau_0) &= \left\{(b, \Psi) \mid b \in H_P^{\frac{1}{2} + \varepsilon, q - 2\varepsilon} \cap H_P^{\frac{1}{2}, q}, \Psi \in  H_P^{\frac{1}{2} + \varepsilon, q-1-2\varepsilon} \cap H_P^{-\frac{1}{2}, q +1}\right\} \nonumber \\
W_P(\tau_0) &= \left\{(b, \Psi) \mid b \in H_P^{-\frac{1}{2} + \varepsilon, q - 2\varepsilon} \cap H_P^{-\frac{1}{2}, q}, \Psi \in H_P^{-\frac{1}{2} + \varepsilon, q - 1 - 2\varepsilon} \cap H_P^{-\frac{1}{2}, q- 1}\right\}. \nonumber
\end{align} Note that in the function space $U(\tau_0),$ $\Psi$ has the space-time regularity of a solution of the heat equation with $H^{q-1}$ initial data, and $b$ has one more spatial derivative than $\Psi$. This is sensible since the curvature of a solution of (\ref{ymf}) obeys a heat-type equation. \par
For \(T \in (0, \infty]\), we also set 
\begin{align}\label{shorttimeexistence:localhilbertspace}
    U^+(T) = \left\{(a, \Omega) \mid  (a, \Omega) \in U(\tau_0) \ \forall \ \tau_0 < T\right\}.
\end{align} We say \((a, \Omega)\) is a solution of (\ref{shorttimeexistence:yangmillssystem}) if \((a, \Omega) \in U(\tau_0)\), \((a(0), \Omega(0)) = (a_0, \Omega_0)\), and \((a, \Omega)\) satisfies for all 
\begin{align*}
    (\varphi, \psi) \in C_c^\infty\left(M \times (0, \tau_0), \Omega^1(\mathfrak{g}_E) \oplus \Omega^2(\mathfrak{g}_E)\right)
\end{align*} that
\begin{align*}
    &\int_0^{\tau_0} \hspace{-0.6em} \int_M \left\langle a, -\frac{\partial \varphi}{\partial t}\right\rangle +  \langle \Omega, D_{\text{ref}} \varphi\rangle =  \int_0^{\tau_0} \hspace{-0.6em} \int_M \langle a \# \Omega, \varphi\rangle \\
    &\int_0^{\tau_0} \hspace{-0.6em} \int_M \left\langle \Omega, -\frac{\partial \psi}{\partial t}\right\rangle + \langle \Omega, \nabla_{\text{ref}}^\ast \nabla_{\text{ref}} \psi \rangle =  \int_0^{\tau_0} \hspace{-0.6em} \int_M \langle \Upsilon, \psi\rangle, 
\end{align*} where \(\Upsilon\) is the RHS of the second equation in (\ref{shorttimeexistence:yangmillssystem}). \par
We now begin the proof of existence of solutions to (\ref{shorttimeexistence:yangmillssystem}). We will solve for \((a, \Omega)\) of the form \((a_1, \Omega_1) + (a_2, \Omega_2)\), where \((a_1, \Omega_1)\) solves a homogeneous initial-value problem and \((a_2, \Omega_2)\) solves an inhomogeneous system. In the sequel, we take \(\tau_0 \leq 1\). First we solve the initial-value problem
 \begin{empheq}[left=\empheqlbrace]{equation}\label{shorttimeexistence:ivp}
  \begin{split}
\frac{d}{dt} a_1 &+ D_{\text{ref}}^\ast \Omega_1 = 0 \\
\frac{d}{dt} \Omega_1  &+ \nabla_{\text{ref}}^\ast \nabla_{\text{ref}} \Omega_1 = 0 \\
a_1(0) &= a_0 \\
\Omega_1(0) &= \Omega_0,
 \end{split}
\end{empheq} with initial data \((a_0, \Omega_0) \in X\) and \((a_1, \Omega_1) \in U(\tau_0)\). In the sequel, \(C\) will denote a positive constant, depending on the geometry of \(M\), \(\nabla_{\text{ref}}\), and \(q\), and \(C\) may increase from line to line. Since \(\nabla_{\text{ref}}\) is smooth, \cite[Prop. D.9]{radethesis} implies that the initial-value problem
\begin{align*}
&\frac{d}{dt} \Omega_1  + \nabla_{\text{ref}}^\ast \nabla_{\text{ref}} \Omega_1 = 0 \\
&\Omega_1(0) = \Omega_0
\end{align*} has a unique weak solution \(\Omega_1(t) \in H^{\frac{1}{2}-r, q-1 + 2r}([0, \tau_0]), r \in \R\). In particular, for each \(\tau_0 \in (0, 1]\), the map \(\Omega_0 \to \Omega_1(t)\) yields operators
\begin{align*}
H^{q-1} &\to H^{\frac{1}{2}-r, q-1 + 2r}([0, \tau_0]) \\
H^{q-1} &\to C^0([0, \tau_0], H^{q-1}),
\end{align*} with norms bounded respectively by \(C\max\{1, \tau_0^r\}\) and \(C\). Thus, \(\Omega_1 \in  H^{\frac{1}{2} + \varepsilon, q-1-2\varepsilon} \cap H^{-\frac{1}{2}, q +1}\). Once we obtain \(\Omega_1\), we solve for \(a_1\) by
\begin{align*}
    a_1 = a_0 - \int_0^t D_0^\ast \Omega_1(s) \, ds \in H^{\frac{1}{2} + \varepsilon, q - 2\varepsilon} \cap H^{\frac{1}{2}, q}.
\end{align*} Overall, we have obtained an operator 
\begin{align*}
    M : X &\to U(\tau_0) \\
    (a_0, \Omega_0) &\to (a_1, \Omega_1),
\end{align*} and by 
 \cite[Prop. D.7 and D.9]{radethesis} we have
\begin{align*}
    \|M\| \leq C \tau_0^{-\varepsilon}.
\end{align*} \par
Now that we have \((a_1, \Omega_1)\), we solve the inhomogeneous system 
\begin{empheq}[left=\empheqlbrace]{equation}\label{inhomogeneoussystem}
  \begin{split}
     \hspace{0.6em}
\partial_t a_2 + D_{\text{ref}}^\ast \Omega_2 &= (a_1 + a_2) \# (\Omega_1 + \Omega_2)  \\
(\partial_t + \nabla_{\text{ref}}^\ast \nabla_{\text{ref}})\Omega_2 &= F_{\text{ref}} \# (\Omega_1 + \Omega_2) + Rm \# (\Omega_1 + \Omega_2) \\ 
&+ \nabla_{\text{ref}} (a_1 + a_2) \# (\Omega_1 + \Omega_2) +  (a_1 + a_2) \# \nabla_{\text{ref}}(\Omega_1 + \Omega_2) \\
&+ (a_1 + a_2) \# (a_1 + a_2) \# (\Omega_1 + \Omega_2) \\
a_2(0) &= 0 \\
\Omega_2(0) &= 0
\end{split}
\end{empheq} via \cite[Lemma B.5]{radethesis}. To do this, we will show that the operators \(L, Q_1, Q_2, Q_3\) from \cite[pg. 15]{radethesis}, given by
\begin{align*}
    L(b, \Psi) &:= \left(\frac{db}{dt} + D_{\text{ref}}^\ast \Psi, \frac{d\Psi}{dt} + \nabla_{\text{ref}}^\ast \nabla_{\text{ref}} \Psi\right) \\
    Q_1(b, \Psi) &:= (0, F_{\text{ref}} \# \Psi + Rm \# \Psi) \\
    Q_2(b, \Psi) &:= (b \# \Psi, b \# \nabla_{\text{ref}} \Psi + \nabla_{\text{ref}} b \# \Psi) \\
    Q_3(b, \Psi) &:= (0, b \# b \# \Psi),
\end{align*} extend to maps
\begin{align*}
    L &: U_P(\tau_0) \to W_P(\tau_0) \\
    Q_1 &: U(\tau_0) \to W_P(\tau_0) \\
    Q_2 &: \text{Sym}^2U(\tau_0) \to W_P(\tau_0) \\
    Q_3 &: \text{Sym}^3U(\tau_0) \to W_P(\tau_0),
\end{align*} with \(L\) invertible, and satisfy the operator norm bounds:
\begin{empheq}[left=\empheqlbrace]{equation}\label{shorttimeexistence:estimatesforexistence}
\begin{split}
    \|L^{-1}\| &\leq C \\
    \|Q_1\| &\leq C \tau_0^{1- \varepsilon} \\
    \|Q_2\| &\leq C \tau_0^{3\varepsilon} \\
    \|Q_3\| &\leq C \tau_0^{4\varepsilon}.
\end{split}
\end{empheq} \par
By \cite[Prop. D.7 and D.8]{radethesis}, \(\frac{d}{dt}\) extends to a bounded, invertible operator
\begin{align*}
    H_P^{\frac{1}{2} + \varepsilon, q - 2\varepsilon} \cap H_P^{\frac{1}{2}, q} \to H_P^{-\frac{1}{2} + \varepsilon, q - 2\varepsilon} \cap H_P^{-\frac{1}{2}, q}
\end{align*}and \(\frac{d}{dt} + \nabla_{\text{ref}}^\ast \nabla_{\text{ref}}\) extends to a bounded, invertible operator
\begin{align*}
    H_P^{\frac{1}{2} + \varepsilon, q-1-2\varepsilon} \cap H_P^{-\frac{1}{2}, q +1} \to H_P^{-\frac{1}{2} + \varepsilon, q - 1 - 2\varepsilon} \cap H_P^{-\frac{1}{2}, q- 1},
\end{align*} and the norms of these operators and their inverses are bounded by \(C\). Moreover, \(D_{\text{ref}}^\ast\) maps
\begin{align*}
     H_P^{-\frac{1}{2}, q +1} \to H_P^{-\frac{1}{2}, q}.
\end{align*} Next, by interpolation
\begin{align*}
    H_P^{\frac{1}{2} + \varepsilon, q-1-2\varepsilon} \cap H_P^{-\frac{1}{2}, q +1} \hookrightarrow H^{-\frac{1}{2} + \varepsilon, q + 1 -2\varepsilon},
\end{align*} and we further have that \(D_{\text{ref}}^\ast\) maps
\begin{align*}
    H^{-\frac{1}{2} + \varepsilon, q + 1 -2\varepsilon} \to H^{-\frac{1}{2} + \varepsilon, q -2\varepsilon}.
\end{align*} Hence \(D_{\text{ref}}^\ast\) is a bounded operator
\begin{align*}
     H_P^{\frac{1}{2} + \varepsilon, q-1-2\varepsilon} \cap H_P^{-\frac{1}{2}, q +1} \to H_P^{-\frac{1}{2} + \varepsilon, q - 2\varepsilon} \cap H_P^{-\frac{1}{2}, q},
\end{align*}with norm bounded by \(C\). 
\begin{comment}
Moreover, \(D_{\text{ref}}^\ast\) clearly is a bounded operator
\begin{align*}
     H_P^{\frac{1}{2} + \varepsilon, q-1-2\varepsilon} \cap H_P^{-\frac{1}{2}, q +1} \to H_P^{-\frac{1}{2} + \varepsilon, q - 2\varepsilon} \cap H_P^{-\frac{1}{2}, q},
\end{align*}with norm bounded by \(C\). 
\end{comment}
Since 
\begin{align*}
    L^{-1} = 
    \begin{pmatrix}
        (\frac{d}{dt})^{-1} & -(\frac{d}{dt})^{-1}D_{\text{ref}}^\ast (\frac{d}{dt} + \nabla_{\text{ref}}^\ast \nabla_{\text{ref}})^{-1} \\
        0 & (\frac{d}{dt} + \nabla_{\text{ref}}^\ast \nabla_{\text{ref}})^{-1}
    \end{pmatrix},
\end{align*} we deduce that \(\|L^{-1}\| \leq C\). \par 
Next we bound the norm of \(Q_1\). Since \(F_{\text{ref}}\) and \(Rm\) are smooth in space and constant in time, they are in \(H^{\frac{1}{2} + \varepsilon, q+\frac{n}{2}}\) with \(\|F_{\text{ref}}\|_{H^{\frac{1}{2} + \varepsilon, q+ \frac{n}{2}}} + \|Rm\|_{H^{\frac{1}{2} + \varepsilon, q+\frac{n}{2}}} \leq  C \tau_0^{-\varepsilon}\) \cite[pg. 43]{radethesis}. Also, by interpolation we have \(\Omega \in H^{\frac{1}{2}, q-1}\) \cite[(73)]{radethesis}.
By hypothesis,
\begin{align*}
    &2q+ \frac{n}{2}-1 > 0, \\
    &q - 1 \leq \min\left\{q-1, q + \frac{n}{2},  2q -1\right\}.
\end{align*} It follows from the Sobolev multiplication theorem \cite[Prop. D.7]{radethesis} that \(Q_1\) maps \(H^{\frac{1}{2} + \varepsilon, q + \frac{n}{2}} \times H^{\frac{1}{2}, q - 1} \to H^{-\frac{1}{2} + \varepsilon, q - 1}\). By \cite[(70)]{radethesis},
\begin{align*}
    Q_1(a, \Omega) \in H^{-\frac{1}{2} + \varepsilon, q - 1}
    \begin{cases}
        \inj H^{-\frac{1}{2} + \varepsilon, q-1 - 2\varepsilon} =  H_P^{-\frac{1}{2} + \varepsilon, q-1-2\varepsilon}; \\
        = H_P^{-\frac{1}{2}+\varepsilon, q-1} \inj H_P^{-\frac{1}{2}, q-1}.
    \end{cases}
\end{align*} Moreover, by the last statement of \cite[Prop. D.7]{radethesis}, we have the estimate
\begin{align*}
    \|Q_1(a, \Omega)\|_{W_P(\tau_0)} &\leq C \tau_0(\|F_{\text{ref}}\|_{H^{\frac{1}{2} + \varepsilon, q+\frac{n}{2}}} + \|Rm\|_{H^{\frac{1}{2} + \varepsilon, q+\frac{n}{2}}})\|(0, \Omega)\|_{U(\tau_0)} \\
    &\leq C \tau_0^{1 - \varepsilon}\|(a, \Omega)\|_{U(\tau_0)}.
\end{align*} \par
Next we consider \(Q_2\). By interpolation, \(\Omega \in H^{-\frac{1}{2}+4\varepsilon, q+1-8\varepsilon}\). Since \(8\varepsilon < \min\{1, \mu\}\), we have
the Sobolev multiplication map \(H^{\frac{1}{2}, q} \times H^{-\frac{1}{2} + 4\varepsilon, q+1-8\varepsilon} \to H^{-\frac{1}{2} + \varepsilon, q}\). Hence
\begin{align*}
    &a \# \Omega \in H^{-\frac{1}{2} + \varepsilon, q}
    \begin{cases}
        \inj H^{-\frac{1}{2} + \varepsilon, q - 2\varepsilon} = H_P^{-\frac{1}{2} + \varepsilon, q-2\varepsilon} ; \\
        = H_P^{-\frac{1}{2} + \varepsilon, q}\inj H_P^{-\frac{1}{2}, q}. 
    \end{cases}
\end{align*} Moreover, we have the estimate 
\begin{align*}
    \|(a \# \Omega, 0)\|_{W_P(\tau_0)} \leq C \tau_0^{3\varepsilon}\|(a, \Omega)\|_{U(\tau_0)}^2.
\end{align*} Furthermore, we have \(\nabla_{\text{ref}} a \in H^{\frac{1}{2}, q-1}\) and \(\nabla_{\text{ref}} \Omega \in H^{-\frac{1}{2} + 4\varepsilon, q-8\varepsilon}\). By our assumptions on \(q\) and \(\varepsilon\),
we have the Sobolev multiplication maps \(H^{\frac{1}{2}, q-1} \times H^{-\frac{1}{2} + 4\varepsilon, q-8\varepsilon} \to H^{-\frac{1}{2} +\varepsilon, q-1}\) and \(H^{\frac{1}{2}, q} \times H^{-\frac{1}{2} + 4\varepsilon, q+1-8\varepsilon} \to H^{-\frac{1}{2} + \varepsilon, q-1}\). Hence
\begin{align*}
    \nabla_{\text{ref}} a \# \Omega, \ a \# \nabla_{\text{ref}} \Omega\in H^{-\frac{1}{2} + \varepsilon, q-1}
    \begin{cases}
        \inj H^{-\frac{1}{2} + \varepsilon, q-1-2\varepsilon} = H_P^{-\frac{1}{2} + \varepsilon, q-1-2\varepsilon}; \\
        = H_P^{-\frac{1}{2} + \varepsilon, q-1} \inj H_P^{-\frac{1}{2}, q-1}.
    \end{cases}
\end{align*} Moreover, we have the bound
\begin{align*}
    \|(0, \nabla_{\text{ref}} a \# \Omega + a \# \nabla_{\text{ref}} \Omega)\|_{W_P(\tau_0)} \leq C\tau_0^{3\varepsilon} \|(a, \Omega)\|_{U(\tau_0)}^2.
\end{align*} Overall, we have
\begin{align*}
    \|Q_2\| \leq C\tau_0^{3\varepsilon}.
\end{align*} \par 
Finally, we consider \(Q_3\). We have \(a \in H^{\frac{1}{2} + \varepsilon, q-2\varepsilon}\), and by interpolation \(\Omega \in H^{-\frac{1}{2} + 3\varepsilon, q + 1 -6\varepsilon}\). By our assumptions on \(q\) and \(\varepsilon\),
 we have the Sobolev multiplication maps \(H^{\frac{1}{2} + \varepsilon, q-2\varepsilon} \times H^{-\frac{1}{2} + 3\varepsilon, q+1-6\varepsilon} \to H^{-\frac{1}{2} + \varepsilon, q-2\varepsilon}\) and \(H^{\frac{1}{2} + \varepsilon, q - 2\varepsilon} \times H^{-\frac{1}{2} + \varepsilon, q-2\varepsilon} \to H^{-\frac{1}{2} + \varepsilon, q - 1}\). Hence
\begin{align*}
    &a \# a \# \Omega \in H^{-\frac{1}{2} + \varepsilon, q-1}
    \begin{cases}
        \inj H^{-\frac{1}{2} + \varepsilon, q-1-2\varepsilon} = H_P^{-\frac{1}{2} + \varepsilon, q - 1 -2\varepsilon} ; \\
        = H_P^{-\frac{1}{2} + \varepsilon, q-1} \inj H_P^{-\frac{1}{2}, q-1}.
    \end{cases}
\end{align*}Moreover, we have the bound 
\begin{align*}
    \|Q_3\| \leq C \tau_0^{4\varepsilon}.
\end{align*}\par 
It now follows from (\ref{shorttimeexistence:estimatesforexistence}) that the operators \(A_0, \dots, A_3\) from \cite[pg. 19]{radethesis}, which are given by applying the $Q_i$ to $(a_1, \Omega_1)$ and $(a_2, \Omega_2)$ such that $A_j$ is $j$-homogeneous in $(a_2, \Omega_2),$ satisfy
\begin{align*}
    \|A_0\| &\leq C(\tau_0^{1-\varepsilon}\tau_0^{-\varepsilon} + \tau_0^{3\varepsilon}\tau_0^{-2\varepsilon} + \tau_0^{4\varepsilon}\tau_0^{-3\varepsilon})(1 + \|(a_0, \Omega_0)\|_X)^3 \\
    &\leq C\tau_0^\varepsilon(1 + \|(a_0, \Omega_0)\|_X)^3 \\
    \|A_1\|  &\leq C(\tau_0^{1-\varepsilon} + \tau_0^{3\varepsilon}\tau_0^{-\varepsilon} + \tau_0^{4\varepsilon}\tau_0^{-2\varepsilon})(1 + \|(a_0, \Omega_0)\|_X)^2 \\
    &\leq C\tau_0^{2\varepsilon}(1 + \|(a_0, \Omega_0)\|_X)^2 \\
    \|A_2\| &\leq C(\tau_0^{3\varepsilon} + \tau_0^{4\varepsilon}\tau_0^{-\varepsilon})(1 + \|(a_0, \Omega_0)\|_X) \\
    &\leq C \tau_0^{3\varepsilon}(1 + \|(a_0, \Omega_0)\|_X) \\
    \|A_3\| &\leq C\tau_0^{4\varepsilon}.
\end{align*} Hence the left-hand side of \cite[(24), pg. 18]{radethesis} is bounded above by
\begin{align*}
     C\tau_0^{\varepsilon}(m^{-1}(1 + \|(a_0, \Omega_0)\|_X)^3 + (1+ \|(a_0, \Omega_0)\|_X)^2 + m(1 + \|(a_0, \Omega_0)\|_X) + m^2).
\end{align*} Therefore, if we take \(m\) equal to say \(\tau_0^{-\frac{\varepsilon}{4}}\) and
\begin{align*}
    \tau_0 < \frac{(C(1 + \|(a_0, \Omega_0)\|_X)^3)^{-\frac{2}{\varepsilon}}}{2\|L^{-1}\|},
\end{align*} then by \cite[Lemma B.5]{radethesis}, the inhomogeneous system (\ref{inhomogeneoussystem}) has a solution \((a_2, \Omega_2) \in U_P(\tau_0)\). Moreover, we have the bound \(\|(a_2, \Omega_2)\|_{U_P(\tau_0)} \leq m\), and this solution is unique in \(B_m(U_P(\tau_0))\). Since \(m \to \infty\) as \(\tau_0 \searrow 0\), the same argument as on \cite[pg. 21]{radethesis} yields uniqueness of solutions to the system (\ref{shorttimeexistence:yangmillssystem}). \par We next follow the argument on \cite[pg. 20-21]{radethesis} to show that
\begin{align*}
    a &\in C^0([0, \tau_0], H^q) \\
    \Omega &\in C^0([0, \tau_0], H^{q-1}) \cap H^{0, q}.
\end{align*} By \cite[Prop. D.9]{radethesis}, \(\Omega_1 \in C^0([0, \tau_0], H^{q-1}) \cap H^{0, q}\). Consequently, it follows from integrating (\ref{shorttimeexistence:ivp}) in time that
\begin{align*}
    (1 + \nabla_{\text{ref}}^\ast \nabla_{\text{ref}})\int_0^t \Omega_1(s) \, ds = \int_0^t \Omega_1(s) \, ds + \Omega_0 - \Omega_1 \in C^0([0, \tau_0], H^{q-1}).
\end{align*} Thus by \cite[Prop. D.1]{radethesis}, \(\int_0^t \Omega_1(s) \, ds \in C^0([0, \tau_0], H^{q+1})\), whence \(a_1 \in C^0([0, \tau_0], H^q)\). Next we consider the regularity of \((a_2, \Omega_2)\). Since \((a, \Omega) \in U(\tau_0)\), the Sobolev multiplication maps obtained above when estimating the norms of \(Q_1, Q_2, Q_3\) imply that
\begin{align*}
    &a \# \Omega \in H^{-\frac{1}{2} + \varepsilon, q} \\
    & (F_{\text{ref}} + Rm)\# \Omega, (\nabla_{\text{ref}} a) \# \Omega, a \# \nabla_{\text{ref}} \Omega, a \# a \# \Omega \in H^{-\frac{1}{2} + \varepsilon, q-1}. 
\end{align*} Hence
\begin{align*}
    &\frac{d}{dt} a_2 + D_{\text{ref}}^\ast \Omega_2 \in H^{-\frac{1}{2} + \varepsilon, q} \\
    &\frac{d}{dt} \Omega_2 + \nabla_{\text{ref}}^\ast \nabla_{\text{ref}} \Omega_2 \in H^{-\frac{1}{2} + \varepsilon, q-1}.
\end{align*} Thus by \cite[Prop. D.8]{radethesis}, \(\Omega_2 \in  H^{\frac{1}{2} + \varepsilon, q-1} \cap  H^{-\frac{1}{2} + \varepsilon, q+1}\). By Sobolev embedding \cite[(74)]{radethesis}, \(\Omega_2 \in C^0([0, \tau_0], H^{q-1}) \supset  H^{\frac{1}{2} + \varepsilon, q-1}\). By interpolation \(\Omega_2 \in H^{\varepsilon, q} \subset H^{0, q}\). Furthermore, \(D_{\text{ref}}^\ast \Omega_2 \in H^{-\frac{1}{2} + \varepsilon, q}\), so that by integrating in time and Sobolev embedding, \(a_2 \in C^0([0, \tau_0], H^q)\). Thus we deduce that \((a, \Omega) = (a_1, \Omega_1) + (a_2, \Omega_2)\) has the claimed regularity properties. By bootstrapping, \((a, \Omega)\) is smooth on \(M \times [0, \tau_0]\) if \((a_0, \Omega_0)\) is smooth. \par
Finally, it follows from the implicit function theorem \cite[Proof of Lemma B.5, pg. 19]{radethesis} that given \(K > 0\), there exists a uniform \(\tau_0 > 0\) such that for all \((a_0, \Omega_0) \in B_K(0) \subset X\), the maps
\begin{align}
    &Y_1 : B_K(0) \to U(\tau_0) \\
    &Y_2 : B_K(0) \to C^0([0, \tau_0], X) \label{shorttimeexistence:smoothdependence}
\end{align} sending initial data to the corresponding solution of (\ref{shorttimeexistence:yangmillssystem}) are smooth. \par 
We still need to show that if \(\Omega_0 = F_{A_0}\), then \(\Omega(t) = F_{A(t)}\). We just need to verify that if \(q > \frac{n}{2}-1\), then the curvature of an \(H^q\) connection is in \(H^{q-1}\), so since we also have \(q \geq 1\), the argument in section B.3 of \cite{radethesis} holds. We expand
\begin{align*}
    F_{A_0} = F_{\text{ref}} + D_{\text{ref}}a_0 + a_0 \wedge a_0,
\end{align*} so since \(q > \frac{n}{2}-1\), the \(a_0 \wedge a_0\) term is in \(H^{q-1}\), and thus the claim follows.
 \end{proof}

 

\vspace{5mm}



\subsection{Proof of item (\ref{wellposedness:wellposedness}) of Theorem \ref{thm:wellposedness}} \label{item1proof}
 Since \(T_0 < \min\{T, \infty\}\), we may set
   \begin{align*}
       K := \sup_{0 \leq t \leq T_0} \|A(t)\|_{H^q}+1 < \infty.
   \end{align*} By Theorem \ref{thm:shorttimeexistence}, there exists \(\tau > 0\) such that for any connection \(B_0 \in H^q\) with \(\|B_0\|_{H^q} \leq K\), the solution of (\ref{ymf}) starting at \(B_0\) exists on \([0, \tau]\). By shrinking \(\tau\), we may WLOG assume \(\frac{T_0}{\tau}\) is an integer \(N \geq 1\). \par 
   Let \(\varepsilon > 0\). By the short-time well-posedness (\ref{shorttimeexistence:smoothdependence}),  we have the following. For each integer \(0 \leq i \leq N-1\), there exists an \(H^q\) neighborhood \(U_i\) of \(A(i\tau)\) such that for any \(B_0 \in U_i\), the solution \(B(t)\) of (\ref{ymf}) with \(B(0) = B_0\) exists on \([0, \tau]\) and satisfies 
   \begin{align*}
       \sup_{i\tau \leq t \leq (i+1)\tau} \|B(t-i\tau) - A(t)\|_{H^q} < \varepsilon.
   \end{align*} If \(N = 1\), then this is precisely item (\ref{wellposedness:wellposedness}) of Theorem \ref{thm:wellposedness} and we are done. If \(N \geq 2\), we have by (\ref{shorttimeexistence:smoothdependence}) again that we may choose a neighborhood \(V_{N-2} \subset U_{N-2}\) of \(A((N-2)\tau)\) such that for 
   \(B_0 \in V_{N-2}\), \(B(\tau) \subset V_{N-1} := U_{N-1}\). By induction and (\ref{shorttimeexistence:smoothdependence}), we obtain, for \(0 \leq i \leq N-2\), neighborhoods \(V_i \subset U_i\) of \(A(i\tau)\) such that for \(B_0 \in V_i\), \(B(\tau) \in V_{i+1}\). \par
   Let \(\delta > 0\) be small enough so that \(B_\delta(A_0) \subset V_0\). We claim that this choice of \(\delta\) yields item (\ref{wellposedness:wellposedness}). To see this, let \(A_0' \in B_{\delta}(A_0)\). Then the solution of (\ref{ymf}) \(A'(t)\) starting at \(A_0'\) exists on \([0, \tau]\), \(\sup_{0 \leq t \leq \tau} \|A'(t) - A(t)\|_{H^q} < \varepsilon\), and \(A'(\tau) \subset V_1\). Then by definition of \(V_i, U_i\), we have by induction that for \(0 \leq i \leq N\), \(A'(t)\) exists on \([0, i\tau]\), \(A'(i\tau) \in V_i\), and \(\sup_{0 \leq t \leq i\tau} ||A'(t) - A(t)\|_{H^q} < \varepsilon\). Taking \(i = N\), and recalling that \(N\tau = T_0\), we deduce item (\ref{wellposedness:wellposedness}) of Theorem \ref{thm:wellposedness} in general, as desired.
%\end{proof}




%\subsection{Proofs of items (\ref{wellposedness:gaugeequivalenttosmooth}-\ref{wellposedness:blowupcharacterization}) of Theorem \ref{thm:wellposedness} }


\vspace{5mm}




\subsection{Sobolev-distance estimate under (\ref{ymf})} \label{sobolevdistest}

To prove items (\ref{wellposedness:gaugeequivalenttosmooth}) and (\ref{wellposedness:blowupcharacterization}) of Theorem \ref{thm:wellposedness}, we will use the following adaptation of \cite[Prop. 3.3]{instantons} to higher dimensions. We note that this result does not require compactness of $M$. 
\begin{thm}\label{thm:Hkestimates}
    Let \(k\) be an integer such that
    \begin{align*}
        k \geq 
        \begin{cases}
            0 & n \leq 4 \\
            \frac{n}{2} - 1 & n \geq 5.
        \end{cases}
    \end{align*}
   
    Suppose that \(A(t)\) is an \(H^q\) solution of (\ref{ymf}), for \(q > k\), on $\LB 0, T \right)$ such that for times \(t_1 < t_2 \in [0, T)\) and constants positive constants \(\tau, K\) we have
    \begin{align}\label{Hkestimates:curvaturebound}
        \sup_{t \in [t_1 - \tau, t_2]} \|F(t)\|_\infty \leq K.
    \end{align} Set
    \begin{align*}
        \rho := \int_{t_1 - \tau}^{t_2} \|D^\ast F(t)\| \, dt .
    \end{align*} 
    Given $\Omega \Subset M,$ there exists a constant \(C_{(\ref{Hkestimates:Hkbounds})} \geq 1\), depending on \(k, \tau, K\), and the geometry of $M$ near $\Omega,$ such that
    \begin{align}\label{Hkestimates:Hkbounds}
        \|A(t_2) - A(t_1)\|_{H^{k+1}(\Omega)} \leq C_{(\ref{Hkestimates:Hkbounds})}\left(1+\rho^{(k+2)!} \right) \left(1 + \|A(t_i)\|_{H^k}^{(k+1)!} \right) \rho, \quad i = 1, 2.
    \end{align}
\end{thm}
\begin{proof}
By item (\ref{wellposedness:wellposedness}) of Theorem \ref{thm:wellposedness}, it suffices to prove (\ref{Hkestimates:Hkbounds}) when \(A(t)\) is smooth. The argument of \cite[Prop. 3.3]{instantons} works for \(n \leq 4\). In general dimension, we do the following. Let \(r \geq 0\) be an integer, and set
\begin{align*}
    f_r(t) := \max_{0 \leq i \leq r} \|\nabla^{(i)} D^\ast F(t)\|_{L^\infty(\Omega_1)}.
\end{align*} By the curvature bound (\ref{Hkestimates:curvaturebound}), one can easily pass from an $L^1$ bound on $\| D^*F(t) \|$ to supremum bounds on all higher derivatives in the manner of \cite[Prop. 3.2]{instantons} or \cite[Lemma 2.2]{waldronuhlenbeck}, to obtain
\begin{align}\label{Hkestimates:rhobound}
    \int_{s_1}^{s_2} f_r(t) \, dt \leq C_{r, \tau, K, M} \rho, \quad t_1 \leq s_1 < s_2 \leq t_2.
\end{align}
We now use (\ref{Hkestimates:rhobound}) to estimate $A(t)$ in the Sobolev distances. Since \(\Omega\) is compactly contained in \(M\), there exists 
\begin{comment}
    by Sard's theorem applied to a cutoff for \Omega_1
\end{comment} 
an open set \(\Omega_1\) such that \(\Omega\) has \(C^1\) boundary and \(\Omega \Subset \Omega_1 \Subset M\). In the sequel, all Lebesgue/Sobolev spaces will be on \(\Omega_1\) unless otherwise noted. Let
     \begin{align*}
         q_r &:=
         \begin{cases}
             \frac{n}{r+1} & 0 \leq r \leq \frac{n}{2}-1 \\
             2 & r > \frac{n}{2}-1.
         \end{cases} \\
        \tilde{r} &:= \max\left\{r, \frac{n}{2}\right\}.
     \end{align*}
For \(0 \leq r \leq \frac{n}{2}-1\), define the norm
    \begin{align*}
        \|B\|_{X_r} := \max_{0 \leq i \leq r} \|\nabla_{\text{ref}}^{(i)} B\|_{q_i}.
    \end{align*}
    We will first establish the estimate
    \begin{align}\label{Hkestimates:lowinductionestimate}
        \|\nabla_{\text{ref}}^{(r)}(A(s_2) - A(s_1))\|_{q_r} \leq C_{r, \Omega_1} \left( 1 + \sup_{s_1 \leq s \leq s_2} \|A(s)\|_{X_{r-1}}^r \right) \int_{s_1}^{s_2} f_r(s) \, ds, \ \ t_1 \leq s_1 < s_2 \leq t_2,
    \end{align}
    for \(0 \leq r \leq \frac{n}{2}-1\), and the estimate
    \begin{align}\label{Hkestimates:highinductionestimate}
        \|\nabla_{\text{ref}}^{(r)}(A(s_2) - A(s_1))\|_2 \leq C_{r, \Omega_1} \left(1 + \sup_{s_1 \leq s \leq s_2} \|A(s)\|_{H^{r-1}}^r \right) \int_{s_1}^{s_2} f_r(s) \, ds, \ \  t_1 \leq s_1 < s_2 \leq t_2,
    \end{align} for \(r > \frac{n}{2}-1\).
    
   We compute
    \begin{align*}
        \nabla_{\text{ref}}^{(r)}(A(s_2) - A(s_1)) = \int_{s_1}^{s_2} \frac{d}{dt} \nabla_{\text{ref}}^{(r)} A(t) \, dt &= -\int_{s_1}^{s_2} \nabla_{\text{ref}}^{(r)} D^\ast F(t)  \,dt.
    \end{align*}  
    \begin{comment}
   For notational convenience, we let \(\#\) denote  multilinear (not necessarily commutative or associative) operations depending only on \(r\) whose precise forms and non-commutativity/non-associativity will not be germane.
    \end{comment}
    For all \(r \geq 0\), we have
    \begin{align}\label{Hkestimates:expression}
        \nabla_{\text{ref}}^{(r)} D^\ast F = \sum \nabla_{\text{ref}}^{(\ell_1)} A \# \cdots \# \nabla_{\text{ref}}^{(\ell_j)} A \# \nabla_A^{(m)} D^\ast F,
    \end{align} where \(j\) and the indices \(m, \ell_1, \dots, \ell_j\) are nonnegative and satisfy 
    \begin{align}\label{Hkestimates:indexsum}
        \sum_{i=1}^j \ell_i + j + m\leq r.
    \end{align}
The \(k = 0\) case of (\ref{Hkestimates:expression}) is clear, and we have
    \begin{align*}
        \nabla_{\text{ref}}\sum \nabla_{\text{ref}}^{(\ell_1)} A \# \cdots \# \nabla_{\text{ref}}^{(\ell_j)} A \# \nabla_A^{(m)} D^\ast F = &\sum \nabla_{\text{ref}} \left( \nabla_{\text{ref}}^{(\ell_1)} A \# \cdots \# \nabla_{\text{ref}}^{(\ell_j)} A \right) \# \nabla_A^{(m)} D^\ast F \\
        &+ \sum \nabla_{\text{ref}}^{(\ell_1)} A \# \cdots \# \nabla_{\text{ref}}^{(\ell_j)} A \# \nabla_{\text{ref}} \nabla_A^{(m)} D^\ast F.
    \end{align*} Using the Leibniz rule and the identity \(\nabla_{\text{ref}} = \nabla_A + A \#\) on the second term, we can obtain (\ref{Hkestimates:expression}) by induction.
    
Next, we will use the Sobolev and H\"older's inequalities to prove (\ref{Hkestimates:lowinductionestimate}) and (\ref{Hkestimates:highinductionestimate}).
      Now, by (\ref{Hkestimates:expression}) and H{\"o}lder's inequality, we have
     \begin{align*}
         \|\nabla_{\text{ref}}^{(r)} D^\ast F(s) \|_{q_r} &\leq \sum \|\nabla_{\text{ref}}^{(\ell_1)} A \# \cdots \# \nabla_{\text{ref}}^{(\ell_j)} A \# \nabla_A^{(m)} D^\ast F\|_{q_r} \\
        &\leq C_{\Omega_1}\left(1 + \sum_{j \geq 1} \|\nabla_{\text{ref}}^{(\ell_1)} A \# \cdots \# \nabla_{\text{ref}}^{(\ell_j)} A\|_{q_r} \right)f(s),
     \end{align*}
     where the indices are still subject to (\ref{Hkestimates:indexsum}).
     We will now estimate the individual terms $\|\nabla_{\text{ref}}^{(\ell_1)} A \# \cdots \# \nabla_{\text{ref}}^{(\ell_j)} A\|_{q_r}$ appearing in the sum for $j \geq 1.$
     
     \begin{claim}
     For each \(r \geq 0\), we have
    \begin{align}\label{Hkestimates:sumclaim}
        \sum_{\ell_i > r-1- \frac{n}{2}} \frac{n-2(\tilde{r}-1 - \ell_i)}{2n}
        \begin{cases}
            \leq \frac{1}{q_r}, & 0 \leq r \leq \frac{n}{2}, \text{ or } r > \frac{n}{2} \text{ and } j = 1, \\
            < \frac{1}{q_r}, & r > \frac{n}{2} \text{ and } j \geq 2.
        \end{cases}
    \end{align}
    \end{claim}
    \begin{claimproof}
    \emph{Case 1.} \(r \leq \frac{n}{2}\). By definition of \(\tilde{r}\), We have
    \begin{align*}
        \sum_{\ell_i > r - 1 - \frac{n}{2}} \frac{n-2(\tilde{r}-1 - \ell_i)}{2n} & = \sum_{i = 1}^j \frac{\ell_i+1}{n}.
    \end{align*}
    In view of (\ref{Hkestimates:indexsum}), we have 
        \begin{align*}
        \sum_{i = 1}^j \frac{\ell_i+1}{n} \leq \frac{r}{n} 
        \begin{cases}
            \leq \frac{r+1}{n} = \frac{1}{q_r} & r \leq \frac{n}{2}- 1, \\
            \leq \frac{1}{2} = q_r & \frac{n}{2}-1 < r \leq \frac{n}{2}.
        \end{cases}
    \end{align*}
        \emph{Case 2.} \(r > \frac{n}{2} \text{ and } j = 1\). Then \(\ell_j = r-1\), so
    \begin{align*}
        \sum_{\ell_i > r - 1 - \frac{n}{2}} \frac{n-2(\tilde{r}-1-\ell_i)}{2n} = \frac{n-2(r-1-(r-1))}{2n} = \frac{1}{2} = \frac{1}{q_r}.
    \end{align*}
    \emph{Case 3.} Finally suppose that \(r > \frac{n}{2} \text{ and } j \geq 2\). Let \(J := |\{\ell_i \mid \ell_i >  r - 1 - \frac{n}{2}\}|\). It follows from (\ref{Hkestimates:indexsum}) that either \(J \geq 2\), or \(J = 1\) and \(\ell_1 \leq r-2\). In the first case
    \begin{align*}
         \sum_{\ell_i  > r - 1 - \frac{n}{2}} \frac{n-2(\tilde{r}-1-\ell_i)}{2n} &= J\frac{n-2r}{2n} + \sum_{\ell_i  > r - 1 - \frac{n}{2}} \frac{\ell_i+1}{n}  \\
         &\leq J\frac{n-2r}{2n} + \frac{r}{n} \\
         &= \frac{1}{2} + (J-1)\frac{n-2r}{2n} \\
         &< \frac{1}{2} = \frac{1}{q_r},
    \end{align*} and in the second case
    \begin{align*}
        \frac{n-2(r-1-\ell_1)}{2n} \leq \frac{n-2(r-1-(r-2))}{2n} < \frac{1}{2} = \frac{1}{q_r}.
    \end{align*}
    \end{claimproof}
    We can now prove (\ref{Hkestimates:lowinductionestimate}) and (\ref{Hkestimates:highinductionestimate}). From H{\"o}lder's inequality, if \(1 \leq p_1, \cdots, p_s, q \leq \infty\) are such that \(\sum_{i = 1}^s \frac{1}{p_i} \leq \frac{1}{q}\), then there exists a constant \(C\), depending on \(\text{Vol}(\Omega_1), p_1, \cdots, p_s, q\), such that for functions \(f_1, \cdots, f_s\),
    \begin{align}\label{Hkestimates:Holdersinequality}
        \left\|\prod_{i = 1}^s f_i\right\|_q \leq C\prod_{i = 1}^s\|f_i\|_{p_i}.
    \end{align} First note that if \(\ell_i > r - 1 - \frac{n}{2}\), then
    \begin{align}\label{Hkestimates:validLebesgueexponent}
        0 < \frac{n-2(\tilde{r}-1 - \ell_i)}{2n} \leq 1.
    \end{align} The strict lower bound of \(0\) follows for \(r > \frac{n}{2}\), whence \(\tilde{r} = r\), by the assumption on \(\ell_i\) since \(\frac{n-2(\tilde{r}-1 - \ell_i)}{2n} = \frac{n-2(r-1 - \ell_i)}{2n}\), and follows for \(r \leq \frac{n}{2}\), whence \(\tilde{r} = \frac{n}{2}\), since \(1 + \ell_i > 0\). The upper bound of \(1\) follows from \(\tilde{r} - 1 - \ell_i \geq 0\) by (\ref{Hkestimates:indexsum}). \par
    Now, if \(r \leq \frac{n}{2}\), then \(\ell_i > r - 1 - \frac{n}{2}\) for all \(i\), and moreover \(\ell_i \leq \frac{n}{2}-1\) by (\ref{Hkestimates:indexsum}). Therefore, (\ref{Hkestimates:sumclaim}), (\ref{Hkestimates:Holdersinequality}), and (\ref{Hkestimates:validLebesgueexponent}) yield
    \begin{align*}
        \|\nabla_{\text{ref}}^{(\ell_1)} A \# \cdots \# \nabla_{\text{ref}}^{(\ell_j)} A\|_{q_r} &\leq  C_{r, \Omega_1} \prod_{i=1}^j \|\nabla_{\text{ref}}^{(\ell_i)} A\|_{\frac{2n}{n-2(\tilde{r}-1-\ell_i)}} = C_{r,  \Omega_1} \prod_{i=1}^j \|\nabla_{\text{ref}}^{(\ell_i)} A\|_{\frac{n}{\ell_i+1}}.
    \end{align*} The bound (\ref{Hkestimates:lowinductionestimate}) now follows the definition of the \(X_r\) norm and the fact that \(j \leq r\) by (\ref{Hkestimates:indexsum}). If \(r > \frac{n}{2}\) and \(j = 1\), then (\ref{Hkestimates:indexsum}) implies that the term is \(\nabla_{\text{ref}}^{(r-1)} A\), and clearly
    \begin{align}\label{Hkestimates:onlyoneterm}
        \|\nabla_{\text{ref}}^{(r-1)} A\|_2 \leq \|A\|_{H^{r-1}}.
    \end{align} On the other hand, suppose \(r > \frac{n}{2}\) and \(j \geq 2\). Let
    \begin{align*}
        G_r := \max\left\{\ell_1, \cdots, \ell_j \text{ satisfying } (\ref{Hkestimates:indexsum}) \mid \sum_{\ell_i > r - 1 -\frac{n}{2}} \frac{n-2(r-1-\ell_i)}{2n}\right\}.
    \end{align*} Note that \(\frac{1}{2} - G_r > 0\) since \(G_r\) is the maximum of a finite set of numbers which are strictly less than \(\frac{1}{2}\) by (\ref{Hkestimates:sumclaim}). Moreover, \(G_r \geq 0\) by the left inequality in (\ref{Hkestimates:validLebesgueexponent}). Let
    \begin{align*}
        \frac{1}{p_r} := \frac{1}{r}\left(\frac{1}{2} - G_r\right). 
    \end{align*} Note that \(p_r\) depends only on \(n\) and \(r\), and moreover \(1 \leq p_r < \infty\) since \(0 \leq G_r < \frac{1}{2}\). Thus by (\ref{Hkestimates:indexsum}),
    \begin{align*}
        \sum_{\ell_i \leq r - 1 - \frac{n}{2}} \frac{1}{p_r} + \sum_{\ell_i > r-1-\frac{n}{2}} \frac{n-2(r-1-\ell_i)}{2n} &\leq \frac{r}{p_r} + \sum_{\ell_i > r-1-\frac{n}{2}} \frac{n-2(r-1-\ell_i)}{2n}  \\
        &\leq \frac{1}{2}.
    \end{align*}Consequently, 
    \begin{align*}
        \|\nabla_{\text{ref}}^{(\ell_1)} A \# \cdots \# \nabla_{\text{ref}}^{(\ell_j)} A\|_{q_r} &\leq  C_{r, \Omega_1} \prod_{\ell_i \leq r-1-\frac{n}{2}} \|\nabla_{\text{ref}}^{(\ell_i)} A\|_{p_r} \prod_{\ell_i >r-1-\frac{n}{2}} \|\nabla_{\text{ref}}^{(\ell_i)} A\|_{\frac{2n}{n-2(r-1-\ell_i)}}.
    \end{align*} 
     Since the reference connection is metric-compatible and since \(\Omega_1\) has \(C^1\) boundary, it follows from Kato's inequality and the Sobolev inequality on functions that we have the following bounds for \(B \in 
     \Gamma((\otimes^l T^\ast M) \otimes \mathfrak{g}_E), l \geq 0\):
    \begin{align}\label{Hkestimates:sobolevinequalities}
        \|\nabla_{\text{ref}}^{(\ell)} B\|_{\frac{2n}{n-2(r-1-\ell)}} &\leq C_{\Omega_1} \|B\|_{H^{r-1}}, \quad \ell > r - 1 - \frac{n}{2} \\
        \|\nabla_{\text{ref}}^{\ell} B\|_{p} &\leq C_{p, \Omega_1} \|B\|_{H^{\ell + \lceil\frac{n}{2}\rceil}}, \quad 1 \leq p < \infty. \nonumber
    \end{align} By integrality of \(\ell_i\) and \(r\), \(\ell_i \leq r-1 - \frac{n}{2}\) implies \(\ell_i \leq r-1 - \lceil\frac{n}{2}\rceil\). Hence
    \begin{align*}
         \prod_{\ell_i \leq r-1-\frac{n}{2}} \|\nabla_{\text{ref}}^{(\ell_i)} A\|_{p_r} \prod_{\ell_i >r-1-\frac{n}{2}} \|\nabla_{\text{ref}}^{(\ell_i)} A\|_{\frac{2n}{n-2(r-1-\ell_i)}} \leq C_{r,  \Omega_1}\left(1 + \|A\|_{H^{r-1}}^r\right).
    \end{align*} Taking into account (\ref{Hkestimates:onlyoneterm}), the bound (\ref{Hkestimates:highinductionestimate}) now follows. \par
    We finally prove the desired estimate (\ref{Hkestimates:Hkbounds}), first with \(t_i = t_1\). Any dependence on \(\Omega_1\) is a dependence on \(\Omega\). Then for \(t \in [t_1, t_2]\)
    \begin{align*}
        \sup_{t_1 \leq t \leq t_2} \|A(t)\|_{q_0} &\leq \|A(t_1)\|_{q_0} + \sup_{t_1 \leq t \leq t_2} \|A(t) - A(t_1)\|_{q_0} \\
        &\leq \|A(t_1)\|_{X_0} + \sup_{t_1 \leq t \leq t_2} \|A(t) - A(t_1)\|_{q_0} \\ 
        &\leq \|A(t_1)\|_{X_0} + C_{\tau, K, \Omega} \rho. 
    \end{align*} Suppose that for some \(r\), with \(0 \leq r \leq \frac{n}{2}-2\), we have the estimate
    \begin{align*}
        \sup_{t_1 \leq t \leq t_2} \|A(t)\|_{X_{r}} \leq C_{r, \tau, K, \Omega}\left(1+\rho^{(r+1)!} \right) \left(1+\|A(t_1)\|_{X_{r}}^{r!} \right).
    \end{align*} Then by (\ref{Hkestimates:lowinductionestimate})
    \begin{align*}
        \sup_{t_1 \leq t \leq t_2} \|\nabla_{\text{ref}}^{(r+1)} A(t)\|_{q_{r+1}} &\leq \|\nabla_{\text{ref}}^{(r+1)} A(t_1)\|_{q_{r+1}} + \sup_{t_1 \leq t \leq t_2} \|\nabla_{\text{ref}}^{(r+1)}(A(t) - A(t_1))\|_{q_{r+1}} \\
    &\leq \|A(t_1)\|_{X_{r+1}} + \sup_{t_1 \leq t \leq t_2} \|\nabla_{\text{ref}}^{(r+1)}(A(t) - A(t_1))\|_{q_{r+1}} \\
    &\leq \|A(t_1)\|_{X_{r+1}} + C_{r+1, \tau, K, \Omega} (1 + \sup_{t_1 \leq t \leq t_2} \|A(t)\|_{X_{r}}^{r+1})\rho  \\
    &\leq \|A(t_1)\|_{X_{r+1}} + C_{r+1, \tau, K, \Omega} (1 + ((1+\rho^{(r+1)!})(1+\|A(t_1)\|_{X_{r}}^{r!}))^{r+1})\rho  \\
    &\leq C_{r+1, \tau, K, \Omega}\left(1+\rho^{(r+2)!} \right)\left(1+\|A(t_1)\|_{X_{r+1}}^{(r+1)!} \right).
    \end{align*} Thus by induction
    \begin{align*}
        \sup_{t_1 \leq t \leq t_2} \|A(t)\|_{X_{\lfloor\frac{n}{2}\rfloor-1}} \leq C_{\tau, K, \Omega} \left(1 + \rho^{\lfloor\frac{n}{2}\rfloor!} \right) \left(1 + \|A(t_1)\|_{X_{\lfloor\frac{n}{2}\rfloor-1}}^{(\lfloor\frac{n}{2}\rfloor-1)!} \right).
    \end{align*} Observe that \(q_r \geq 2\) for all \(r\). Hence
    \begin{align*}
        \|B\|_{H^{\lfloor\frac{n}{2}\rfloor-1}} \leq C_\Omega \|B\|_{X_{\lfloor\frac{n}{2}\rfloor-1}}.
    \end{align*} On the other hand, the Sobolev inequalities (\ref{Hkestimates:sobolevinequalities}) imply
    \begin{align*}
        \|B\|_{X_{\lfloor\frac{n}{2}\rfloor-1}} \leq C_\Omega \|B\|_{H^{\lceil\frac{n}{2}\rceil-1}}.
    \end{align*} Hence
    \begin{align*}
        \sup_{t_1 \leq t \leq t_2} \|A(t)\|_{H^{\lfloor\frac{n}{2}\rfloor-1}} \leq C_{\tau, K, \Omega} \left(1+\rho^{\lfloor\frac{n}{2}
    \rfloor!} \right) \left(1 + \|A(t_1)\|_{H^{\lceil\frac{n}{2}\rceil-1}}^{(\lfloor\frac{n}{2}\rfloor-1)!} \right).
    \end{align*} The bound (\ref{Hkestimates:Hkbounds}) now follows by a similar induction as just carried out, using (\ref{Hkestimates:highinductionestimate}) in place of (\ref{Hkestimates:lowinductionestimate}), combined with the obvious bounds \(\|B\|_{W^{\ell, s}(\Omega)} \leq \|B\|_{W^{\ell, s}(\Omega_1)} \leq \|B\|_{W^{\ell, s}(M)}\). The case \(i = 2\) follows by an essentially identical argument, starting from
    \begin{align*}
        \sup_{t_1 \leq t \leq t_2} \|A(t)\|_{q_0} \leq \|A(t_2)\|_{q_0} + \sup_{t_1 \leq t \leq t_2} \|A(t_2) - A(t)\|_{q_0}.
    \end{align*}
\end{proof}
\begin{comment}
\begin{proof}[Proof of item (\ref{wellposedness:gaugeequivalenttosmooth}) of Theorem \ref{thm:wellposedness}]
    We will show the existence of a gauge transformation as claimed in item (\ref{wellposedness:gaugeequivalenttosmooth}) via a compactness argument. \par
    \(C^\infty(\mathcal{A}_E)\) is dense in \(H^q(\mathcal{A}_E)\), so there exists a sequence of smooth connections \(A_i\) converging to \(A_0\) in \(H^q\). By item (\ref{wellposedness:wellposedness}) of Theorem \ref{thm:wellposedness}, we may assume that the solution \(\tilde{A}_i(t)\) of (\ref{ymf}) with \(\tilde{A}_i(0) = A_i\) exists on \([0, \tau]\), and moreover that the \(\tilde{A}_i(t)\) converge in \(C^0([0, \tau], H^q)\) to \(A(t)\). In particular,
    \begin{align*}
        \sup_{i, t \in [0, \tau]} \|\tilde{A}_i(t)\|_{H^q} < \infty,
    \end{align*} and hence
    \begin{align*}
        \sup_{i, t \in [0, \tau]} \|F_i(t)\|_{H^{q-1}} < \infty.
    \end{align*}Since \(q > \frac{n}{2}-1\), \(H^{q-1} \inj L^p\) for some \(p \in (\frac{n}{2}, \infty]\). Given a uniform \(L^p\) bound on the curvature for such \(p\), epsilon-regularity, e.g. \cite[Lemma 2.2]{waldronuhlenbeck}, yields that for each \(\delta > 0\)
    \begin{align}\label{gaugeequivalenttosmoothprop:curvaturebound}
        \sup_{i, t \in [\delta, \tau]} \|F_i(t)\|_\infty < \infty.
    \end{align} Moreover, since \(\tau \leq 1\), the global energy identity and the convergence of \(A_i \to A_0\) yields
    \begin{align}\label{gaugeequivalenttosmoothprop:tensionbound}
        \sup_i \int_0^{\tau} \|D_i^\ast F_i(t) \| \, dt < \infty.
    \end{align}Then by standard strong Uhlenbeck compactness arguments for Yang-Mills flow, e.g. \cite[Theorem 1.3]{waldronuhlenbeck}, we can pass to a subsequence, also labeled by \(i\), such that there exist smooth gauge transformations \(u_i\) for which \(u_i(\tilde{A}_i(\tau))\) converge in \(C^\infty\) to some smooth connection. \par 
    Let \(t_m := \frac{\tau}{m}\). Note that the bounds (\ref{gaugeequivalenttosmoothprop:curvaturebound}) and (\ref{gaugeequivalenttosmoothprop:tensionbound}) are gauge-invariant, and so they also hold for the \(u_i(\tilde{A}_i(t))\). Since \(u_i(\tilde{A}_i(\tau))\) is smoothly convergent, it follows from Prop. \ref{thm:Hkestimates}, with \(t_2 = \tau\) and \(t_i = t_2\) in the statement of the proposition, that for each \(m\) and for each \(k \geq \lceil\frac{n}{2}\rceil\)
    \begin{align}\label{gaugeequivalenttosmoothprop:Hkbounds}
        \sup_{i, t_m \leq t \leq \tau} \|u_i(\tilde{A}_i(t))\|_{H^k} &\leq \sup_{i, t_m \leq t \leq \tau}(\|u_i(\tilde{A}_i(\tau))\|_{H^k} + \|u_i(\tilde{A}_i(t)) - u_i(\tilde{A}_i(\tau))\|_{H^k}) < \infty.
    \end{align} Now, we may exchange time derivatives of \(u_i(\tilde{A}_i)\) with spatial derivatives using the flow equation. Hence, (\ref{gaugeequivalenttosmoothprop:Hkbounds}) and Sobolev embedding yields that for each \(m\) and each \(k\)
    \begin{align*}
        \sup_i \|u_i(\tilde{A}_i)\|_{C^k(M \times [t_m, \tau])} < \infty
    \end{align*}
    Thus for each \(m\), Arzela-Ascoli implies that we may pass to a subsequence, also labeled by \(i\), such that \(u_i(\tilde{A}_i)\) converges in \(C^\infty(M \times [t_m, \tau])\) to a family of connections \(B_m(t)\), which is also a solution of (\ref{ymf}). Note that \(B_m(t) = B_n(t)\) on their common interval of definition by uniqueness of solutions to (\ref{ymf}). Passing to a subsequence, \(u_i(\tilde{A}_i)\) converge in \(C_{\text{loc}}^\infty(M \times (0, \tau])\) to a smooth solution \(B(t)\) of Yang-Mills flow. \par
    Finally, since the \(\tilde{A}_i(\tau)\) themselves converge in \(H^q\), we have as in \cite[\textsection 2.3.17]{donkron} that the \(u_i\) converge to some gauge transformation \(u_\infty\) in \(H^{q+1}\). Then since
    \begin{align*}
        u_i(\tilde{A}_i(t)) - u_\infty(A(t)) = u_i(\tilde{A}_i(t) - A(t)) + (u_i - u_\infty)(A(t)),
    \end{align*} we deduce that \(B(t) = u_\infty(A(t))\) for \(t \in (0, \tau]\). By the uniqueness of solutions to (\ref{ymf}), the same holds for all \(t\) for which \(A(t)\) is defined, and \(B(t)\) remains smooth. Thus \(u_\infty\) satisfies the conclusions of item (\ref{wellposedness:gaugeequivalenttosmooth}) of Theorem \ref{thm:wellposedness}, as desired.
\end{proof}
\end{comment}



\vspace{5mm}



\subsection{Proof of item (\ref{wellposedness:gaugeequivalenttosmooth}) of Theorem \ref{thm:wellposedness}} \label{gaugeeqtosmth}
    We will show the existence of a gauge transformation as claimed in item (\ref{wellposedness:gaugeequivalenttosmooth}) via a compactness argument. \par
   Since \(C^\infty(\mathcal{A}_E)\) is dense in \(H^q(\mathcal{A}_E)\), there exists a sequence of smooth connections \(A_i\) converging to \(A_0\) in \(H^q\). By item (\ref{wellposedness:wellposedness}) of Theorem \ref{thm:wellposedness}, we may assume that the solution \(\tilde{A}_i(t)\) of (\ref{ymf}) with \(\tilde{A}_i(0) = A_i\) exists on \([0, \tau]\), and moreover that the \(\tilde{A}_i(t)\) converge in \(C^0([0, \tau], H^q)\) to \(A(t)\). In particular,
    \begin{align*}
        \sup_{i, t \in [0, \tau]} \|\tilde{A}_i(t)\|_{H^q} < \infty,
    \end{align*} and hence
    \begin{align*}
        \sup_{i, t \in [0, \tau]} \|F_i(t)\|_{H^{q-1}} < \infty.
    \end{align*}Since \(q > \frac{n}{2}-1\), \(H^{q-1} \inj L^p\) for some \(p \in (\frac{n}{2}, \infty]\). Given a uniform \(L^p\) bound on the curvature for such \(p\), $\varepsilon$-regularity, e.g. \cite[Lemma 2.2]{waldronuhlenbeck}, yields that for each \(\delta > 0\)
    \begin{align}\label{gaugeequivalenttosmoothprop:curvaturebound}
        \sup_{i, t \in [\delta, \tau]} \|F_i(t)\|_\infty < \infty.
    \end{align} Moreover, since \(\tau \leq 1\), the global energy identity and the convergence of \(A_i \to A_0\) yields
    \begin{align}\label{gaugeequivalenttosmoothprop:tensionbound}
        \sup_i \int_0^{\tau} \|D_i^\ast F_i(t) \| \, dt < \infty.
    \end{align}Then by standard strong Uhlenbeck compactness arguments for Yang-Mills flow, e.g. \cite[Theorem 1.3]{waldronuhlenbeck}, we can pass to a subsequence, also labeled by \(i\), such that there exist smooth gauge transformations \(u_i\) for which \(u_i(\tilde{A}_i(\tau))\) converge in \(C^\infty\) to some smooth connection. \par 
    Let \(t_m := \frac{\tau}{m}\). Note that the bounds (\ref{gaugeequivalenttosmoothprop:curvaturebound}) and (\ref{gaugeequivalenttosmoothprop:tensionbound}) are gauge-invariant, and so they also hold for the \(u_i(\tilde{A}_i(t))\). Since \(u_i(\tilde{A}_i(\tau))\) is smoothly convergent, it follows from Theorem \ref{thm:Hkestimates}, with \(t_2 = \tau\) and \(t_i = t_2\) in the statement, that for each \(m\) and for each \(k \geq \lceil\frac{n}{2}\rceil\)
    \begin{align}\label{gaugeequivalenttosmoothprop:Hkbounds}
        \sup_{i, t_m \leq t \leq \tau} \|u_i(\tilde{A}_i(t))\|_{H^k} &\leq \sup_{i, t_m \leq t \leq \tau} \left(\|u_i(\tilde{A}_i(\tau))\|_{H^k} + \|u_i(\tilde{A}_i(t)) - u_i(\tilde{A}_i(\tau))\|_{H^k} \right) < \infty.
    \end{align} Now, we may exchange time derivatives of \(u_i(\tilde{A}_i)\) with spatial derivatives using the flow equation. Hence, (\ref{gaugeequivalenttosmoothprop:Hkbounds}) and Sobolev embedding yields that for each \(m\) and each \(k\)
    \begin{align*}
        \sup_i \|u_i(\tilde{A}_i)\|_{C^k(M \times [t_m, \tau])} < \infty.
    \end{align*}
    Thus for each \(m\), Arzela-Ascoli implies that we may pass to a subsequence, also labeled by \(i\), such that \(u_i(\tilde{A}_i)\) converges in \(C^\infty(M \times [t_m, \tau])\) to a family of connections \(B_m(t)\), which is also a solution of (\ref{ymf}). Note that \(B_m(t) = B_n(t)\) on their common interval of definition by uniqueness of solutions to (\ref{ymf}). Passing to a subsequence, \(u_i(\tilde{A}_i)\) converge in \(C_{\text{loc}}^\infty(M \times (0, \tau])\) to a smooth solution \(B(t)\) of Yang-Mills flow. \par
    Finally, since the \(\tilde{A}_i(\tau)\) themselves converge in \(H^q\), we have as in \cite[\textsection 2.3.17]{donkron} that the \(u_i\) converge to some gauge transformation \(u_\infty\) in \(H^{q+1}\). Then since
    \begin{align*}
        u_i(\tilde{A}_i(t)) - u_\infty(A(t)) = u_i(\tilde{A}_i(t) - A(t)) + (u_i - u_\infty)(A(t)),
    \end{align*} we deduce that \(B(t) = u_\infty(A(t))\) for \(t \in (0, \tau]\). By the uniqueness of solutions to (\ref{ymf}), the same holds for all \(t\) for which \(A(t)\) is defined, and \(B(t)\) remains smooth. Thus \(u_\infty\) satisfies the conclusions of item (\ref{wellposedness:gaugeequivalenttosmooth}) of Theorem \ref{thm:wellposedness}, as desired. \qed
%\end{proof}



\vspace{5mm}



\subsection{Proof of item (\ref{wellposedness:blowupcharacterization}) of Theorem \ref{thm:wellposedness}} \label{ltecriterion}
    Suppose for %a 
    contradiction that
    \begin{align*}
        \limsup_{t \nearrow T} \|F(t)\|_\infty < \infty.
    \end{align*} By item (\ref{wellposedness:gaugeequivalenttosmooth}) of Theorem \ref{thm:wellposedness}, we may gauge transform \(A(t)\) so that \(u(A(t))\) is smooth. The pointwise norm of the curvature is gauge-invariant, so we still have \(\limsup_{t \nearrow T}\ \|F_{u(A(t))}\|_\infty < \infty\). It then follows from Theorem \ref{thm:Hkestimates} and finiteness of \(T\) that for each \(k \geq 0\), the \(H^k\) norm of \(u(A(t))\) remains bounded as \(t \nearrow T\). Thus the \(H^q\) norm of \(A(t)\) remains bounded as \(t \nearrow T\), so by Theorem \ref{thm:shorttimeexistence}, \(A(t)\) exists on \([0, T + \varepsilon]\) for some \(\varepsilon > 0\). This contradicts the maximality of \(T\), as desired. \qed
%\end{proof}



%{\color{red} Where is the proof of Item (2) complete?}


\vspace{10mm}


\section{Continuity at infinite time}

Concerning infinite-time behavior, we will prove the following standard result, cf. \cite[Prop. 7.4]{rade}. See also Kozono-Maeda-Naito \cite{kozonomaedaasymptoticstability, kozonomaedanaitostablemanifold} for early work on asymptotic stability of Yang-Mills flow.
\begin{thm}\label{thm:continuityatinfinitetime}
    Let \(q\) be such that \(q \geq 1\) and \(q > \lceil\frac{n}{2}-1\rceil\). %Let $A_\infty \in H^q(\calA_E)$ be a Yang-Mills connection.
    Suppose that \(A_0 \in H^q(\mathcal{A}_E)\) is such that the solution $A(t)$ of (\ref{ymf}) starting at \(A_0\) exists for all time and converges in \(H^q\) to a Yang-Mills connection \(A_\infty =: A(\infty)\) as \(t \to \infty\). There exists \(\delta_\infty > 0 \) (depending on \(A_\infty, q,\) and \(M\)) such that given any $\eps > 0,$ there exists \(\delta > 0\) (depending on \(A_0, q,\) and \(M\)) as follows.
    
    Let \(A'_0 \in B_{\delta}(A_0)\) and consider the solution $A'(t)$ of (\ref{ymf}) with $A'(0) = A_0'.$ Then exactly one of the following alternatives holds:
    \begin{enumerate}
        \item There exists \(0 < \tau < \infty\) such that $A'(t)$ is defined on $\LB 0 , \tau \RB$ and \(\mathcal{YM}(A'(\tau)) \leq \mathcal{YM}(A_\infty) - \delta_\infty\), or
        \item The solution \(A'(t)\) exists for all time and converges in \(H^q\) to a Yang-Mills connection $A'(\infty)$ as \(t \to \infty\), with \begin{align}\label{continuityatinfinitetime:continuityestimate}
            \| A'(t) - A(t) \|_{H^q} < \eps
        \end{align}
        for all $0 \leq t \leq \infty.$ % Moreover, given \(\varepsilon > 0\)  there exists \(\delta, T > 0\) such that for all \(A_1\) satisfying this convergence property, \(\|A_1 - A_0\|_{H^q} < \delta\) implies
        %\begin{align*}
        %    \|A(t) - A\|_{H^q} < \varepsilon, \ \ \ t \geq T.
        %\end{align*}
    \end{enumerate}
\end{thm}

In \textsection \ref{lojasiewiczineq}, we first describe the setup of an $H^{-1}$ \L ojasiewicz inequality in general dimension, which we then prove in Lemma \ref{lemma:lojasiewicz}. We also specialize to the case of flat connections in Lemma \ref{lemma:uniformlojasiewicz}. Theorem \ref{thm:continuityatinfinitetime} is proven in \textsection \ref{continuityatinftime} using Lemma \ref{lemma:lojasiewicz}.

\vspace{5mm}




\subsection{\L ojasiewicz inequalities} \label{lojasiewiczineq} We will use a variant of the \L ojasiewicz-Simon inequality proved by R\aa de in dimension less than four and by B. Yang \cite[Lemma 12]{byanguniqueness} in general dimension. First we describe our setup. In the setup, \(C\) denotes a constant, depending only on \(M\), which may increase from line to line. \par
Given an \(H^q\) connection \(A\), with \(q > \frac{n}{2}-1\) and \(q \geq 1\), we define the following quantities, cf. \cite[(2.10) in Def. 2.3]{taubespathconnected}. For \(b \in H^1\left(T^\ast M \otimes \mathfrak{g}_E\right)\), set
\begin{align*}
    \|b\|_{H_A^1} &:= \sqrt{\left\|\nabla_A b\right\|^2 + \|b\|^2} \\
    \mathcal{V}(A) &:= \sup_{\substack{b \in H^1\left(T^\ast M \otimes \mathfrak{g}_E\right), \\ b \neq 0}} \frac{\left|\int \langle F_A, D_A b\rangle\right|}{\|b\|_{H_A^1}}.
\end{align*} Note that \(\|b\|_{H_A^1}\) and the expression in the above supremum are finite since, by Sobolev multiplication, \(D_A\) and \(\nabla_A\) map \(H^1\) to \(L^2\). Moreover, the supremum itself is finite by Cauchy-Schwarz and the pointwise domination \(|D_A b| \lesssim |\nabla_A b|\). \par
Moreover, we claim that \(\mathcal{V}\) is continuous with respect to the \(H^q\) topology.
\begin{claimproof}
    Let \(A_1, A_2 \in H^q(\mathcal{A}_E)\). We first establish equivalence of \(\|\cdot\|_{H_{A_1}^1}\) and \(\| \cdot \|_{H_{A_2}^1}\) if \(A_1, A_2\) are close in \(H^q\). Note that
\begin{align*}
   | \|\nabla_{A_1} b\|^2 - \|\nabla_{A_2} b\|^2| &\leq (\|\nabla_{A_1} b\| + \|\nabla_{A_2} b\|)\|(A_1-A_2) \# b\| \\
        &\leq C\left(\|b\|_{H_{A_1}^1} + \|b\|_{H_{A_2}^1}\right)\|A_1-A_2\|_n \|b\|_{\frac{2n}{n-2}}.
\end{align*} By Sobolev embedding and the Kato inequality trick, see e.g. \cite[(1.16)]{instantons},
\begin{align*}
    \|b\|_{\frac{2n}{n-2}} \leq C\min\left\{\|b\|_{H_{A_1}^1}, \|b\|_{H_{A_2}^1}\right\}.
\end{align*} Hence for \(\|A_1 - A_2\|_{H^q}\) sufficiently small such that, by Sobolev embedding, \(C^2\|A_1 - A_2\|_n \leq \varepsilon < 1\),
\begin{align}\label{lojasiewiczclaim:normequivalence}
    \frac{1-\varepsilon}{1+\varepsilon}\|b\|_{H_{A_2}^1}^2 \leq \|b\|_{H_{A_1}^1}^2 \leq \frac{1+\varepsilon}{1-\varepsilon}\|b\|_{H_{A_2}^1}^2.
\end{align} Thus we have equivalence of the two norms. \par
 We next show \begin{align}\label{lojasiewiczclaim:D*Fequivalence}
    \left|\int \langle F_{A_1}, D_{A_1} b\rangle - \int \langle F_{A_2}, D_{A_2} b\rangle\right| \leq \varepsilon\min\left\{\|b\|_{H_{A_1}^1}, \|b\|_{H_{A_2}^1}\right\}.
\end{align} We compute
\begin{align*}
    \int \langle F_{A_1}, D_{A_1} b\rangle - \int \langle F_{A_2}, D_{A_2} b\rangle & = \text{I} + \text{II} + \text{III} + \text{IV} + \text{V},
\end{align*}where  
\begin{align*}
    \text{I} &= \int \langle F_{A_2}, (A_1 - A_2)\wedge b\rangle \\
    \text{II} &= \int \langle D_{A_2}(A_1 - A_2), D_{A_2} b\rangle \\
    \text{III} &= \int \langle D_{A_2}(A_1-A_2), (A_1-A_2) \wedge b\rangle \\
    \text{IV} &= \int \langle (A_1 - A_2) \wedge (A_1-A_2), D_{A_2} b\rangle \\
    \text{V} &= \int \langle (A_1 - A_2) \wedge (A_1-A_2), (A_1 - A_2) \wedge b\rangle.
\end{align*} We estimate
\begin{align*}
    |\text{I}| &\leq C\|F_{A_2}\|\|A_1 - A_2\|_n\|b\|_{\frac{2n}{n-2}} \\
    |\text{II}| &\leq C(\|\nabla_{\text{ref}}(A_1 - A_2)\| + \|A_2\|_4\|A_1 - A_2\|_4)\|b\|_{H_{A_2}^1} \\
    |\text{III}| &\leq C(\|\nabla_{\text{ref}}(A_1 - A_2)\| + \|A_2\|_4\|A_1 - A_2\|_4)\|A_1 -A_2\|_n\|b\|_{\frac{2n}{n-2}} \\
    |\text{IV}| &\leq C\|A_1 - A_2\|_4^2\|b\|_{H_{A_2}^1} \\
    |\text{V}| &\leq C\|A_1 - A_2\|_4^2\|A_1 - A_2\|_n\|b\|_{\frac{2n}{n-2}}.
\end{align*} Since \(q \geq 1\) and \(q > \frac{n}{2}-1\), we have 
\begin{align*}
    &\|F_{A_2}\| + \|A_2\|_4 \leq C_{q, M}\|A_2\|_{H^q} \\
    &\|\nabla_{\text{ref}}(A_1 - A_2)\| + \|A_1 - A_2\|_4 + \|A_1 - A_2\|_n \leq C_{q, M} \|A_1 - A_2\|_{H^q}.
\end{align*}Hence if \(\|A_1 - A_2\|_{H^q}\) is small enough,
\begin{align*}
    |\text{I}| + |\text{II}| + |\text{III}| + |\text{IV}| + |\text{V}| \leq \varepsilon\|b\|_{H_{A_2}^1}.
\end{align*} Thus by symmetry of the preceding computations
\begin{align*}
    \left|\int \langle F_{A_1}, D_{A_1} b\rangle - \int \langle F_{A_2}, D_{A_2} b\rangle\right| \leq \varepsilon\min\left\{\|b\|_{H_{A_1}^1}, \|b\|_{H_{A_2}^1}\right\},
\end{align*} which is the bound (\ref{lojasiewiczclaim:D*Fequivalence}) that we wanted to show. \par
Continuity of \(\mathcal{V}\) now readily follows from (\ref{lojasiewiczclaim:normequivalence}) and (\ref{lojasiewiczclaim:D*Fequivalence}), so the claim holds, as desired.
\end{claimproof}
\par We further have that \(\mathcal{V}\) is invariant under the action of \(H^{q+1}\) gauge transformations, which follows from
\begin{align*}
    \frac{\left|\int \langle F_{u(A)}, D_{u(A)} b\rangle\right|}{\left(\left\|\nabla_{u(A)} b\right\|^2 + \|b\|^2\right)^{\frac{1}{2}}} = \frac{\left|\int \langle u(F_A), u(D_A u^{-1}(b))\rangle\right|}{\left(\left\|u\left(\nabla_A u^{-1}(b)\right)\right\|^2 + \|b\|^2\right)^{\frac{1}{2}}} = \frac{\left|\int \langle F_A, D_A u^{-1}(b)\rangle\right|}{\left(\left\|\nabla_A u^{-1}(b)\right\|^2 + \|u^{-1}(b)\|^2\right)^{\frac{1}{2}}}.
\end{align*} Lastly, if we also have \(F_A \in H^q\), whence \(D_A^\ast F_A \in H^{q-1} \subset L^2\), then note that
\begin{align*}
    \mathcal{V}(A) \leq \left\|D_A^\ast F_A\right\|.
\end{align*}\par
We now give our variant of the R{\aa}de-Yang-Simon {\L}ojasiewicz inequality.
\begin{lemma}\label{lemma:lojasiewicz}
Let \(q > \frac{n}{2}-1\) and \(q \geq 1\). Given a fixed Yang-Mills connection $A_1 \in H^q(\calA_E),$ there exists a gauge-invariant $H^q$-neighborhood \(U\) of \(A_1\) and a constant \(\theta \in \left(0, \frac{1}{2}\right)\) 
%a constant \(\theta \in (0, \frac{1}{2}]\), and a constant \(C > 0\), 
such that for \(A \in U\),
\begin{align}\label{lojasiewicz:lojasiewiczinequality}
    \left|\mathcal{YM}(A) - \mathcal{YM}(A_1)\right|^{1 - \theta} \leq \mathcal{V}(A).
\end{align}
\end{lemma}
\begin{proof}
This is a straightforward adaptation of the arguments of \cite[Prop. 7.2]{rade} and \cite[Lemma 12]{byanguniqueness}, based on \cite[Theorem 3]{leonsimon}, which we include for the sake of completeness.
In the sequel, \(C\) denotes a positive constant, whose dependencies include \(A_1, E, q\), and the geometry of \(M\), which may increase from line to line. \par
    
    We first reduce proving (\ref{lojasiewicz:lojasiewiczinequality}) to a simpler setting. Since the desired inequality (\ref{lojasiewicz:lojasiewiczinequality}) is gauge-invariant, and since \(H^q\) Yang-Mills connections are smooth modulo \(H^{q+1}\) gauge, we may WLOG assume that \(A_1\) is smooth. Furthermore, the standard Coulomb gauge-fixing argument, e.g. \cite[Prop. 2.3.4]{donkron}, implies that there is an \(H^q\) neighborhood \(U\) of \(A_1\) such that if \(A \in U\), then \(A\) is gauge-equivalent to a connection $B$  which satisfies \(D_{A_1}^\ast(B - A_1) = 0\). Again by gauge-invariance, it suffices to prove (\ref{lojasiewicz:lojasiewiczinequality}) for \(A\) satisfying the gauge-fixing condition. Finally, by continuity of \(\mathcal{V}\), there exists an \(H^q\)-neighborhood of \(A_1\) such that for \(A\) in the neighborhood
    \begin{align*}
        \|D_A^\ast F_A\|_{H^{-1}} \leq 2 \mathcal{V}(A),
    \end{align*} so it suffices to prove (\ref{lojasiewicz:lojasiewiczinequality}) with \(\|D_A^\ast F_A\|_{H^{-1}}\) in place of \(\mathcal{V}(A)\). \par
    Now that we have reduced the proof of (\ref{lojasiewicz:lojasiewiczinequality}) to an easier setting, we next carry out the Lyapunov-Schmidt reduction step of Simon's argument. In the sequel, we change our reference connection to be the smooth connection \(A_1\), and hence define Sobolev norms with respect to \(A_1\). Since \(M\) is compact, the norms induced by \(\nabla_{\text{ref}}\) and \(\nabla_{A_1}\) are equivalent. \par 
    By Sobolev multiplication, the mapping \(a \mapsto \mathcal{YM}(A)\), where \(\nabla_A = \nabla_{A_1} + a\), is a well-defined, polynomial function on \(H^q\). 
    We define maps
    \begin{align*}
        M &: H^q \to H^{q-2} \\
        M(a) &:= D_A^\ast F_A, \\
        \tilde{M} &: H^q \to H^{q-2} \\
        \tilde{M}(a) &:= D_A^\ast F_A + D_{A_1}D_{A_1}^\ast a \\
        L_A &: H^q \to H^{q-2} \\
        L_A(b) &:= D_A^\ast D_Ab  + D_{A_1} D_{A_1}^\ast b+ (-1)^{3n+1}\ast (b \wedge \ast F_A).
    \end{align*}  Observe that \(L_A\) is the linearization of \(\tilde{M}\) at \(a\), i.e.,
    \begin{align*}
        L_A = (\nabla \tilde{M})(a).
    \end{align*}Now, \(L_{A_1}\) is self-adjoint and elliptic, so by elliptic theory
    \begin{align*}
        \mathcal{H} := \text{ker}(L_{A_1})
    \end{align*} is of finite dimension and consists of smooth sections. Define 
    \begin{align*}
        \Pi &: H^q \to \mathcal{H} \\
        \Pi(a) &:= \text{proj}_{\mathcal{H}}^{L^2}(a) \\
        N &: H^q \to H^{q-2} \\
        N &:= \tilde{M} + \Pi.
    \end{align*} Note that the linearization of \(N\) at \(a = 0\), i.e \((\nabla N)(0)\), equals \(L_{A_1} + \Pi\), which is invertible. Thus as in 
    \cite{leonsimon}, there exist open sets \(U_1 \subset H^q, U_2 \subset H^{q-2}\), with \(0 \in U_i\), and a diffeomorphism
    \begin{align*}
        \Psi : U_2 \to U_1
    \end{align*} which is inverse to \(N\). Next, set
    \begin{align*}
        V := \text{ker}(D_{A_1}^\ast) \subset H^q,
    \end{align*} and let \(\psi_1, \dots, \psi_m\) be an \(L^2\)-orthonormal basis of \(\mathcal{H} \cap V\). For \(\delta > 0\) small enough, we have a real-analytic \cite[(2.12)]{leonsimon} function
    \begin{align*}
        F &: B_\delta(0) \subset \R^m \to \R \\
        F(\xi) &= \mathcal{YM}\left(A_1 + \Psi\left(\sum_{j = 1}^m \xi^j \psi_j\right)\right).
    \end{align*} This completes the Lyapunov-Schmidt reduction step. \par
    Now that we have the existence of the local inverse \(\Psi\) of \(N\), we next show the following smoothing estimates for \(\Psi\), which are analogues of \cite[(2.9) and (2.10)]{leonsimon}: for \(a, b \in U_1\)
    \begin{align}
        \|\Psi(a) - \Psi(b)\|_{H^1} &\leq C\|a - b\|_{H^{-1}}, \label{lojasiewicz:PsiH1Lipschitz} \\
        \|\Psi(a) - \Psi(b)\|_{H^q} &\leq C\|a - b\|_{H^{q-2}}. \label{lojasiewicz:PsiHqLipschitz}
    \end{align} We first consider (\ref{lojasiewicz:PsiH1Lipschitz}). By elliptic theory we have for \(v \in H^q\)
    \begin{align*}
        \|v\|_{H^1} \leq C\|(L_{A_1}+\Pi)v\|_{H^{-1}}.
    \end{align*} Then by Sobolev multiplication, for a given \(\varepsilon > 0\), we may shrink \(U_1, U_2\) so that for \(a \in U_1\) and \(v \in H^q\)
    \begin{align*}
        \|((\nabla \tilde{M})(a) - L_{A_1})(v)\|_{H^{-1}} &\leq C\|(\nabla_{A_1} a) \# v + a \# \nabla_{A_1} v + a \# a \# v\|_{H^{-1}} \\
        &\leq C(\|\nabla_{A_1} a\|_{H^{q-1}} + \|a \# a\|_{H^{q-1}})\|v\|_{H^1} + C\|a\|_{H^q}\|\nabla_{A_1} v\|_{H^0} \\
        &\leq C(\|a\|_{H^{q}} + \|a\|_{H^{q}}^2)\|v\|_{H^1} + C\|a\|_{H^q}\|v\|_{H^1} \\
        &\leq \varepsilon \|v\|_{H^1}.
    \end{align*} Hence
    \begin{align*}
         \|(L_{A_1}+\Pi)v\|_{H^{-1}} &\leq \|((\nabla N)(a))v\|_{H^{-1}} + \|((\nabla \tilde{M})(a) - L_{A_1})(v)\|_{H^{-1}} \\
         &\leq \|((\nabla N)(a))v\|_{H^{-1}} + \frac{1}{2C}\|v\|_{H^1}.
    \end{align*} We deduce that for \(w \in H^{q-2}\)
    \begin{align*}
        \|((\nabla N)(a))^{-1} w\|_{H^1} \leq C\|w\|_{H^{-1}}.
    \end{align*} We may further shrink \(U_1, U_2\) so that \(U_2\) is an open ball about \(0\), which is convex and hence \(\Psi\) is defined on \(a + t(b-a)\) for \(a, b \in U_2, t\in [0, 1]\). Then
    \begin{align*}
        \|\Psi(a) - \Psi(b)\|_{H^1} &\leq \int_0^1 t\|((\nabla \Psi)(a + t(b-a)))(b-a)\|_{H^1} \, dt  \\
        &= \int_0^1 t\|((\nabla N)(\Psi(a + t(b-a))))^{-1}(b-a)\|_{H^1} \, dt \\
        &\leq C\|b-a\|_{H^{-1}}.
    \end{align*} Thus (\ref{lojasiewicz:PsiH1Lipschitz}) holds. \par 
    The other smoothing estimate (\ref{lojasiewicz:PsiHqLipschitz}) follows from a similar Sobolev multiplication argument. Thus we have the desired smoothing estimates for \(\Psi\). \par 
    Now that we have smoothing estimates for \(\Psi\), we proceed with proving the {\L}ojasiewicz inequality, which %we recall
    we only need to do for \(a \in U_1 \cap V\). We will do this, as in \cite{leonsimon}, by establishing the following two inequalities: given \(a \in U_1 \cap V\), upon which we have \(\xi \in B_{\delta}(0)\) defined by \(\Pi(a) = \sum_{j = 1}^m \xi^j \psi_j\),
    \begin{align}
        &|\mathcal{YM}(A_1 + a) - \mathcal{YM}(A_1 + \Psi(\Pi(a))| \leq C\|M(a)\|_{H^{-1}}^2 \label{lojasiewicz:upperbound} \\
        &|(\nabla F)(\xi)| \leq C\|M(a)\|_{H^{-1}} \label{lojasiewicz:lowerbound}.
    \end{align} \par
    We first consider (\ref{lojasiewicz:upperbound}). We compute as in as in \cite[(2.31)]{leonsimon}
    \begin{align*}
        &|\mathcal{YM}(A_1+\Psi(\Pi(a))) - \mathcal{YM}(A_1 + a)| \\
        = &\left| \int_0^1 \frac{d}{dt}\mathcal{YM}(A_1 + a + t(\Psi(\Pi(a)) - a)) \, dt \right| \\
        = & \left| \int_0^1 \int_M \langle M(a + t(\Psi(\Pi(a)) - a)), \Psi(\Pi(a)) - a\rangle \, dt\right| \\\leq &\int_0^1 \|M(a + t(\Psi(\Pi(a))-a)\|_{H^{-1}} \|\Psi(\Pi(a)) - a\|_{H^1} \, dt \\
        = &\int_0^1 \|M(a + t(\Psi(\Pi(a))-\Psi(N(a)))\|_{H^{-1}} \|\Psi(\Pi(a)) - \Psi(N(a))\|_{H^1} \, dt.
    \end{align*} Denoting
    \begin{align*}
        b = t(\Psi(\Pi(a)) - \Psi(N(a))),
    \end{align*} we next estimate \(\|M(a+b)\|_{H^{-1}}\). We compute
    \begin{align*}
        M(a+b) &= D_A^\ast F_A + (D_A^\ast D_A b + \nabla_A b \# b + F_A \# b + b \# b \# b) \\
        &= M(a) \\
        &+ (D_{A_1}^\ast D_{A_1} b + (\nabla_{A_1} a) \# b + a\#\nabla_{A_1} b + a \# a \# b + (\nabla_{A_1} b) \# b + a \# b \# b + F_A \# b + b \# b \# b).
    \end{align*} Since \(q > \frac{n}{2}-1\) and \(q \geq 1\), we may estimate the \(H^{-1}\) norm of the terms in the parentheses using Sobolev multiplication to obtain
    \begin{align*}
        \|D_{A_1}^\ast D_{A_1} b\|_{H^{-1}} &\leq C\|b\|_{H^1} \\
        \|(\nabla_{A_1} a) \# b \|_{H^{-1}} &\leq C\|\nabla_{A_1} a\|_{H^{q-1}}\|b\|_{H^1} \leq C\|a\|_{H^q}\|b\|_{H^1} \\
        \|a \# \nabla_{A_1} b\|_{H^{-1}} &\leq C\|a\|_{H^q} \|\nabla_{A_1} b\|_{H^0} \leq C\|a\|_{H^q}\|b\|_{H^1} \\
        \|a \# a \# b\|_{H^{-1}} &\leq C\|a \# a\|_{H^{q-1}} \|b\|_{H^1} \leq C\|a\|_{H^q}^2\|b\|_{H^1} \\
        \|(\nabla_{A_1} b) \# b\|_{H^{-1}} &\leq C\|b\|_{H^q} \|b\|_{H^1} \\
        \|a \# b \# b\|_{H^{-1}} &\leq C\|a\|_{H^q}\|b\|_{H^q}\|b\|_{H^{-1}} \\
        \|F_A \# b\|_{H^{-1}} &\leq C\|F_A\|_{H^{q-1}} \|b\|_{H^1} \leq C(1 +  \|a\|_{H^q}^2)\|b\|_{H^1} \\
        \|b \# b \# b\|_{H^{-1}} &\leq C\|b\|_{H^q}^2\|b\|_{H^1}.
    \end{align*}Thus if \(\|a\|_{H^q} + \|b\|_{H^q} \leq 1\),
    \begin{align}\label{lojasiewicz:D*FH-1continuous}
        \|M(a + b) - M(a)\|_{H^{-1}} \leq C\|b\|_{H^1}.
    \end{align} Now, we can make \(\|a\|_{H^q}\) small by shrinking \(U_1, U_2\). Moreover, since we assumed \(a \in V\), we have \(M(a) = \tilde{M}(a) = N(a) - \Pi(a)\), so the smoothing estimate (\ref{lojasiewicz:PsiHqLipschitz}) yields
    \begin{align*}
        \|b\|_{H^q} \leq C\|\Pi(a) - N(a)\|_{H^{q-2}} = C\|M(a)\|_{H^{q-2}}.
    \end{align*} Thus, since \(A_1\) is Yang-Mills, we can make \(\|M(a)\|_{H^{q-2}}\) small by shrinking \(U_1, U_2\). We may assume \(\|a\|_{H^q} + \|b\|_{H^q} \leq 1\) and consequently the bound
    \begin{align*}
        \|M(a + b)\|_{H^{-1}} \leq C(\|M(a)\|_{H^{-1}} + \|b\|_{H^1}).
    \end{align*} Therefore, the other smoothing estimate (\ref{lojasiewicz:PsiH1Lipschitz}) yields
    \begin{align*}
        |\mathcal{YM}(A_1 + \Psi(\Pi(a))) - \mathcal{YM}(A_1+a)| &\leq
        C(\|M(a)\|_{H^{-1}} + \|b\|_{H^1})\|\Psi(\Pi(a)) - \Psi(N(a))\|_{H^1} \\
        &\leq C(\|M(a)\|_{H^{-1}} + \|M(a)\|_{H^{-1}})\|M(a)\|_{H^{-1}}\\
        &= C\|M(a)\|_{H^{-1}}^2.
    \end{align*} This establishes (\ref{lojasiewicz:upperbound}). \par 
    We next establish (\ref{lojasiewicz:lowerbound}). Given \(\xi \in B_\delta(0) \subset \R^m\), set \(\iota(\xi) = \sum_{j = 1}^m \xi^j \psi_j\). Observe that
    \begin{align*}
        |(\partial_i F)(\xi)| &= \left|\int \langle M(\Psi(\iota(\xi))), \xi^i((d\Psi)(\iota(\xi)))(\psi_i)\rangle\right| \\
        &\leq |\xi|\|M(\Psi(\iota(\xi)))\|_{H^{-1}} \|((d\Psi)(\iota(\xi)))(\psi_i)\|_{H^1} \\
        &\leq C|\xi|\|M(\Psi(\iota(\xi)))\|_{H^{-1}}.
    \end{align*} Now, \(a \in V\), then \(\Pi(a) \in \mathcal{H} \cap V\), so there exists \(\xi\) such that \(\iota(\xi) = \Pi(a)\). We compute as in \cite[(2.20)]{leonsimon} that
    \begin{align*}
        |(\nabla F)(\xi)| &\leq C\|M(\Psi(\Pi(a))\|_{H^{-1}} \\
        &\leq C(\|M(a)\|_{H^{-1}}  + \|M(\Psi(\Pi(a))) - M(a)\|_{H^{-1}}).
    \end{align*} By (\ref{lojasiewicz:D*FH-1continuous}), we have that if 
    \begin{align*}
        \|a\|_{H^q} + \|\Psi(\Pi(a)) - a\|_{H^q} \leq 1,
    \end{align*} then
    \begin{align*}
        \|M(\Psi(\Pi(a))) - M(a)\|_{H^{-1}} \leq \|\Psi(\Pi(a)) - a\|_{H^1}.
    \end{align*} We already have that \(\|a\|_{H^q}\) is small. Moreover, since \(a = \Psi(N(a))\), and since \(N(a) - \Pi(a) = M(a)\) by the assumption \(a \in V\), the smoothing estimate (\ref{lojasiewicz:PsiHqLipschitz}) yields
    \begin{align*}
        \|\Psi(\Pi(a)) - a\|_{H^q} = \|\Psi(\Pi(a)) - \Psi(N(a))\|_{H^q} \leq C\|\Pi(a) - N(a)\|_{H^{q-2}} \leq C\|M(a)\|_{H^{q-2}},
    \end{align*} which is also small. Hence by the other smoothing estimate (\ref{lojasiewicz:PsiH1Lipschitz}),
    \begin{align*}
        |\nabla F(\xi)| &\leq C(\|M(a)\|_{H^{-1}}  + \|\Psi(\Pi(a))) - a\|_{H^1}) \\
        &\leq C(\|M(a)\|_{H^{-1}} + \|\Pi(a) - N(a)\|_{H^{-1}}) \\
        &= C\|M(a)\|_{H^{-1}}.
    \end{align*}This establishes (\ref{lojasiewicz:lowerbound}). \par
    Finally, now that we have (\ref{lojasiewicz:upperbound}) and (\ref{lojasiewicz:lowerbound}), the same computation as in \cite[(2.22)]{leonsimon} yields
    \begin{align*}
        |\mathcal{YM}(A_1+a) - \mathcal{YM}(A_1)|^{1-\theta} \leq C\|M(a)\|_{H^{-1}}.
    \end{align*} Since \(\mathcal{YM}(A_1 + a) \to \mathcal{YM}(A_1)\) uniformly in \(a\), we may take \(\theta\) smaller and shrink \(U_1\) to remove the constant \(C\) in the estimate. The lemma thus follows.
\end{proof}


\begin{comment}
\begin{rmk}
    {\color{red} Maybe make a remark here that you can use the flow to do the $q = n/2 - 1$ case, assuming that there is a good existence theory for Yang-Mills flow? I.e. you can use Struwe to do the Lojasiewicz in an $H^1$ neighborhood on a 4-manifold?}
\end{rmk}



\begin{proof}
   Since \(H^q\) Yang-Mills connections are smooth modulo \(H^{q+1}\) gauge for \(q > \frac{n}{2} - 1\), and since the desired inequality is gauge-invariant, we may assume that \(A_1\) smooth. By \cite{byanguniqueness}, we have that if \(A\) is sufficiently close to \(A_1\) in \(C^{k, \alpha}\), then
   \begin{align}\label{lojasiewicz:byanglojasiewicz}
       \left|\|F_A\|^2 - \|F_{A_1}\|^2\right|^{1-\theta} \leq 2\|D_A^\ast F_A\|.
   \end{align} Suppose for a contradiction that there does not exist \(\varepsilon > 0\) such that if \(A\) is a smooth connection in \(B_\varepsilon(A_1) \subset H^q\), then \(A\) satisfies
   \begin{align*}
       \left|\|F_A\|^2 - \|F_{A_1}\|^2\right|^{1-\theta} \leq 4\|D_A^\ast F_A\|.
   \end{align*} Then there is a sequence of smooth connections \(A_i \in B_{\frac{1}{i}}(A_1) \subset H^q\) for which
   \begin{align}\label{lojasiewicz:thingtocontradict}
       \left|\|F_A\|^2 - \|F_{A_1}\|^2\right|^{1-\theta} > 4\|D_A^\ast F_A\|.
   \end{align} Since \(A_1\) is a static solution of (\ref{ymf}), item (4) of Theorem \ref{thm:wellposedness} and the convergence of \(A_i \to A_1\) in \(H^q\) implies that for \(i\) large enough, the solutions \(A_i(t)\) of (\ref{ymf}) starting at \(A_i\) exist on \([0, 2]\) and converge in \(C^0([0, 2], H^q)\). In particular, since the hypotheses on \(q\) implies \(A \mapsto D_A^\ast F_A\) is continuous, we have
   \begin{align}\label{lojasiewicz:tensiongoingtozero}
       \|D_{A_i}^\ast F_{A_i}\| + \int_0^2 \hspace{-0.6em} \int |D_i^\ast F_i|^2 \to 0.
   \end{align} Furthermore, since \(H^{q-1} \inj L^p\) for some \(p > \frac{n}{2}\), we have
   \begin{align}\label{lojasiewicz:uniformLp}
       \limsup_{i \to \infty} \sup_{0 \leq t \leq 2} \|F_i(t)\|_p \leq C < \infty.
   \end{align} In the sequel, \(C\) is some positive constant which may increase but will be independent of \(i\). It follows from strong Uhlenbeck compactness that there exists a subsequence of the \(A_i\), also label by \(i\), and smooth gauge transformations \(u_i\) such that \(u_i(A_i(t))\) converges in \(C_{\text{loc}}^\infty((0, 2] \times M)\) to a smooth Yang-Mills connection \(A_\infty\). Since \(C^\infty\) convergence implies \(C^{k, \alpha}\) convergence, we in particular have that for any fixed \(t > 0\), \(u_i(A_i(t))\) satisfies (\ref{lojasiewicz:byanglojasiewicz}) for \(i\) large enough. By gauge-invariance, so does \(A_i(t)\). In view of (\ref{lojasiewicz:uniformLp}), parabolic Moser iteration, e.g. {\color{red} ESTIMATES PAPER} implies
   \begin{align*}
       \|F_i(t)\|_\infty \leq Ct^{-\frac{n}{2p}}.
   \end{align*} From the proof of \cite[Prop. 3.2]{instantons}, we have for \(i\) large enough and \(t \in (0, 2]\).
   \begin{align*}
       (\partial_t - \Delta)|D_i F_i|^2 \leq Ct^{-\frac{n}{2p}}|D_i F_i|^2
   \end{align*} Integrating over \(M\) and using that \(p > \frac{n}{2}\), we deduce that 
   \begin{align*}
       \|D_i^\ast F_i(t)\|^2 \leq e^{Ct^{1-\frac{n}{2p}}} \|D_{A_i}^\ast F_{A_i}\|^2.
   \end{align*} Since \(\theta \in (0, \frac{1}{2}]\) and since \(\|D_{A_i}^\ast F_{A_i}\| \leq 1\) for \(i\) large enough due to (\ref{tensiongoingtozero}), we have
   \begin{align*}
       |\|F_{A_i}\|^2 - \|F_{A_1}\|^2|^{1-\theta} &\leq |\|F_i(t)\|^2 - \|F_{A_1}\|^2|^{1-\theta} + |\|F_{A_i}\|^2 - \|F_i(t)\|^2|^{1-\theta} \\
       &\leq 2\|D_i^\ast F_i(t)\| + \left(\int_0^t \|D_i^\ast F_i(s)\|^2 \, ds \right)^{1-\theta} \\
       &\leq 2e^{Ct^{1-\frac{n}{2p}}}\|D_{A_i}^\ast F_{A_i}\| + Ct\|D_{A_i}^\ast F_{A_i}\|^{2(1-\theta)} \\
       &\leq (2e^{Ct^{1-\frac{n}{2p}}} + Ct^{1-\theta})\|D_{A_i}^\ast F_{A_i}\|.
   \end{align*} Fixing \(t > 0\) sufficiently small so that the term in the parentheses in the last line is less than \(4\), we find that (\ref{lojasiewicz:thingtocontradict}) is contradicted for \(i\) sufficiently large, as desired. \par 
   Hence the desired inequality holds for smooth connections, and consequently for \(H^q\) connections by approximation.
   \end{proof}
\end{comment}
We also have the following uniform statement for the flat connections.

\begin{lemma}\label{lemma:uniformlojasiewicz}
    Let \(q > \lceil\frac{n}{2}-1\rceil\) and \(q \geq 1\).
    Given a \( (G, \rho) \)-bundle \(E \to M\), let $\mathcal{F} \subset H^q \left(\mathcal{A}_E\right)$
    %\mathscr{A}_E \right)$ 
    denote the set of flat connections. There exists \(\theta_0 \in (0, \frac12) \) and a gauge-invariant $H^q$-neighborhood \(U \supset \mathcal{F}\) such that for all \(A \in U\), we have
    \begin{align}\label{uniformlojasiewicz:est}
    \mathcal{YM}(A)^{1 - \theta_0} \leq \mathcal{V}(A).
\end{align} 
Moreover, the solution of (\ref{ymf}) starting at \(A \in U\) exists for all time and converges polynomially to a flat connection in $H^q$.
\end{lemma}
\begin{proof}
    Suppose for %a 
    contradiction that there does not exist \(\theta \in (0, \frac{1}{2})\) such that every \(B \in \mathcal{F}\) has an \(H^q\) neighborhood \(U_B\) for which the inequality (\ref{uniformlojasiewicz:est}) holds. Then there exists \(\theta_i \searrow 0\) and a sequence of flat connections \(B_i\) such that for each \(i\), the inequality (\ref{uniformlojasiewicz:est}) with \(\theta_0 = \theta_i\) fails on every neighborhood of \(B_i\). Since the \(B_i\) are all flat, strong Uhlenbeck compactness implies that we may pass to a subsequence, also labeled by \(i\), such that there exist gauge transformations \(u_i \in H^{q+1}\) for which \(u_i(B_i)\) converge in \(H^q\) to some flat connection \(B_\infty\). By Lemma \ref{lemma:lojasiewicz}, there exists \(\theta_\infty \in (0, \frac{1}{2})\) and a neighborhood \(U_{B_\infty}\) of \(B_\infty\) such that (\ref{uniformlojasiewicz:est}) holds on \(U\) with \(\theta_0 = \theta_\infty\). Thus for \(i\) sufficiently large, (\ref{uniformlojasiewicz:est}) holds on some neighborhood \(U\) of \(u_i(B_i)\). But (\ref{uniformlojasiewicz:est}) is gauge-invariant, so (\ref{uniformlojasiewicz:est}) holds on the neighborhood \(u_i^{-1}(U_{B_{\infty}})\) of \(B_i\), which is a contradiction, as desired. Thus there exists \(\theta_0 \in (0, \frac{1}{2})\) such that every \(B \in \mathcal{F}\) has a neighborhood \(U_B\) for which (\ref{uniformlojasiewicz:est}) holds with \(\theta_0\). \par 
    By Theorem \ref{thm:continuityatinfinitetime}, we may assume \(U_B\) is such that for all \(A \in U_B\), the solution of (\ref{ymf}) starting at \(A\) exists for all time and converges polynomially in \(H^q\) to a flat connection. Thus the open set
    \begin{align*}
        U := \bigcup_{B \text{ flat, } u \in H^{q+1}(\mathcal{G}_E)} u(U_B)
    \end{align*}is a neighborhood of the flat connections satisfying the assertions of the lemma, as desired. 
\end{proof}



\vspace{5mm}



\subsection{Proof of Theorem \ref{thm:continuityatinfinitetime}} \label{continuityatinftime}
\begin{proof} 
    By Theorem \ref{thm:wellposedness}(1), for $\delta$ is sufficiently small, $A'(T_0)$ will be contained in the neighborhood $B_{\eps}(A_\infty)$ for $T_0$ large enough; we can therefore reduce to the case that the given solution $A(t) \equiv A_\infty$ is static. In particular, given $\eps > 0,$ it suffices to determine $\delta > 0$ such that for $A'_0 \in B_\delta(A_\infty),$ either alternative (1) holds for $A'(t)$ or else
    \begin{equation}\label{alternate(2)}
        \| A'(t) - A_\infty \| < \eps
    \end{equation}
    for all $t$ and $A'(t)$ converges, which is equivalent to (2). 

    Suppose that for \(\delta_\infty > 0\) to be determined, we have
    \begin{align}\label{continuityatinfinitetime:energydoesntdropmuch}
        \mathcal{YM}(A'(t)) > \mathcal{YM}(A_\infty) - \delta_\infty
    \end{align} for $t \in \LB 0 , \tau \right),$ where $\tau \leq \infty$ is the maximal such time. The two cases are distinguished by (1) $\tau < \infty$ and (2) $\tau = \infty.$ %$\mathcal{YM}(A(t)) < \mathcal{YM}(A_\infty)$ for $t > z \geq 0,$ or else (2) $\mathcal{YM}(A(t)) \geq \mathcal{YM}(A_\infty)$ for all $t.$
    

    Let \(\kappa > 0\) be such that the {\L}ojasiewicz inequality (\ref{lojasiewicz:lojasiewiczinequality}) holds on \(B_{3\kappa}(A_\infty)\). 
    %We first wish to show that for \(A_0'\) close to \(A_0\), if case (1) of Theorem \ref{thm:continuityatinfinitetime} does not hold then \(A'(t)\) is eventually contained in \(B_\kappa(A_\infty)\). To do this, we collect some estimates on such a solution.
%Further suppose that for \(t\) in some interval \([a, b]\), \(A'(t)\) is contained in \(B_\kappa(A_\infty)\), so that \(A'(t)\) satisfies the {\L}ojasiewicz inequality (\ref{lojasiewicz:lojasiewiczinequality}). We let \(z \in (a, b)\) denote a time such that \(\mathcal{YM}(A'(z)) = \mathcal{YM}(A_\infty)\), if such a \(z\) exists.
We assume that $\delta < \kappa$ is small enough that $A'(t) \subset B_{\kappa}(A_\infty)$ for $0 \leq t \leq 1.$ Now, supposing that \(A'(t) \in B_{3\kappa}(A_\infty)\) for \(t\) in some interval \([a, b] \subset \LB 0, \tau \right) \), %Since the Yang-Mills energy is nonincreasing along the flow, 
we compute
    \begin{align*}
        \int_a^{b} & \left\|D_{A'(t)}^\ast F_{A'(t)}\right\| \, dt \\
        &= \int_a^z \frac{-\frac{d}{dt}\left(\mathcal{YM}(A'(t)) - \mathcal{YM}(A_\infty)\right)}{\left\|D_{A'(t)}^\ast F_{A'(t)}\right\|} \, dt + \int_z^{b} \frac{\frac{d}{dt}\left(\mathcal{YM}(A_\infty) - \mathcal{YM}(A'(t))\right)}{\left\|D_{A'(t)}^\ast F_{A'(t)}\right\|} \, dt  \\
        &\leq \int_a^z \frac{-\frac{d}{dt} (\mathcal{YM}(A'(t))-\mathcal{YM}(A_\infty))}{(\mathcal{YM}(A'(t)) - \mathcal{YM}(A_\infty))^{1-\theta}} \, dt + \int_z^{b} \frac{\frac{d}{dt} (\mathcal{YM}(A_\infty) - \mathcal{YM}(A'(t)))}{(\mathcal{YM}(A_\infty) - \mathcal{YM}(A'(t)))^{1-\theta}} \, dt \\
        &\leq \frac{1}{\theta}\left((\mathcal{YM}(A'(a)) - \mathcal{YM}(A_\infty)))^\theta + (\mathcal{YM}(A_\infty) - \mathcal{YM}(A'(b)))^{\theta}\right) \\
        &< \frac{1}{\theta}\left((\mathcal{YM}(A'_0) - \mathcal{YM}(A_\infty))^\theta + \delta_{\infty}^{\theta}\right).
    \end{align*}
    Assuming that $\delta$ and $\delta_\infty$ are sufficiently small, Theorem \ref{thm:Hkestimates} shows that in fact $A'(t)$ remains inside $B_{2\kappa}(A_\infty)$ for all $t \in \LB 0, \tau \right).$ The rest of the argument is a standard application of the \L ojasiewicz-Simon inequality and Gronwall's inequality to $\mathcal{YM}(A'(t)),$ which we omit.
    \begin{comment}
    Furthermore, the \(H^q\) norm of \(A'(t)\) stays bounded as long as \(A'(t)\) remains in \(B_\kappa(A_\infty)\). Thus by epsilon-regularity, if \(a + 1 < b\), then there exists \(K > 0\), depending only on \(A_\infty, q,\) and the geometry of \(M\), such that
    \begin{align*}
        \sup_{a+1 \leq t \leq b} \left(\|A(t)\|_{H^q} + \|F(t)\|_\infty\right) \leq K.
    \end{align*} Finally, let \(C_{(\ref{Hkestimates:Hkbounds})}\) be the constant from Theorem \ref{thm:Hkestimates} with \(K\) as above, \(k = \lfloor q \rfloor, \tau = 1,\) and \(\Omega = M\). With these bounds in hand, we show that if \(A_0'\) is close to \(A_0\) and \(A'(t)\) satisfies (\ref{continuityatinfinitetime:energydoesntdropmuch}), then \(A'(t)\) indeed is contained in \(B_\kappa(A_\infty)\) for all time. \par
    Choose \(\delta_\infty > 0\) such that 
    \begin{align*}
        \frac{\delta_\infty^\theta}{\theta} < \frac{\min\{1, \kappa\}}{20C_{(\ref{Hkestimates:Hkbounds})}(1+K^{(k+1)!})}. 
    \end{align*} Now, since \(A(t)\) converges to \(A_\infty\), there exists \(T > 1\) such that the following hold:
    \begin{enumerate}
        \item[(i)] \(A(t) \in B_{\frac{\kappa}{5}}(A_\infty)\) for \(t \geq T\)
        \item[(ii)] \(\max\{\mathcal{YM}(A(T)) - \mathcal{YM}(A_\infty), 0\} < \frac{\delta_\infty}{2}\).
    \end{enumerate} It thus follows from item (\ref{wellposedness:wellposedness}) of Theorem \ref{thm:wellposedness} that there exists \(\delta_1 > 0\) such that if \(A_0' \in B_{\delta_1}(A_0)\), then we have the following:
    \begin{enumerate}
        \item \(A'(t)\) exists on \([0, T+2]\); 
        \item \(A'([T, T+2]) \subset B_{\frac{\kappa}{4}}(A_\infty)\); 
        \item \(\max\{\mathcal{YM}(A'(T)) - \mathcal{YM}(A_\infty), 0\} < \delta_\infty\).
    \end{enumerate} Suppose \(A'(t)\) satisfies (\ref{continuityatinfinitetime:energydoesntdropmuch}) as long as \(A'(t)\) exists. Let \(T' > T + 2\) be the maximal time such that \(A'(t)\) is contained in \(B_\kappa(A_\infty)\) for \(t \in [T, T')\). By choice of \(\delta_\infty\), we have \(\rho < 1\). Thus by (\ref{Hkestimates:Hkbounds}) with \(k = \lfloor q \rfloor\), so that \(\lceil \frac{n}{2}-1\rceil \leq k \leq q \leq k+1\), we have
    \begin{align*}
        \|A'(T+2) - A'(T')\|_{H^q} \leq 2C_{(\ref{Hkestimates:Hkbounds})}(1+K)^{(k+1)!}\rho < \frac{\kappa}{4}
    \end{align*}  Hence
    \begin{align*}
        \|A'(T') - A_\infty\|_{H^q} \leq \|A'(T') - A'(T+2)\|_{H^q} + \|A'(T+2) - A_\infty\|_{H^q} < \frac{\kappa}{2}.
    \end{align*} Thus by continuity, \(T' = \infty\). Then by standard arguments, the sequence \(A'(n)\) converges modulo gauge in \(H^q\) to a Yang-Mills connection \(\tilde{A}_\infty'\). Since \(A'(n) \in B_{\kappa}(A_\infty)\) for all \(n\), gauge-invariance of the {\L}ojasiewicz ineqaulity implies that \(\mathcal{YM}(\tilde{A}_\infty') = \mathcal{YM}(A_\infty)\). In particular, \(\mathcal{YM}(A(t)) \geq \mathcal{YM}(A_\infty)\) for all \(t\). Therefore, \(A(t)\) converges to a Yang-Mills connection in \(B_{\kappa}(A_\infty)\) in view of Theorem \ref{thm:Hkestimates} and the polynomial decay of \(\mathcal{YM}(A(t)) - \mathcal{YM}(A_\infty)\), e.g. see \cite[Prop. 7.4]{rade}. \par 
    We next establish the estimate (\ref{continuityatinfinitetime:continuityestimate}). Given \(\varepsilon > 0\), choose \(\tilde{T}\) such that \(\|A(t) - A_\infty\|_{H^q} < \frac{\varepsilon}{8}\) for \(t \geq \tilde{T}\) and such that
    \begin{align*}
        \frac{1}{\theta}(\mathcal{YM}(A(\tilde{T})) - \mathcal{YM}(A_\infty))^\theta < \frac{\min\{1, \varepsilon\}}{20C_{(\ref{Hkestimates:Hkbounds})}(1 + K^{(k+1)!})}.
    \end{align*} Then by item (\ref{wellposedness:wellposedness}) of Theorem \ref{thm:wellposedness}, we have for \(\delta\) small enough that 
    \begin{align*}
        \frac{1}{\theta}(\mathcal{YM}(A'(\tilde{T})) - \mathcal{YM}(A_\infty))^\theta < \frac{\min\{1, \varepsilon\}}{10C_{(\ref{Hkestimates:Hkbounds})}(1 + K^{(k+1)!})}
    \end{align*} and \(\|A'(\tilde{T}) - A_\infty\|_{H^q} < \frac{\varepsilon}{4}\). Since \(\mathcal{YM}(A'(t)) \geq \mathcal{YM}(A_\infty)\) for all \(t\), the preceding computations imply that for \(t \geq \tilde{T}\)
    \begin{align*}
        \|A'(t) - A(t)\|_{H^q} &\leq \|A'(t) - A'(\tilde{T})\|_{H^q} + \|A'(\tilde{T}) - A_\infty\|_{H^q} +\|A(t) - A_\infty\|_{H^q} \\
        &< \varepsilon.
    \end{align*} Hence (\ref{continuityatinfinitetime:continuityestimate}) holds for \(t \geq \tilde{T}\). On the other hand, by item (\ref{wellposedness:wellposedness}) of Theorem \ref{thm:wellposedness}, we may shrink \(\delta\) so that (\ref{continuityatinfinitetime:continuityestimate}) holds on \([0, \tilde{T}]\), as desired. \par
    By the preceding arguments, we have for \(A_0'\) satisfying the hypotheses of the proposition, \(A'(t)\) exists as long as \(\mathcal{YM}(A'(t)) \geq \mathcal{YM}(A_\infty) - \delta_\infty\), i.e. \(A'(t)\) cannot blow up before its energy drops below \(\mathcal{YM}(A_\infty) - \delta_\infty\). Therefore, the only other option is case (1) of Theorem \ref{thm:continuityatinfinitetime}, as desired.
    \end{comment}
\end{proof}

\begin{comment}   
\begin{proof}
    In view of Item \ref{wellposedness:wellposedness} of Theorem \ref{thm:wellposedness}, we WLOG assume \(A_0 = A_\infty\) and \(\mu = \delta_1\). \par 
    Since \(A_\infty\) is in particular a static solution of (\ref{ymf}), we may assume the following for \(\delta_1\) sufficiently small:
    \begin{enumerate}
        \item \(A(t)\) exists on \([0, 2]\) and \(A([0, 2]) \subset B_{2\delta_1}(A_\infty\).
        \item For some \(\kappa \geq 3\delta_1\), the {\L}ojasiewicz inequality (\ref{lojasiewicz:lojasiewiczinequality}) holds on \(B_\kappa(A_\infty)\).
    \end{enumerate} \par 
   Suppose that \(A(t)\) satisfies \(\mathcal{YM}(A(t)) \geq \mathcal{YM}(A_\infty) - \delta_1\) as long as \(A(t)\) exists. We first argue that if \(\delta_1\) is small enough, \(A(t)\) exists and is contained in \(B_{\kappa}(A_\infty)\) for all time. Let \(T' > 2\) denote the maximal time that \(A(t)\) is contained in \(B_\kappa(A_\infty)\). We let \(z \in (0, T')\) denote a time for which \(\mathcal{YM}(A(z)) = \mathcal{YM}(A_\infty)\), if such a \(z\) exists. Since the Yang-Mills energy is nonincreasing along the flow, we compute
    \begin{align*}
        \int_0^{T'} \|D^\ast F(t)\| \, dt &= \int_0^z \frac{-\frac{d}{dt}\left(\mathcal{YM}(A(t)) - \mathcal{YM}(A_\infty)\right)}{\|D^\ast F(t)\|} \, dt + \int_z^T \frac{\frac{d}{dt}\left(\mathcal{YM}(A_\infty) - \mathcal{YM}(A(t))\right)}{\|D^\ast F(t)\|} \, dt  \\
        &\leq \int_0^z -\frac{\frac{d}{dt} (\mathcal{YM}(A(t))-\mathcal{YM}(A_\infty))}{(\mathcal{YM}(A(t)) - \mathcal{YM}(A_\infty))^{1-\theta}} \, dt + \int_z^{T'} \frac{\frac{d}{dt} (\mathcal{YM}(A_\infty) - \mathcal{YM}(A(t)))}{(\mathcal{YM}(A_\infty) - \mathcal{YM}(A(t)))^{1-\theta}} \, dt \\
        &\leq \frac{1}{\theta}\left((\mathcal{YM}(A_1) - \mathcal{YM}(A_\infty)))^\theta + (\mathcal{YM}(A_\infty) - \mathcal{YM}(A(T')))^{\theta}\right) \\
        &\leq\frac{1}{\theta}\left((\mathcal{YM}(A_1) - \mathcal{YM}(A_\infty)))^\theta + \delta_1^{\theta}\right) \\
        &=: \rho.
    \end{align*} Furthermore, since the \(H^q\) norm of \(A(t)\) remains bounded as long as \(A(t)\) remains in \(B_\kappa(A_\infty)\), we have by epsilon-regularity that for some \(K> 0\)    \begin{align*}
        \sup_{1 \leq t \leq T} \left(\|A(t)\|_{H^q} + \|F(t)\|_\infty\right) \leq K.
    \end{align*} Thus by Cor. \ref{cor:Hkestimateslowregularity} with \(k = \lfloor q \rfloor\), so that \(\lceil \frac{n}{2}-1\rceil \leq k \leq q \leq k+1\), we have that if \(\delta_1\) is sufficiently small such that \(\rho \leq 1\), then
    \begin{align*}
        \|A(2) - A(T)\|_{H^q} \leq C_{q, K, M} \rho.
    \end{align*} Taking \(\delta_1\) smaller such that 
    \begin{align*}
        \rho < \frac{1}{C_{q, K, M}}\frac{\min\{1, \kappa\}}{10},
    \end{align*} we have
    \begin{align*}
        \|A(T') - A_\infty\|_{H^q} \leq \|A(T') - A(2)\|_{H^q} + \|A(2) - A_\infty\|_{H^q} < \frac{\kappa}{3}.
    \end{align*} Thus by continuity, \(T' = \infty\). Hence by standard arguments, the sequence \(A(n)\) converges to a Yang-Mills connection \(A_\infty'\) in \(H^q\). Since \(A(n) \in B_{\kappa}(A_\infty)\), we have \(\mathcal{YM}(A_\infty') = \mathcal{YM}(A_\infty)\). In particular, \(\mathcal{YM}(A(t)) \geq \mathcal{YM}(A_\infty)\) for all \(t\), and hence convergence of \(A(t)\) in \(t\) follows from the fact that \(\mathcal{YM}(A(t)) - \mathcal{YM}(A_\infty)\) is polynomially decaying, e.g. see \cite[Prop. 7.4]{rade}. \par  
    By the preceding arguments, we have for \(A_1\) satisfying the hypotheses of the proposition, \(A(t)\) exists as long as \(\mathcal{YM}(A(t)) \geq \mathcal{YM}(A_\infty) - \delta_1\), i.e. \(A(t)\) cannot blow up before its energy drops below \(\mathcal{YM}(A_\infty)\). Therefore, the only other option is item (1) of this proposition.
\end{proof}
\end{comment}



\subsection*{Data availability and conflict of interest statement}

This paper has no associated data and the authors have no conflicts of interest.


\bibliographystyle{amsinitial}
\bibliography{biblio}


  
\end{document}







\begin{comment}

Gauge theory in dimension four is the classical setting for this problem. An instanton is, up to orientation, a solution of the anti-self-duality equation
$$F_A^+ = 0.$$
As in all dimensions, a Yang-Mills connection is a solution of the equation
$$D_A^*F_A = 0,$$
which states that $A$ is a critical point of the Yang-Mills functional
$$\YM(A) = \frac12 \int_M |F_A|^2\, dV.$$
Since instantons minimize the Yang-Mills functional, % among all connections on a given bundle,
they naturally solve the Yang-Mills equation. %the question is then whether any higher critical points exist.
However, even for the case $M = S^4$ and structure group $\SU(2),$ the existence of higher critical points remained unknown until 1989, when Sibner, Sibner, and Uhlenbeck \cite{} constructed non-instanton solutions of the Yang-Mills equations. This demonstrates the difficulty of the problem. %This shows that the converse statement can only hold under further assumptions. %Several other examples of non-instanton $\SU(2)$-Yang-Mills connections on $S^4$ have been constructed, but it is unknown whether this list is exhaustive; in particular, we still do not know whether any of the known examples attains the lowest nontrivial critical value. % of the Yang-Mills energy on $S^4$.
%This demonstrates the subtlety of the question.

\end{commment}




\subsection{Background} 

Mathematical gauge theory %, particularly as a branch of geometric analysis,
studies both first-order and second-order elliptic systems involving a connection (``gauge field'') on a principal bundle over a Riemannian manifold. Solutions of the first-order system are referred to as \emph{instantons}, % ; these are the central geometric objects of the field. 
while solutions of the second-order system are referred to as \emph{Yang-Mills connections}. %; these are the central analytic objects of the field.
%; in particular, these are the critical points of 
%which are only required to satisfy the well-known partial differential equation
%$$D_A^*F_A = 0.$$
%This equation states that $A$ is a critical point of
%the Yang-Mills functional
%$$\YM(A) = \frac12 \int_M |F_A|^2 \, dV.$$
In all cases of interest, the instanton equation directly implies the Yang-Mills equation. %; indeed, instantons typically minimize the above functional.
%Meanwhile, the %question of whether the Yang-Mills equation implies the instanton equation %is much more subtle. %, and
The converse is more subtle, and
has generated a substantial body of work.
%The converse is much more subtle, %, and always requires further assumptions. 
%and has generated a substantial body of work.
\begin{comment}
The question is important because geometric applications typically depend on constructing nonempty moduli spaces of instantons, % on a given base manifold,
whereas variational procedures can succeed only in producing Yang-Mills connections. %on a given base manifold, whereas gauge-theoretic invariants generally depend on the existence of nonempty moduli spaces of instantons.
\end{comment}

\begin{comment}
%A key difficulty in 
Mathematical gauge theory %, particularly as a branch of geometric analysis,
studies both the so-called ``first-order'' and ``second-order'' Yang-Mills equations. Solutions of the former are referred to as \emph{instantons}, % ; these are the central geometric objects of the field. 
while solutions of the latter are referred to as \emph{Yang-Mills connections}. %; these are the central analytic objects of the field.
%; in particular, these are the critical points of 
%which are only required to satisfy the well-known partial differential equation
%$$D_A^*F_A = 0.$$
%This equation states that $A$ is a critical point of
%the Yang-Mills functional
%$$\YM(A) = \frac12 \int_M |F_A|^2 \, dV.$$
In all cases of interest, the (first-order) instanton equation implies the (second-order) Yang-Mills equation. %; indeed, instantons typically minimize the above functional.
However, the converse is much more subtle,
%The converse is much more subtle, %, and always requires further assumptions. 
and has generated a substantial body of work.
%The importance of the question stems from the fact that variational procedures can at best produce Yang-Mills connections, whereas gauge-theoretic invariants generally come from moduli spaces of instantons.
%However, the converse is usually difficult to establish, and has been the subject of a fair amount of work. %a significant body of work. 
\end{comment}

Gauge theory in dimension four provides the classical setting for this problem. An instanton is, up to orientation, a solution of the anti-self-duality equation
$$F_A^+ = 0.$$
A Yang-Mills connection, as in all dimensions, is a solution of the equation
$$D_A^*F_A = 0.$$
The latter %states that $A$ is a critical point of
is simply the Euler-Lagrange equation of the Yang-Mills functional
$$\YM(A) = \frac12 \int_M |F_A|^2\, dV.$$
Since instantons are minimizers of $\YM,$ % among all connections on a given bundle,
they naturally solve the Yang-Mills equation. However, %the question is then whether any higher critical points exist.
even for the trivial $\SU(2)$-bundle over $S^4,$ higher critical points of the Yang-Mills functional turned out to exist %; this is the celebrated result of Sibner, Sibner, and Uhlenbeck
\cite{sibnersuhlenbeck}.
Strong assumptions are therefore necessary to conclude that a solution of the Yang-Mills equation must be an instanton, as required for geometric or physical applications.

%The question remains important because geometric applications typically depend on constructing nonempty moduli spaces of instantons, whereas variational procedures can succeed only in producing Yang-Mills connections. %, as required for geometric or physical applications.  %on a given base manifold, whereas gauge-theoretic invariants generally depend on the existence of nonempty moduli spaces of instantons.. % the existence of higher critical points remained unknown until 1989, when Sibner, Sibner, and Uhlenbeck \cite{} constructed non-instanton solutions of the Yang-Mills equations. %This demonstrates the difficulty of the problem. %This shows that the converse statement can only hold under further assumptions. %Several other examples of non-instanton $\SU(2)$-Yang-Mills connections on $S^4$ have been constructed, but it is unknown whether this list is exhaustive; in particular, we still do not know whether any of the known examples attains the lowest nontrivial critical value. % of the Yang-Mills energy on $S^4$.
%This demonstrates the subtlety of the question.
%The typical result along these lines takes the form of a ``gap theorem:'' a solution $A$ of the second-order Yang-Mills equations, with energy close to that of an instanton, must indeed be an instanton. 
%The most common assumption is that the energy be no more than a small amount above that of an instanton, {\it i.e.} the infimum over all connections on the given bundle. %A ``gap theorem'' then states that $A$ must in fact solve the first-order Yang-Mills equations.
\begin{comment}
Gauge theory in dimension four is the best-known setting for this question. %instance of the gap phenomenon.
Here, the instanton equation is the celebrated anti-self-duality equation
$$F_A^+ = 0,$$
while the general Yang-Mills equation is
$$D_A^*F_A = 0.$$
As in all 
must be anti-self-dual (or self-dual).
\end{comment}
%The typical result along these lines takes the form of a ``gap theorem;'' this states that a solution $A$ of the second-order Yang-Mills equations, with energy close to that of an instanton, must indeed be an instanton.
Positive results in this direction were pioneered by Bourguignon, Lawson, and Simons \cite{bourguignonlawsonsimons,bourguignonlawson}. These either take the form of ``stability theorems,'' stating that a Yang-Mills connection with nonnegative second variation must be an instanton, or ``gap theorems,'' stating that a Yang-Mills connection with appropriately small self-dual curvature is in fact an instanton. We are interested here in the latter type of result. %This note is concerned with the latter.
%The typical positive result takes the form of a ``gap theorem,'' stating that a Yang-Mills connection with nearly minimal energy must be an instanton.

The standard $L^2$ gap theorem on $S^4,$ due to Min-Oo \cite{minoogap}, runs as follows: suppose that $A$ is a Yang-Mills connection on a bundle $E \to S^4,$ satisfying the Yang-Mills equation. If also
$$\| F_A^+ \|_{L^2(S^4)} < \delta,$$
a positive universal constant, then in fact $A$ is an instanton ($F_A^+ \equiv 0$). A proof can also be found in Donaldson-Kronheimer \cite{donkron}, \S 2.3.10.


The $L^2$ gap theorem also holds for 4-manifolds with appropriately positive curvature \cite{gerhardtgap}, for bundles that carry a flat connection \cite{}, or for generic metrics and structure group $\SU(2)$ \cite{} (using the Freed-Uhlenbeck Theorem).
For %bundles with nontrivial first Pontryagin class over
general compact Riemannian 4-manifolds, the gap question remains wide open.\footnote{An ArXiv posting from May 2015 claims to prove the more general statement that the spectrum of the Yang-Mills functional on a compact 4-manifold is discrete, but the argument given there is incorrect.}



\begin{comment}
The Min-Oo gap theorem has been generalized by Gerhardt to 4-manifolds with appropriately positive curvature. The gap theorem is also easy to prove in the case of a flat bundle, and for a generic metric, by appealing to the Freed-Uhlenbeck transversality theorem.
Meanwhile, for bundles with nontrivial first Pontryagin class over general compact Riemannian 4-manifolds, the gap question remains open. %, despite an unresolved claim from May 2015.
%This again demonstrates the subtlety of the problem.
\end{comment}

Gap results have also been proven in higher-dimensional gauge theory, albeit under even stricter restrictions. %under limited circumstances.
Bourguignon-Lawson, ...

\begin{comment}
Note that these gap theorems all share an obvious drawback; namely, that the connection in question is required to solve the (second-order) Yang-Mills equation in the first place. Although variational methods succeed at producing critical points of $\YM$ in low dimensions ($n < 4$), the Yang-Mills equation itself is extremely difficult to solve in higher dimensions ($n \geq 4,$ and especially so for $n > 4$). This is due both to the gauge freedom of the above functional and to the so-called ``bubbling phenomenon;'' meanwhile, from the point-of-view of geometric analysis, these are what make the subject so interesting.
\end{comment}

In this note, we approach the gap question from the perspective of Yang-Mills flow; this is the downward gradient flow of the Yang-Mills functional, given by
$$\frac{\p A}{\p t} = - D_A^*F_A.$$
Under favorable circumstances---indeed, the same as those of 
the gap theorems discussed above---we will show that a general connection with appropriately small curvature deforms smoothly to an instanton under the flow. %In this way, we need not begin with a solution of the Yang-Mills equation in order to solve the instanton equation.
In this sense, we remove the assumption %from several existing gap theorems
that the connection in question solve the Yang-Mills equation {\it a priori}. %, giving us a stronger method to construct instantons. %; hence, in theory, we need not begin with a solution of the Yang-Mills equation in order to solve the instanton equation. %In the opposite direction, we can also prove certain non-existence results for general connections with small curvature on a given bundle.
%This approach cuts in two directions: first,
Meanwhile, in cases where an instanton cannot exist for topological reasons, we %learn that a general connection with small curvature also cannot exist; this gives us 
obtain a gap statement for general connections on the given bundle.

%We now state our results.
%Our results are as follows.
We state our results here. %Since they mostly follow from the author's previous work and that of Dayaprema \cite{}, the proofs will scarcely be longer than the statements.

%On the other hand, we learn that if a connection with energy close to the required energy does exist, then it deforms under Yang-Mills flow to an instanton. We expect the latter type of result to be useful for constructing instantons in certain higher-dimensional circumstances.







%We now state our results.
